% DER_XIIb-a: Grothendieck's text with the commentary set in, then the overview, inside a group that applies the chunk's own definitions
\dgrpart{1}{XIIb-a}{Section 5, part 1 of 2, manuscript pages 59--122}{XIIb-a\enspace\textperiodcentered\enspace Chapter XII, section 5, part 1 of 2\enspace\textperiodcentered\enspace manuscript pages 59--122}{Chapter XII \textperiodcentered\ \textup{XIIb-a}}
\dgrchunkbegin
% DER_XIIb-a: the definitions of the chunk's preamble, applied inside its group
\let\printedpage\relax\newcommand{\printedpage}[1]{\marginpar{\footnotesize p.~#1}}
\let\msp\relax\newcommand{\msp}[1]{\marginpar{\footnotesize [#1]}}
\let\uHom\relax\newcommand{\uHom}{\underline{\mathrm{Hom}}}
\let\Hom\relax\newcommand{\Hom}{\mathrm{Hom}}
\let\Cat\relax\newcommand{\Cat}{\mathrm{Cat}}
\let\Ens\relax\newcommand{\Ens}{\mathrm{Ens}}
\let\id\relax\newcommand{\id}{\mathrm{id}}
\let\Ob\relax\newcommand{\Ob}{\mathrm{Ob}\,}
\let\circled\relax\newcommand{\circled}[1]{\tikz[baseline=(c.base)]\node[draw,circle,inner sep=1pt](c){#1};}
\let\cercle\relax\newcommand{\cercle}[1]{\tikz[baseline=(c.base)]\node[draw,circle,inner sep=1pt](c){#1};}
\let\defeq\relax\newcommand{\defeq}{\overset{\text{déf}}{=}}
\let\upi\relax\newcommand{\upi}{\underline{\pi}}
\let\uCh\relax\newcommand{\uCh}{\underline{\mathrm{Ch}}}
\let\bDelta\relax\newcommand{\bDelta}{\boldsymbol{\Delta}}
\let\hra\relax\newcommand{\hra}{\lhook\joinrel\longrightarrow}

% Reader layout additions only; all edition slips in the body are retained.

% ---- from Der-ChXIIb-012_p001-005.tex
\printedpage{43}
% head: A. Grothendieck : Dérivateurs, Ch. XII, Édition des Universités de Montpellier II et Paris VII ; foot: Version 7 décembre 2002

\msp{page 59}

\section*{V\quad Sommes amalgamées et carrés $W$-cocartésiens dans $\mathrm{Cat}$}

\subsection*{1\quad Généralités sur $U \overset{i}{\hookrightarrow} X \overset{j}{\hookleftarrow} T$, et les sommes amalgamées.}

\subsubsection*{1.1\quad Données de recollement entre un ouvert et un fermé.}

Soit $X$ dans Cat, $U$ un ouvert dans $X$, $T$ le complémentaire fermé. (\textbf{NB} Les ouverts correspondent bien aux ouverts de $\mathrm{Ob}\,X$, topologisé par sa relation de préordre naturelle $x \leq y \iff \mathrm{Hom}_X(x,y) \neq \emptyset$, de même pour les fermés. Les topologies associées aux relations de préordre sont celles où toute intersection d'ouverts est encore ouverte, i.e. toute réunion de fermés est encore fermée. Donc les ouverts d'une telle topologie sont les fermés d'une topologie duale, correspondant à la relation de préordre opposée. La relation de préordre associée à la topologie est
\[
x \leq y \qquad\qquad \iff \qquad\qquad x \in \overline{\{y\}},
\]
la topologie associée au préordre sur $I$ est celle pour laquelle les ouverts sont les réunions d'ensembles de la forme $I/x$ (lesquels forment une base d'ouverts), ou encore les \textit{`cribles'} ($x \in U$, $y \leq x \Longrightarrow y \in U$), les fermés sont les réunions d'ensembles $x\backslash I$, c'est-à-dire les cocribles. Pour toute partie $A$ de $I$, il y a un plus petit ouvert contenant $A$, qui est aussi son plus petit voisinage, c'est la réunion des $I/x$ pour $x \in A$. Et dualement il y a un plus grand fermé contenu dans $A$, dont le complémentaire est le plus petit ouvert contenant $B = I\backslash A$, donc $\bigcup_{x\in B} I/x$. L'adhérence de $A$ est la réunion des adhérences $x\backslash I$ de ses points.)

\msp{page 60}

\dgrblock{dgrb507}{An open–closed cut remembers a profunctor}
\begin{dcomment}{exposition}{p59}{An open--closed cut remembers a profunctor}
Grothendieck begins with a cut of a category which permits arrows to cross in only one direction. This is the reason that the subsequent pushout calculation is manageable. An open subcategory $U\subset X$ is a full sieve: if an arrow has its target in $U$, its source belongs to $U$. The complementary full subcategory $T$ is therefore a cosieve, and there are no arrows from $T$ to $U$. Once the compositions internal to $U$ and $T$ are known, the remaining information is precisely the action of these compositions on the sets of arrows from $U$ to $T$.

There is a reversal in the displayed specialization convention on this page. With $x\leq y$ meaning that $\Hom_X(x,y)$ is nonempty and with the opens taken to be downsets, one has
\[
 \overline{\{x\}}=\{y\mid x\leq y\},\qquad
 x\leq y\ \Longleftrightarrow\ y\in\overline{\{x\}}.
\]
The edition really prints $x\in\overline{\{y\}}$. That formula is incompatible with its following description of opens as sieves; it is a source slip, not a transcription error. For the walking arrow $0\to1$, the opens are $\varnothing,\{0\},\{0,1\}$ and the closure of $\{0\}$ is the whole space. All the subsequent formulas for closure and boundary use the downset convention and are correct with it. In particular, the smallest open containing $A$ is its downward closure, while its ordinary closure is its upward closure. No antisymmetry is required: categorically isomorphic objects, and more generally objects connected by arrows in both directions, cannot be distinguished by these opens.

For completeness, the reconstruction in Proposition~1 can be written without a choice of compositions. Given
\[
 H:U^{\mathrm{op}}\times T\longrightarrow\Ens,
\]
form a category $\operatorname{Coll}(H)$ whose objects are the tagged disjoint union of those of $U$ and $T$, whose hom sets within each part are the original ones, whose hom sets from $T$ to $U$ are empty, and whose hom sets from $u\in U$ to $t\in T$ are $H(u,t)$. Define composition with $a:u'\to u$ by $h\mapsto H(a,1)(h)$ and composition with $b:t\to t'$ by $h\mapsto H(1,b)(h)$. The two operations commute by bifunctoriality. Every composable string contains at most one arrow crossing from $U$ to $T$, so these identities, together with associativity in the two pieces, prove associativity in the collage. The identities follow from the identity laws for $H$.

A functor preserving the two parts consists of functors $f_U,f_T$ and a natural transformation $H\to H'(f_U^{\mathrm{op}}\times f_T)$. The naturality squares are exactly the equations expressing preservation of composition involving a crossing arrow. This proves full faithfulness of the correspondence of Proposition~1; the collage construction proves essential surjectivity. With tagged objects there is a specified inverse up to canonical isomorphism, which explains the editors' preference for an equivalence over an unqualified isomorphism of categories. Today $H$ is called a profunctor and $\operatorname{Coll}(H)$ its collage. The point of the terminology is that a category can be assembled from its internal categories and a bimodule of arrows between them.
\end{dcomment}


La situation $(X,U,T)$ se reconstitue, à isomorphisme unique près, par la connaissance de

\begin{enumerate}[label={}]
\item[$1^\circ$)] Les deux catégories $U$, $T$.
\item[$2^\circ$)] Le bifoncteur $\displaystyle\left\{\begin{aligned}
&H : U^\circ \times T \longrightarrow \mathrm{Ens}\\[1ex]
&H(u,t) = \mathrm{Hom}_X(u,t).
\end{aligned}\right.$
\end{enumerate}

De plus, ces données $1^\circ$) $2^\circ$) en sont soumises à aucune restriction. On a, plus exactement % [sic] en sont -> ne sont, negation required by the sense

\textbf{Proposition 1.} \textit{Un isomorphisme} \texttt{[plutôt une équivalence de catégories]} \textit{entre la catégorie $\mathcal{T}$ dont les objets sont les triples $(X,U,T)$, avec $X \in \mathrm{Ob}\,\mathrm{Cat}$, $U \in \mathrm{Ouv}(X)$, $T \in \mathrm{Ferm}(X)$, $U$ et $T$ complémentaires l'un de l'autre, les morphismes de $(X,U,T)$ dans $(X',U',T')$ étant les foncteurs $f : X \longrightarrow X'$ qui induisent $U \longrightarrow U'$, $T \longrightarrow T'$ (d'où
\printedpage{44}
% head: A. Grothendieck : Dérivateurs, Ch. XII, Édition des Universités de Montpellier II et Paris VII ; foot: Version 7 décembre 2002
$U = f^{-1}(U')$, $T = f^{-1}(T')$), avec la catégorie $\mathcal{T}'$ des triples $(U,T,H)$, de deux catégories $U$, $T$ et d'un bifoncteur}
\[
H : U^\circ \times T \longrightarrow \mathrm{Ens}
\]
\hspace*{2em}\textit{(ou encore un préfaiscau $U^\circ \longrightarrow T^\vee$ sur $U$, donc un objet de} $\uHom(U^\circ,T^\vee)$\textit{)),} % [sic] préfaiscau -> préfaisceau, misspelling

\textit{un morphisme de $(U,T,H)$ dans $(U',T',H')$ étant un triple}
\[
(f_U, f_T, u) \qquad\qquad
\left\{\begin{array}{l}
\left.\begin{array}{llll}
f_U : & U & \longrightarrow & U' \\[1ex]
f_T : & T & \longrightarrow & T'
\end{array}\right\}\text{des foncteurs}\\[2ex]
u : H \longrightarrow H' \circ (f_U^\circ \times f_T),\\
\text{i.e. } H(u,t) \longrightarrow H'(f_U(u), f_T(t))\\
\text{un homomorphisme fonctoriel}\\
\text{(flèche dans } (U \times T^\circ)^\wedge \simeq \uHom(U^\circ,T^\vee)\text{)),}
\end{array}\right.
\]
\textit{composition des flèches évidente.}

Si $U \hookrightarrow X \hookleftarrow T$ est un triple, et
\[
f_{U'} : U' \longrightarrow U,\ f_{T'} : T' \longrightarrow T
\]
des flèches dans Cat, il y a une façon privilégiée de `recoller' $U'$ et $T'$, de façon à ce que l'objet $(X',U',T')$ s'envoie dans $(X,U,T)$ en induisant

\msp{page 61}

$f_{U'}$, $f_{T'}$:
\[
\begin{tikzcd}[column sep=large, row sep=large]
U' \arrow[r, hook, "i'"] \arrow[d, "f_{U'}"'] & X' \arrow[d, "f"] & T' \arrow[l, hook', "j'"'] \arrow[d, "f_{T'}"] \\
U \arrow[r, hook, "i"] & X & T, \arrow[l, hook', "j"']
\end{tikzcd}
\]
elle s'obtient à l'aide des bifoncteurs
\[
\left\{\begin{aligned}
&H'(u',t') \overset{\text{déf}}{=} H(f_{U'}(u'), f_{T'}(t'))\\
&\text{i.e. } H' = H \circ (f_{U'}^\circ \times f_{T'}).
\end{aligned}\right.
\]
Ce triple, i.e. cet objet de $\mathcal{T}$, représente le foncteur suivant sur $\mathcal{T}$ :
\[
(X'',U'',T'') \longmapsto \mathrm{Hom}_{\mathrm{Cat}}(U'',U) \times \mathrm{Hom}_{\mathrm{Cat}}(T'',T),
\]
i.e. on a pour tout objet $(X'',U'',T'')$ de $\mathcal{T}$ une bijection canonique, fonctorielle en cet objet
\[
\mathrm{Hom}_{\mathcal{T}}((X'',U'',T''),(X',U',T')) \simeq \mathrm{Hom}_{\mathrm{Cat}}(U'',U) \times \mathrm{Hom}_{\mathrm{Cat}}(T'',T). \tag*{$(*)$}
\]

\printedpage{45}
% head: A. Grothendieck : Dérivateurs, Ch. XII, Édition des Universités de Montpellier II et Paris VII ; foot: Version 7 décembre 2002

Nul doute qu'on a même
\[
\uHom_{\mathcal{T}}((X'',U'',T''),(X',U',T')) \simeq \uHom_{\mathrm{Cat}}(U'',U) \times \uHom_{\mathrm{Cat}}(T'',T), \tag*{$(**)$}
\]
tout aussi fonctoriel, dont $(*)$ se déduit en passant aux ensembles d'objets. On dit que le triple $(X',U',T')$ est le triple \textit{induit} par le triple $(X,U,T)$, via les foncteurs $f_{U'}$, $f_{T'}$.

\textbf{NB} Si $f_{U'}$, $f_{T'}$ sont pleinement fidèles, alors il en est de même de $f : X' \longrightarrow X$. Nous utiliserons cette opération de recollement privilégié surtout quand $U' = U$, $f_{U'} = \mathrm{id}_U$, ou inversement $T' = T$, $f_{T'} = \mathrm{id}_T$.

\subsubsection*{1.2\quad Relation avec le recollement des topos.}

La situation d'un triple
\[
(\mathcal{X},\mathcal{U},\mathcal{T})
\]
d'un topos $\mathcal{X}$, d'un ouvert $\mathcal{U}$ de $\mathcal{X}$, et d'un

\msp{page 62}

\dgrblock{dgrb508}{The induced collage has a relative universal property}
\begin{dcomment}{exposition}{p61}{The induced collage has a relative universal property}
The induced triple is meant to express change of the two categories on which the gluing data live. Its universal property has to remember a map to the original triple. Without that map the displayed formulas $(*)$ and $(**)$ are false; both formulas are reproduced faithfully from the edition.

Write $\mathcal T$ for the category of triples and $P:\mathcal T\to\Cat\times\Cat$ for $(X,U,T)\mapsto(U,T)$. Let $b=(f_{U'},f_{T'}):(U',T')\to(U,T)$ and let $X'$ be the collage of $b^*H$. For any triple $Y=(Y,V,S)$ the correct bijection is
\[
\begin{split}
 \Hom_{\mathcal T}(Y,X')\simeq
 &\Hom_{\mathcal T}(Y,X)\mathop{\times}_{
 \Hom(V,U)\times\Hom(S,T)}\\[-2pt]
 &\hspace{18mm}\bigl(\Hom(V,U')\times\Hom(S,T')\bigr).
\end{split}
\]
On the right one chooses a morphism $g:Y\to X$ and functors $a:V\to U'$, $c:S\to T'$ such that $g|_V=f_{U'}a$ and $g|_S=f_{T'}c$. The map on mixed hom sets of $g$ then has type
\[
 H_Y\longrightarrow H(f_{U'}a-,f_{T'}c-)
                 =(b^*H)(a-,c-).
\]
It defines a unique lift $Y\to X'$, and restriction of a lift recovers exactly these data. Thus $X'\to X$ is a cartesian arrow of $P$, not an object representing arbitrary pairs of functors to $U$ and $T$.

The same argument works on natural transformations. A transformation between two lifts is a pair of transformations on $V,S$ satisfying the naturality equations for every mixed arrow, and these are precisely the equations in the corresponding strict pullback of functor categories. This is the correctly typed version of $(**)$. One may express the statement in the language of a split fibration of gluing data, using the literal restriction operation on bifunctors; the corresponding assertion for unlabelled collages is invariant under the canonical isomorphisms.

The missing datum cannot be supplied by merely inserting primes on the right of the printed formulas. Take $U=T=U'=T'=e$, both change-of-base functors the identity, and $H=\{a,b\}$. Then $X'=X$ is the category with two parallel arrows from the open object to the closed object. There are two part-preserving functors from the walking arrow to $X'$, whereas there is only one pair of functors from its two singleton parts to $U',T'$. The choice of crossing arrow is precisely the information lost by the printed right-hand side. The following observation about full faithfulness is unaffected: on the two diagonal sorts of hom sets it follows from $f_{U'},f_{T'}$, and on mixed hom sets it follows from the defining equality $H'=b^*H$.
\end{dcomment}


sous-topos fermé complémentaire, est toute similaire (\textbf{NB} Si $\mathcal{X} = \mathrm{Top}(X)$, avec $X$ dans Cat, alors $\mathrm{Ouv}(\mathcal{X}) \simeq \mathrm{Ouv}(X)$, puisque les ouverts correspondent aux sous-faisceaux du faisceau final sur $\mathcal{X}$, i.e. du préfaisceau final sur $X$. Cette bijection est compatible avec le passage aux complémentaires.) La théorie des topos dit que les couples $(\mathcal{X},\mathcal{U},\mathcal{T})$ \texttt{[plutôt triples]} se reconstituent, à équivalence définie à isomorphisme unique près, à l'aide d'un triple
\[
\left\{\begin{aligned}
&\mathcal{U},\ \mathcal{T} \text{ deux topos}\\[1ex]
&\varphi : \mathrm{Fais}\,(\mathcal{U}) \longrightarrow \mathcal{T} \text{ un foncteur exact à gauche (`accessible'),}
\end{aligned}\right.
\]
ce foncteur $\varphi$ étant obtenu, quand on part de
\[
\mathcal{U} \xrightarrow{\;i\;} \mathcal{X} \xleftarrow{\;j\;} \mathcal{T},
\]
comme
\[
\varphi = j^*i_*.
\]
De plus, le triple $(\mathcal{U},\mathcal{T},\varphi)$ peut être choisi arbitrairement, avec la seule restriction (pratiquement anodine, car toujours satisfaite pour tous les foncteurs naturels qu'on rencontre) que $\varphi$ soit (en plus d'être exact à gauche) `accessible'.

Dans le cas actuel $\mathcal{X} = \mathrm{Top}(X)$, $\mathcal{U} = \mathrm{Top}(U)$, $\mathcal{T} = \mathrm{Top}(T)$, le foncteur $\varphi = j^*i_*$ est non seulement exact à gauche, mais même il commute aux $\varprojlim$ quelconques, puisqu'il en est ainsi (toujours) de $i_*$, et aussi (parce que le $j$ topossique est induit par le $j : T \longrightarrow X$ catégorique) de $j^*$. En fait, on
\printedpage{46}
% head: A. Grothendieck : Dérivateurs, Ch. XII, Édition des Universités de Montpellier II et Paris VII ; foot: Version 7 décembre 2002
\msp{page 63}
trouve ceci :

\dgrblock{dgrb509}{The presheaf case is already present in SGA 4}
\begin{dcomment}{context}{p62}{The presheaf case is already present in SGA 4}
Grothendieck's comparison with toposes recovers a particularly exact antecedent. In \citet[Expos\'e~IV, Theorem~9.5.4]{SGA4I1972}, an open subtopos and its closed complement reconstruct the ambient topos through the accessible left exact gluing functor. Exercise~9.5.10 of the same expos\'e then treats presheaf toposes: its parts~(a)--(c) identify opens with sieves, introduce the cross-hom bifunctor, and characterize the presheaf case by preservation of products by the gluing functor. Thus Proposition~2 is a return to a construction already formulated there, now serving a calculation of categorical pushouts and their homotopy properties. The letters $i,j$ are interchanged between the SGA convention and the present pages.

The category of sheaves in the glued topos is the category of triples
\[
 (F,G,\theta),\qquad F\in\mathcal U,\quad G\in\mathcal T,
 \quad\theta:G\longrightarrow\varphi(F).
\]
The comparison arrow records how a section on the closed side approaches the open side. For $X=\operatorname{Coll}(H)$, a presheaf $E$ on $X$ restricts to $F=E|_U$ and $G=E|_T$. Its values on the mixed arrows are maps
\[
 G(t)\times H(u,t)\longrightarrow F(u)
\]
natural in $u,t$. Transposing in the first variable produces $G(t)\to\Hom_{U^\wedge}(H(-,t),F)$, exactly the required gluing arrow. Conversely these maps reconstruct the restriction maps of $E$ on all arrows; their naturality supplies the mixed composition laws. This directly explains why the ordinary categorical collage and the topos gluing construction agree. The comparison is not an analogy between two shapes of diagrams: both constructions encode the same mixed restriction maps.
\end{dcomment}


\textbf{Proposition 2 :} \textit{Considérons une situation de recollement de deux topos $\mathcal{U}$, $\mathcal{T}$ par $\varphi : \mathrm{Fais}\,\mathcal{U} \longrightarrow \mathrm{Fais}\,\mathcal{T}$, foncteur exact à gauche. Supposons que $\mathcal{U} = \mathrm{Top}(U)$, $\mathcal{T} = \mathrm{Top}(T)$, avec $U$ et $T$ dans $\mathrm{Cat}$. Pour que le topos recollé $\mathcal{X}$ soit équivalent à un topos $\mathrm{Top}(X)$, avec $X$ dans $\mathrm{Cat}$, il faut et il suffit que le foncteur de recollement $\varphi$ commute aux $\varprojlim$ quelconques, ou ce qui revient au même, qu'il ait un adjoint à gauche :}
\[
\psi : \underbrace{\mathrm{Fais}\,\mathcal{T}}_{= T^\wedge} \longrightarrow \underbrace{\mathrm{Fais}\,\mathcal{U}}_{= U^\wedge}.
\]

On s'attend dès lors que la donnée de recollement $\varphi$, ou ce qui revient au même le foncteur $\psi$ ayant un adjoint à droite, i.e. commutant aux $\varinjlim$ quelconques, revient à celle d'une donnée de recollement $H$
\[
H : U^\circ \times T \longrightarrow \mathrm{Ens}
\]
pour $U$ et $T$. Mais on sait bien que les foncteurs $\psi : T^\wedge \longrightarrow \mathcal{C}$ dans une catégorie $\mathcal{C}$ avec $\varinjlim$ quelconques, qui commutent auxdites $\varinjlim$, forment une catégorie équivalente (par le foncteur $\psi \longmapsto \psi|T(\subseteq T^\wedge)$) à celle des foncteurs
\[
\psi_0 : T \longrightarrow \mathcal{C} \quad\ (= \underbrace{U^\wedge}_{\uHom(U^\circ,\mathrm{Ens})} \text{ en l'occurrence}),
\]
\msp{page 64}
et la donnée d'un tel foncteur $\psi_0$ revient bien à celle d'un $H$. Donc

\textbf{Corollaire.} \textit{La catégorie des données de recollement topossique $\varphi : U^\wedge \longrightarrow T^\wedge$ entre} $\mathrm{Top}(U)$, $\mathrm{Top}(T)$\textit{, telles que le topos recollé soit un `topos} paradisiaque\textit{', i.e. équivalent à un topos de la forme $X^\wedge$ (avec $X$ dans} Cat\textit{), est équivalente à celle des données de recollement $H : U^\circ \times T \longrightarrow$} Ens \textit{entre $U$ et $T$. La donnée topossique $\varphi$ associée à $H$ est le foncteur adjoint à droite de $T^\wedge \longrightarrow U^\wedge$ déduit de $T \xrightarrow{H} U^\wedge$ par prolongement commutant aux $\varinjlim$, ou encore elle est donnée par}
\[
\varphi(\underbrace{F}_{\in U^\wedge})(t) \simeq \mathrm{Hom}_{U^\wedge}(H(t),F),
\]
\textit{où dans le second membre $H$ est interprété comme un foncteur $T \longrightarrow U^\wedge$. Le topos recollé par $\varphi$ est canoniquement équivalent à} $\mathrm{Top}(X)$\textit{, où $X$ est l'objet de} Cat \textit{obtenu par recollement à l'aide de $H$.}

Du point de vue de la théorie des faisceaux, cela donne l'interprétation de $H$, comme exprimant le foncteur
\[
\varphi = j^*i_* : U^\wedge \longrightarrow T^\wedge
\]
par la formule précédente.

\printedpage{47}
% head: A. Grothendieck : Dérivateurs, Ch. XII, Édition des Universités de Montpellier II et Paris VII ; foot: Version 7 décembre 2002

\msp{page 65}

\dgrblock{dgrb510}{The adjunction which characterizes the presheaf gluing}
\begin{dcomment}{exposition}{p64}{The adjunction which characterizes the presheaf gluing}
The extra preservation of limits isolates those gluing functors that come from actual sets of arrows. Let $H_t=H(-,t)\in U^\wedge$. Its colimit-preserving extension from the Yoneda image of $T$ is
\[
 \psi(G)=\int^{t\in T}G(t)\times H_t,
 \qquad
 \varphi(F)(t)=\Hom_{U^\wedge}(H_t,F).
\]
Here the coend is an ordinary colimit in $U^\wedge$. A map $\psi(G)\to F$ is a family of maps $G(t)\times H_t\to F$ compatible with the arrows of $T$. By the ordinary set-valued adjunction between product and hom, this is exactly a natural transformation $G\to\varphi(F)$. This proves $\psi\dashv\varphi$ and the displayed formula of the corollary. It also proves that $\varphi$ preserves all small limits.

Conversely, suppose that an accessible gluing functor $\varphi:U^\wedge\to T^\wedge$ preserves all small limits. The adjoint functor theorem for locally presentable categories supplies a left adjoint $\psi$. Set $H_t=\psi(y_Tt)$. Every presheaf is the colimit of its representables indexed by its category of elements; since $\psi$ is a left adjoint, this determines $\psi$ from $H$. The preceding adjunction then identifies its right adjoint with $\varphi$. Applying the explicit construction of the preceding block identifies the glued topos with $\operatorname{Coll}(H)^\wedge$. This proves the sufficient direction of Proposition~2 with the size assumption made explicit. The accessibility assumption belongs to the ambient gluing problem; it should not be silently discarded in an abstract claim that every limit-preserving functor has a left adjoint.

For the converse direction of the proposition, let the glued topos be equivalent to $C^\wedge$. Its specified open subtopos corresponds under this equivalence to a subterminal presheaf. Such a presheaf takes the value either $\varnothing$ or a singleton at each object, and its support is a sieve $C_U\subset C$. The complementary full cosieve $C_T$ presents the closed complement. Restriction to $C_T$ preserves all limits, as does the right adjoint to restriction to $C_U$. Their composite is the transported gluing functor, so $\varphi$ preserves all limits. The argument uses the specified open and closed pieces but does not require the given presentations $U,T$ to coincide literally with $C_U,C_T$. Equivalences of presheaf categories preserve the limiting operations involved. This is precisely the content of the presheaf case of the SGA exercise just cited.

There is a variance convention in the word ``category'' in the corollary. For fixed $U,T$, a transformation $H\to K$ induces $\psi_H\to\psi_K$, but it induces a transformation $\varphi_K\to\varphi_H$ on the right adjoints. Thus the asserted equivalence of categories of gluing data uses reversed transformations on the right-adjoint side; with ordinary pointwise transformations it is an antiequivalence. The object classification and every displayed objectwise formula remain unchanged.

A useful test is $U=T=e$ and $H=S$ a set. The collage has two objects and one crossing arrow for each element of $S$. The functors are $\psi(G)=G\times S$ and $\varphi(F)=F^S$. A presheaf on the collage is a pair of sets $F,G$ with an $S$-indexed family of maps $G\to F$. If $S=\varnothing$, there is no compatibility datum and the presheaf topos is the product of two copies of $\Ens$. If $S$ is a singleton, it is the category of arrows of sets. These limiting cases also verify the direction $G\to\varphi(F)$ of the gluing arrow.
\end{dcomment}


\subsubsection*{1.3\quad Flèches dans $\mathrm{Cat}$, de source ou but un objet recollé $(X,U,T)$.}

Quand $(X,U,T)$ est donné par une donnée de recollement
\[
U,\ T,\ H : U^\circ \times T \longrightarrow \mathrm{Ens},
\]
les flèches
\[
f : X' \longrightarrow X
\]
dans Cat se décrivent immédiatement ainsi : la donnée de $f$ définit d'abord
\[
U' = f^{-1}(U), \qquad\qquad T' = f^{-1}(T),
\]
couple d'un ouvert de $X'$ et du fermé complémentaire, de sorte que cette donnée équivaut à celle de

\begin{enumerate}[label={}]
\item[$1^\circ$)] un couple $(U',T')$ dans $X'$
\item[$2^\circ$)] un morphisme (dans $\mathcal{T}$) de $(X',U',T')$ dans $(X,U,T)$, ce qui équivaut par 1.1 à un triple
\[
\left\{\begin{array}{rccl}
f_{U'} : & U' & \longrightarrow & U\\[1ex]
f_{T'} : & T' & \longrightarrow & T\\[1ex]
u : & H_{U',T'} & \longrightarrow & H \circ (f_{U'}^\circ \times f_{T'}),
\end{array}\right.
\]
où on a posé
\[
\left\{\begin{aligned}
&H_{U',T'} : U^{\prime\circ} \times T' \longrightarrow \mathrm{Ens}\\[1ex]
&H_{U',T'}(u',t') = \mathrm{Hom}_{X'}(u',t'),
\end{aligned}\right.
\]
i.e. $H_{U',T'}$ est la donnée de recollement qui conrrespond au triple $(X',U',T')$ dans $\mathcal{T}$. % [sic] conrrespond -> correspond, misspelling
\end{enumerate}

\msp{page 66}

Nous serons plus intéressés par la situation inverse, i.e. par la question de déterminer les foncteurs
\[
\underbrace{X}_{= X_{U,T,H}} \quad \longrightarrow X'
\]
de $X$, obtenu par $(U,T,H)$, dans un objet quelconque de Cat. On trouve tautologiquement

\textbf{Proposition 3.} $\mathrm{Hom}_{\mathrm{Cat}}(X,X')$ \textit{est en correspondance bijective canonique, fonctorielle en $X'$ dans} Cat\textit{, avec l'ensemble des triples}

% ---- from Der-ChXIIb-012_p006-010.tex
\printedpage{48}
% head: A. Grothendieck : Dérivateurs, Ch. XII, Édition des Universités de Montpellier II et Paris VII
% foot: Version 7 décembre 2002

\[
\left\{
\begin{aligned}
&\left.\begin{aligned} f_U &: U \longrightarrow X' \\ f_T &: T \longrightarrow X' \end{aligned}\right\}\text{flèches dans Cat}\\
&u : H \longrightarrow H_{X'}(f_U^{\circ}\times f_T),\ \text{flèche dans } \uHom(U^{\circ}\times T,\Ens),
\end{aligned}
\right.
\tag*{\texttt{[}$(*)$\texttt{]}}
\]
\textit{où pour tout $X'$ dans} Cat\textit{, on a posé}
\[
H_{X'} = ((x,y)\longmapsto \Hom_{X'}(x,y)) : X'^{\circ}\times X' \longrightarrow \Ens.
\]

Sûrement on aura même un isomorphisme de catégories
\[
\uHom(X,X') \simeq \text{catégorie des triples } (f_U,f_T,u) \text{ comme dans } (*).
\]

\subsection*{1.4\quad Construction de sommes amalgamées à partir d'une immersion ouverte où fermée.} % [sic] où fermée -> ou fermée, conjonction « ou » et non adverbe « où »

Pour une construction d'une sorte de somme amalgamée, en termes d'un triple $(X,U,T)$ $+\ \ldots$, \msp{page 67} qui soit symétrique en $U$, $T$, j'ai envie de regarder d'abord la situation ‘mixte’, d'un diagramme en traits pleins
\[
\begin{tikzcd}
U \arrow[r, hook, "i"] \arrow[d, "f_U"'] & X \arrow[d, dotted, "f"] & T \arrow[l, hook', "j"'] \arrow[d, "f_T"] \\
U' \arrow[r, hook, dotted, "i'"] & X' & T', \arrow[l, hook', dotted, "j'"']
\end{tikzcd}
\tag*{$(*)$}
\]
qu'on aimerait compléter ‘canoniquement’, si faire se peut, en un morphisme de triples (i.e. en une flèche dans $\mathcal{T}$). L'idée naturelle, c'est que le $X'$ qui s'y insère serait la limite inductive, dans Cat, du diagramme initial. Si $T'=T$, $f_T=\id_T$, cela signifie donc que $X'$ s'identifie à la somme amalgamée $X\amalg_U U'$, et inversement si $U'=U$, $f_U=\id_U$.

\dgrblock{dgrb511}{Functors out of a collage include the mixed naturality equation}
\begin{dcomment}{exposition}{p66}{Functors out of a collage include the mixed naturality equation}
The next construction will be recognized by mapping out of it. Proposition~3 is therefore the correct place to retain natural transformations, even though its first formulation only mentions sets of functors. A functor $\operatorname{Coll}(H)\to Y$ consists of $a:U\to Y$, $b:T\to Y$ and maps
\[
 \eta_{u,t}:H(u,t)\longrightarrow\Hom_Y(a(u),b(t))
\]
compatible with precomposition in $U$ and postcomposition in $T$. The argument constructing the collage proves both existence and uniqueness of the functor from these data.

A morphism from $(a,b,\eta)$ to $(\bar a,\bar b,\bar\eta)$ is a pair of natural transformations $r:a\Rightarrow\bar a$, $s:b\Rightarrow\bar b$ satisfying, for every $h\in H(u,t)$,
\[
 s_t\,\eta_{u,t}(h)=\bar\eta_{u,t}(h)\,r_u.
\]
The components $r_u$ and $s_t$ define a natural transformation on the collage because its naturality is already known on arrows within either part, while this equation is its naturality on a crossing arrow. Conversely restriction recovers the pair and the equation. Composition and identities are componentwise. This gives the claimed isomorphism of functor categories, rather than merely an equivalence on their objects.

For a functor into a collage the procedure is different in presentation but just as exact. Pulling back the open part specifies a sieve in the source, and its complement specifies the cosieve; these are forced by the functor. The residual data are then a morphism of profunctors. This is why Proposition~3 does not require the images of $a$ and $b$ to be disjoint: the target is arbitrary, and a functor out of the collage may identify objects from the two parts. The universal property of the forthcoming ordinary pushout must allow precisely these functors.
\end{dcomment}


Désignons par $X'$ la somme amalgamée \texttt{[plutôt la limite inductive]}, et explicitons, pour $X''$ quelconque dans Cat, l'ensemble des flèches $g$ de $X'$ dans $X''$, comme foncteur en $X''$. C'est l'ensemble des triples
\[
(g_X, g_{U'}, g_{T'}) \qquad
\left\{\begin{aligned} g_X &: X \longrightarrow X'' \\ g_{U'} &: U' \longrightarrow X'' \\ g_{T'} &: T' \longrightarrow X''\end{aligned}\right.
\qquad \text{avec} \quad
\begin{aligned} g_X\circ i &= g_{U'} f_U \\ g_X \circ j &= g_{T'} f_T. \end{aligned}
\tag*{$(*)$}
\]
\printedpage{49}
% head: A. Grothendieck : Dérivateurs, Ch. XII, Édition des Universités de Montpellier II et Paris VII
% foot: Version 7 décembre 2002
\[
\begin{tikzcd}
U \arrow[r, hook, "i"] \arrow[d, "f_U"'] & X \arrow[dd, "g_X"'] & T \arrow[l, hook', "j"'] \arrow[d, "f_T"] \\
U' \arrow[dr, "g_{U'}"'] & & T' \arrow[dl, "g_{T'}"] \\
& X'' &
\end{tikzcd}
\]

Notons que pour un tel triple, $g_X$ est partiellement connu déjà en termes de $g_{U'}$, $g_{T'}$, par ses restrictions $g_X\circ i$ et $g_X\circ j$ à $U$ et $T$ respectivement, par les formules
\[
g_U = g_{U'}f_U, \qquad g_T = g_{T'}f_T.
\]

\msp{page 68} Donc la donnée d'un triple $(*)$ équivaut à celle d'un triple
\[
(g_{U'}, g_{T'}, g_X) \qquad
\left\{\begin{aligned}
&g_X \text{ une flèche } X\longrightarrow X'' \text{ correspondant à}\\
&g_U = g_{U'}\circ f_U \text{ et } g_T = g_{T'}\circ f_T \text{ déjà donnés (via } g_{U'}\text{, } g_{T'}\text{),}
\end{aligned}\right.
\]
donc en vertu de 1.3. cela correspond à
\[
(g_{U'}, g_{T'}, u) \qquad
\left\{\begin{aligned}
g_{U'} &: U' \longrightarrow X''\\
g_{T'} &: T' \longrightarrow X''\\
u &: H \longrightarrow \underbrace{H_{X''}\left((g_{U'}\circ f_U)^{\circ}\times(g_{T'}\circ f_T)\right)}_{(u,t)\longmapsto \Hom_{X''}(g_{U'}f_U(u),g_{T'}f_T(t))}.
\end{aligned}\right.
\tag*{$(*)$}
\]
On va interpréter la donnée de $u$, en termes d'un foncteur
\[
F^{*} = (H'\longmapsto H'\circ F) : \uHom(U'^{\circ}\times T',\Ens) \longrightarrow \uHom(U^{\circ}\times T,\Ens)
\tag*{$(**)$}
\]
où
\[
F : U^{\circ}\times T \longrightarrow U'^{\circ}\times T' \qquad \text{est donné par } F = f_U^{\circ}\times f_T.
\]
On a un objet $H''=H''_{g_{U'},g_{T'}}$ de la catégorie source de $(**)$ déterminé en termes du couple $(g_{U'},g_{T'})$, par
\[
H''(u',t') = \Hom_{X''}(g_{U'}(u'),g_{T'}(t')),
\]
et le bifoncteur but de la flèche $u$ dans $(*)$ s'interprète comme
\[
F^{*}(H'') = H''\circ F : (u,t)\longmapsto \underbrace{H''(f_U(u),f_T(t))}_{=\Hom(g_{U'}f_U(u),g_{T'}f_T(t))} \quad .
\]
Ainsi, $u$ s'interprète comme une flèche
\[
u : H \longrightarrow F^{*}(H'').
\]
Or le foncteur $F^{*}$ a un adjoint à gauche $F_{!}$, i.e. on a
\[
\Hom(H,F^{*}(H'')) \simeq \Hom(F_{!}(H),H'').
\tag*{$(*{*}*)$}
\]

\printedpage{50}
% head: A. Grothendieck : Dérivateurs, Ch. XII, Édition des Universités de Montpellier II et Paris VII
% foot: Version 7 décembre 2002
\msp{page 69} avec
\[
F_{!} : \uHom(U^{\circ}\times T,\Ens) \longrightarrow \uHom(U'^{\circ}\times T',\Ens).
\]
Posons
\[
H' = F_{!}(H) : U'^{\circ}\times T' \longrightarrow \Ens,
\]
c'est un bifoncteur bien déterminée par la seule donnée du diagramme initial $(*)$ p. 67, ou ce qui revient au même, en termes de % [sic] déterminée -> déterminé, accord avec « bifoncteur » (masculin)
\[
(U,T,H,f_U : U\longrightarrow U', f_T : T\longrightarrow T').
\]

\dgrblock{dgrb512}{Transporting the crossing arrows by a Kan extension}
\begin{dcomment}{exposition}{p68}{Transporting the crossing arrows by a Kan extension}
The two categories $U,T$ can now be changed independently. Their new internal arrows are already present in $U',T'$. What must be supplied are all composites crossing from the first new category to the second, subject to the relations forced by old compositions. This is what the left Kan extension supplies.

Put $F=f_U^{\mathrm{op}}\times f_T$. The ordinary pointwise formula can be written in a form that displays every generator:
\[
 H'(u',t')=
 \int^{u\in U,\,t\in T}
 \Hom_{U'}(u',f_Uu)\times H(u,t)
                  \times\Hom_{T'}(f_Tt,t').
\]
The notation means the quotient of the disjoint union of these products by the relations for composition in $U$ and $T$. An element is represented by a triple
\[
 u'\xrightarrow{\alpha}f_Uu,\qquad
 h\in H(u,t),\qquad f_Tt\xrightarrow{\beta}t'.
\]
For $a:u_0\to u_1$ in $U$ and $b:t_0\to t_1$ in $T$, the defining relations are
\[
 [f_U(a)\alpha,h,\beta]=[\alpha,H(a,1)h,\beta],
 \qquad
 [\alpha,H(1,b)h,\beta]=[\alpha,h,\beta f_T(b)],
\]
whenever the displayed terms are typed. This is the coend formula for $\operatorname{Lan}_F H$; the equivalent comma-category formula is the pointwise Kan-extension theorem \citep[Theorem~6.2.1]{Riehl2016}.

To prove the pushout assertion directly, form $C=\operatorname{Coll}(H')$ and send a crossing arrow $h$ of $X$ to $[1,h,1]$. The relations make this a functor $X\to C$, and the inclusions of $U',T'$ agree with it on $U,T$. Now take compatible functors $X\to Y$, $U'\to Y$, $T'\to Y$. On the generator above there is only one possible extension: it must take the value $\beta\,h\,\alpha$, after applying the respective functors. The two relations hold because the original functors preserve composition. Thus there is a unique functor $C\to Y$ inducing all three given functors. This proves the strict colimit universal property, and also identifies its canonical map from $X$.

Since $C$ was constructed as a collage, $U'$ is a full sieve and $T'$ its full complementary cosieve. Their internal hom sets have not been enlarged by the quotient: no string can leave the open part and subsequently return to it. This supplies the structural content of Theorem~1 beyond existence of a colimit in $\Cat$. It also clarifies the distinction from the earlier induced triple: restriction of $H$ gave a cartesian lift, whereas this left Kan extension gives the universal forward transport. The labels ``pullback'' and ``pushout'' refer to these different universal properties, not to a formal substitution of primes in the same formula.
\end{dcomment}


Et on voit que la donnée d'une flèche $X'\longrightarrow X''$, $X''$ objet variable de Cat, équivaut à celle de
\[
g_{U'} : U'\longrightarrow X'',\ g_{T'} : T'\longrightarrow X'',\ u : H'\longrightarrow H'' = H_{X''}(g_{U'}^{\circ}\times g_{T'}),
\]
c'est à dire, en vertu de 1.3., à celle d'une flèche
\[
X'_0 \longrightarrow X'',
\]
où $X'_0$ désigne l'objet de Cat déduit de $U'$, $T'$ par recollement au moyen de $H'$. Cela montre donc que la limite inductive cherchée n'est autre que $X'_0$. On a ainsi trouvé le

\textbf{Théorème 1.} \textit{La limite inductive $X'$ dans} Cat \textit{du diagramme} $(*)$ \textit{page 67, défini en termes d'un triple $(X,U,T)$ dans $\mathcal{T}$ et de flèches $f_U : U\longrightarrow U'$, $f_T : T\longrightarrow T'$ dans} Cat\textit{, existe (sans mérite)} \msp{page 70} \textit{et les flèches canoniques}
\[
U' \overset{i'}{\hookrightarrow} X' \overset{j'}{\hookleftarrow} T'
\]
\textit{sont des immersions ouvertes et fermées complémentaires l'une de l'autre. La donnée de recollement qui définit $X'$ via $U'$, $T'$ est donnée par}
\[
H' = (\underbrace{f_U^{\circ}\times f_T}_{F})_{!}(H) : U'^{\circ}\times T' \longrightarrow \Ens,
\]
\textit{où $H$ est la donnée de recollement qui définit $X$, i.e. $H : U^{\circ}\times T\longrightarrow$} Ens \textit{est le bifoncteur $(u,t)\longmapsto \Hom_X(u,t)$.}

Ainsi on a, pour $u'$ dans $U'$ et $t'$ dans $T'$
\[
\Hom_{X'}(u',t') = F_{!}(H)(u',t'),
\]
et il reste à expliciter $F_{!}(H)$. Ici $F=f_U^{\circ}\times f_T$ est un foncteur covariant sur $U^{\circ}\times T$, on peut l'interpréter comme un préfaisceau sur $U\times T^{\circ}$, et écrire dès lors
\[
\underbrace{F_{!}(H)(u',t')}_{\Hom_{X'}(u',t')} = \varinjlim_{(u',t')\backslash U\times T^{\circ}} \underbrace{\Hom_X(u,t)}_{H(u,t)} \ .
\tag*{$(*)$}
\]

\printedpage{51}
% head: A. Grothendieck : Dérivateurs, Ch. XII, Édition des Universités de Montpellier II et Paris VII
% foot: Version 7 décembre 2002
La catégorie $(u',t')\backslash U\times T^{\circ}$ sur laquelle on prend la $\varinjlim$ s'explicite comme un produit
\[
(u',t')\backslash U\times T^{\circ} \simeq (u'\backslash U)\times(T/t')^{\circ}.
\]
Les applications
\[
\Hom_X(u,t)\longrightarrow \Hom_{X'}(u',t') \qquad (=F_{!}(H)(u',t'))
\]
\msp{page 71} déduites de cette formule ne sont autres que les suivantes : $u$, $t$ sont données comme objets de $u'\backslash U$, $T/t'$, i.e. avec des flèches
\[
u'\xrightarrow{\alpha} f_U(u), \qquad f_T(t)\xrightarrow{\beta} t',
\]
d'où
\[
\Hom_X(u,t) \xrightarrow{\text{via } f=f_X:X\to X'} \Hom_{X'}(\underbrace{f(u)}_{f_U(u)},\underbrace{f(t)}_{f_T(t)}) \xrightarrow{\text{via } \alpha,\ \beta} \Hom_{X'}(u',t').
\]

Notons pour mémoire que pour tout objet $Z$ de Cat, et tout préfaisceau $G$ sur $Z$, on a la formule
\[
\varinjlim_{Z} G \simeq \pi_0(Z/G),
\]
où $Z/G$ s'interprète en termes de l'immersion canonique $Z\longrightarrow Z^{\wedge}$ $(G\in \Ob Z^{\wedge})$. C'est donc une catégorie fibrée sur $Z$ à fibres discrètes, et on a en fait un isomorphisme entre la sous-catégorie pleine de $\Cat/Z$ formée de telles catégories fibrées à fibres discrètes, avec $Z^{\wedge}$. Dans ce dictionnaire parfait entre préfaisceaux d'ensembles et catégories fibrées à fibres discrètes, le foncteur $\varinjlim$ sur $Z^{\wedge}$ correspond simplement au foncteur $\pi_0$ sur la catégorie des objets au dessus de $Z$ dans Cat, lui-même induit par le foncteur (indépendant de $Z$)
\[
\Cat \xrightarrow{\pi_0} \Ens.
\]

\msp{page 72} On peut donc réécrire la formule fondamentale $(*)$ p. 70, en termes du foncteur $H$ définissant le recollement $X$ de $U$ et de $T$, par
\[
\boxed{\underbrace{\Hom_{X'}(u',t')}_{H'(u',t')} \simeq \pi_0(\underbrace{(u',t')\backslash U\times T^{\circ}/H}_{((u'\backslash U)\times(T/t')^{\circ})/H}).}
\tag*{$(*)$}
\]

Par la suite, nous n'aurons sans doute qu'à utiliser le cas où on a soit $T'=T$, $f_T=\id_T$, soit l'inverse. Nous travaillerons surtout dans le premier cas, en déduisant ce dont nous aurons besoin pour le cas dual, par principe de dualité (automorphisme $X\longmapsto X^{\circ}$ de la
\printedpage{52}
% head: A. Grothendieck : Dérivateurs, Ch. XII, Édition des Universités de Montpellier II et Paris VII
% foot: Version 7 décembre 2002
catégorie Cat). Dans ce cas, dans la formule fondamentale $(*)$ l'élément $t'$ de $T'$ s'écrit comme un élément $t$ de $T$, et la formule se retranscrit comme
\[
\boxed{\Hom_{X'}(u',t) \simeq \pi_0(\underbrace{u'\backslash U/t}_{\text{catégorie des triples } (u,\alpha,\beta'),\ \left\{\begin{aligned} &u\in \Ob U\\ &\alpha : u\longrightarrow t\\ &\beta' : u'\longrightarrow f_U(u)\end{aligned}\right.}).}
\tag*{$(**)$}
\]

C'est la formule que j'avais trouvée directement dans ce cas, sans me soucier encore de faire varier $T$ en même temps que $U$. Je voudrais faire le lien entre les deux formules $(*)$ et $(**)$, par voie topologique, donc établir une bijection canonique, fonctorielle etc.
\[
\pi_0(u'\backslash U/t) \xrightarrow{\sim} \pi_0(((u'\backslash U)\times(T/t)^{\circ})/H).
\]

\msp{page 73} On a sûrement un foncteur
\[
u'\backslash U/t \longrightarrow ((u'\backslash U)\times(T/t)^{\circ})/H
\]
qui associe à un triple
\[
(u,\alpha,\beta') : u\in \Ob U,\ \alpha : u\longrightarrow t,\ \beta' : u'\longrightarrow f_U(u)
\]
le triple $(\bar u,\bar t,\xi)$ $(\xi\in H(\bar u,\bar t))$ dans le second membre, avec $\bar u$ dans $u'\backslash U$, $\bar t$ dans $T/t$, défini ainsi :
\[
\bar u = (u,\beta'),\ \bar t = (t,\id_t), \underbrace{\xi}_{=\alpha} \in H(\bar u,\bar t) = \Hom(u,t).
\]
En somme
\[
u'\backslash U/t = \text{fibre de } ((u'\backslash U)\times(T/t)^{\circ})/H \text{ en l'objet initial } \bar t \text{ de } (T/t)^{\circ}
\]
par la projection naturelle, flèche composée
\[
\begin{tikzcd}
((u'\backslash U)\times(T/t)^{\circ})/H \arrow[d] \\
(u'\backslash U)\times(T/t)^{\circ} \arrow[d, "\pi"] \\
(T/t)^{\circ}.
\end{tikzcd}
\]

\dgrblock{dgrb513}{Why connected components give the ordinary hom set}
\begin{dcomment}{exposition}{p72}{Why connected components give the ordinary hom set}
Grothendieck wants a topological description of a purely ordinary quotient. For a presheaf $G:Z^{\mathrm{op}}\to\Ens$, the category of elements $Z/G$ has objects $(z,g)$ and arrows
\[
 (z,g)\longrightarrow(z',g')\quad\text{given by}\quad
 a:z\longrightarrow z',\qquad G(a)(g')=g.
\]
The colimit of $G$ is the disjoint union of its values modulo the identifications $G(a)(g')\sim g'$. These are exactly the identifications imposed on objects when one passes to connected components of $Z/G$. Hence
\[
 \mathop{\mathrm{colim}}_{Z^{\mathrm{op}}}G\simeq\pi_0(Z/G).
\]
The source writes $\varinjlim_Z G$ for this presheaf colimit. The opposite belongs in the modern indexing category; it is not a correction to the numerical value of the formula.

The same convention explains the indexing category in the preceding page. If $H$ is regarded as a covariant functor on $U^{\mathrm{op}}\times T$, its ordinary left Kan extension is indexed by
\[
 F/(u',t')\simeq (u'\backslash f_U)^{\mathrm{op}}\times(f_T/t').
\]
Equivalently, regard $H$ as a presheaf on $U\times T^{\mathrm{op}}$ and take its presheaf colimit over the opposite category
\[
 (u'\backslash f_U)\times(f_T/t')^{\mathrm{op}}.
\]
This is the convention used in the edition. Its sentence ``Ici $F=\cdots$ ... on peut l'interpr\'eter comme un pr\'efaisceau'' contains a harmless but real letter slip: it is $H$, not the indexing functor $F$, that is the set-valued presheaf.

When $T'=T$, the resulting objects can be compressed to factorizations
\[
 (u,\alpha:u'\to p(u),\gamma:u\to t).
\]
An arrow from $(u,\alpha,\gamma)$ to $(v,\alpha',\gamma')$ is $a:u\to v$ with $p(a)\alpha=\alpha'$ and $\gamma'a=\gamma$. Its connected component is the crossing arrow of the pushout represented by $\gamma\alpha$. Thus the factorization category keeps track of all presentations of a new arrow and all relations between those presentations. Taking $\pi_0$ intentionally forgets its higher homotopy. The later asphericity conditions ask precisely when this forgetting is harmless for the homotopy problem, and do not follow from the existence of the ordinary colimit alone.
\end{dcomment}


Il faut voir que l'inclusion de la fibre donne une bijection sur les $\pi_0$, ou ce qui revient au même, que l'inclusion de la fibre $u'\backslash U$ \textit{de} $\pi$ en l'objet initial $(t,\id_t)$ de $(T/t)^{\circ}$ donne une bijection sur les $\varinjlim$ du foncteur $H$. Donc l'inclusion a intérêt à être asphérique, ou du moins co-cofinalle. % [sic] co-cofinalle -> co-cofinale, orthographe

% ---- from Der-ChXIIb-012_p011-015.tex
% ===================== page 11 =====================
\printedpage{53}
% head: A. Grothendieck : Dérivateurs, Ch. XII, Édition des Universités de Montpellier II et Paris VII

\msp{page 74}
Mais l'inclusion d'un objet initial est toujours asphérique, et le reste par produit avec un objet quelconque (i.e. $u'\backslash U$) de Cat (et même par tout changement de base lisse \dots), d'où la conclusion.

Dans la suite, nous nous intéressons surtout à la situation de sommes amalgamées (où on change de notation un peu, à présent)
\[
\begin{tikzcd}
U \arrow[r, hook, "i"] \arrow[d, "p"'] & X \arrow[d, "q"'] & T \arrow[l, hook', dotted, "j"'] \arrow[d, dotted] \\
U' \arrow[r, hook, "i'"] & X' & T' \simeq T. \arrow[l, hook', dotted]
\end{tikzcd}
\tag*{$(*)$}
\]

Le fermé $T'$ complémentaire de $U'$ dans $X'$ est alors canoniquement isomorphe au fermé $T$ complémentaire de $U$ dans $X$, l'isomorphisme étant induit par $q$. Souvent nous identifierons $T'$ et $T$, ce qui est licite en regardant $X'$ comme déduit du recollement de $U'$ et de $t$ \texttt{[plutôt }$T$\texttt{]}, par le bifoncteur $U'^{\circ}\times T\longrightarrow \mathrm{Ens}$ explicité dans la formule $(**)$ p.~72. Nous nous intéressons surtout à dégager des conditions qui assurent les deux implications suivantes

\begin{itemize}
\item[a)] $i\in W\Longrightarrow i'\in W$ (et même $\Longleftrightarrow$).
\item[b)] $p\in W\Longrightarrow q\in W$ (et même $\Longleftrightarrow$).
\end{itemize}

\msp{page 75}
\textbf{Proposition 4.} \textit{Soit $t\in\Ob T$. Pour que $t\in\Ob\bar{U}$, il faut et il suffit que $q(t)\in\Ob\overline{U'}$. En d'autres termes, le `bord' $\dot{T}=T\cap\bar{U}$ de $T$ dans $X$ est `le même' que son bord dans $X'$.}

\dgrblock{dgrb514}{The one-sided formula has an explicit deformation retraction}
\begin{dcomment}{exposition}{p74}{The one-sided formula has an explicit deformation retraction}
The comparison of the two factorization formulas admits a direct argument which also verifies Grothendieck's cofinality claim with its variance intact. Write $E$ for the category of elements occurring in the general formula when $T'=T$. Its objects are
\[
 (u,\alpha:u'\to p(u),v,\beta:v\to t,h:u\to v).
\]
An arrow to $(u_1,\alpha_1,v_1,\beta_1,h_1)$ consists of
\[
 a:u\to u_1,\quad b:v_1\to v,\qquad
 p(a)\alpha=\alpha_1,\quad \beta b=\beta_1,\quad
 h=b h_1a.
\]
Let $D=u'\backslash U/t$ be the category in the boxed one-sided formula. Define
\[
 R:E\longrightarrow D,\qquad
 R(u,\alpha,v,\beta,h)=(u,\alpha,\beta h).
\]
On arrows $R$ takes $(a,b)$ to $a$; the equality
$\beta_1h_1a=\beta b h_1a=\beta h$ proves that this is an arrow of $D$. The inclusion $J:D\to E$ takes $(u,\alpha,\gamma)$ to $(u,\alpha,t,1_t,\gamma)$, and $RJ=1_D$.

For each object of $E$, the pair $(1_u,\beta)$ gives an arrow
\[
 JR(u,\alpha,v,\beta,h)
       \longrightarrow(u,\alpha,v,\beta,h).
\]
The equations above verify naturality, so $JR\Rightarrow1_E$. Consequently $J$ and $R$ induce inverse homotopy equivalences on nerves, and in particular inverse bijections on $\pi_0$. This proves more than the displayed comparison requires. The construction is natural in $u'$ and $t$, because precomposition of $\alpha$ and postcomposition of $\beta$ commute with $\beta h$.

The reference to the initial object $(t,1_t)$ of $(T/t)^{\mathrm{op}}$ is therefore sound. An initial object is cofinal for the colimit of a \emph{presheaf}; in covariant colimit notation one has passed to the opposite and it becomes terminal. The proof does not assert that an initial object computes arbitrary covariant colimits. This is the same variance discipline needed in the commentary on Chapter~I, manuscript p.~47 (``Cofinality retains an orientation''), but here the explicit retraction identifies the whole category of presentations rather than only invoking the abstract cofinality criterion.
\end{dcomment}


\textsc{Démonstration}. Dire que $t\in\Ob\bar{U}$ signifie aussi
\[
U/t\neq\emptyset
\]
et implique, par continuité, que $q(t)\in\overline{q(U)}\subseteq\overline{U'}$. Inversement, si $q(t)\in\overline{U'}$, i.e. s'il existe $u'\in U'$ tel que
\[
\Hom(u',q(t))(=\pi_0(u'\backslash U/t))\neq\emptyset,
\]
il en résulte que $u'\backslash U/t\neq\emptyset$, a fortiori $U/t$ (qui reçoit $u'\backslash U/t)\neq\emptyset$, q.e.d.

\textbf{Proposition 5.} \textit{Soit}
\[
\begin{tikzcd}
Y \arrow[r, "i"] \arrow[d, "p"'] & X \arrow[d, "q"] \\
Y' \arrow[r, "i'"] & X'
\end{tikzcd}
\]
% foot: Version 7 décembre 2002
% ===================== page 12 =====================
\printedpage{54}
% head: A. Grothendieck : Dérivateurs, Ch. XII, Édition des Universités de Montpellier II et Paris VII
\textit{un carré cocartésien dans \textup{Cat}. Si $i$ est une immersion ouverte (resp. fermée), alors $i'$ aussi, et le carré est aussi cartésien. Idem pour $p$ et $q$. Si $i$ et $p$ sont des immersions ouvertes (\footnote{\textbf{NB} Nous utilisons \textit{`immersion ouverte'}, \textit{`immersion fermée'} au \textit{sens strict}\,: c'est un \textit{monomorphisme} pleinement fidèle dont l'image est une sous-catégorie \textit{pleine,} soit ouverte, soit fermée.}) (resp. \dots), alors la situation est isomorphe à celle d'un objet $X'$ de \textup{Cat}, de deux ouverts (resp. fermés) $X$, $Y'$ de $X$ ayant pour réunion $X'$, pour intersection} % [sic] de $X$ -> de $X'$, X and Y' are opens of X' (their union is X')
\textit{$Y$. Inversement, dans une telle situationm le carré correspondant est cocartésien, i.e. $X$} % [sic] situationm -> situation, (comma), typing error
\texttt{[plutôt }$X'$\texttt{]} \textit{s'identifie à la somme amalgamée}
\[
X'\simeq X\amalg_{Y}Y'.
\]

Donc aussi on a la proposition 15, page 158, et ses corollaires.
% foot: Version 7 décembre 2002

% ===================== page 13 =====================
\printedpage{55}
% head: A. Grothendieck : Dérivateurs, Ch. XII, Édition des Universités de Montpellier II et Paris VII

\section*{2\quad Sommes amalgamées et $W$-équivalences.}

\subsection*{2.1\quad Axiome du carré cocartésien et axiome du cube (ou `Lemme des cinq').}

Nous supposons à présent donné un localiseur fondamental
\[
W\subseteq\mathrm{Fl}(\mathrm{Cat}),
\]
donc $W$ satisfait à W(1, 2 bis, 3). Nous ne supposons pas pour le moment pour les autres axiomes standard W(4) (axiome de fibration), W(5) (axiome de connexité), W(6) (axiome des limites inductives). Mais nous allons introduire à présent deux axiomes de stabilité supplémentaire, de nature différente, concernant les sommes amalgamées.

\begin{description}[leftmargin=5em, style=nextline]
\item[W(7 bis)] (Axiome du carré cocartésien.) Soient $X$ un objet de Cat, $U$, $U'$ deux ouverts de $X$ tels que $U\cup U'=X$, considérons le diagramme cocartésien correspondant (cf. prop. 5 en marge de la page \texttt{[ci-dessus]}) :

\msp{page 76}
\[
\begin{tikzcd}[column sep=small]
& U\cap U' \arrow[dl, hook', "i'_0"'] \arrow[dr, hook, "i_0"] & \\
U \arrow[dr, hook, "i"'] & & U' \arrow[dl, hook', "i'"] \\
& X &
\end{tikzcd}
\tag*{$(*)$}
\]

\dgrblock{dgrb515}{Orbits of crossing arrows and the information discarded}
\begin{dcomment}{examples}{p75}{Orbits of crossing arrows and the information discarded}
The pushout formula can be seen in a small calculation before it is used homotopically. Let $U$ be the discrete category on $u_0,u_1$, let $T=e$ with object $t$, and let $H(u_i,t)$ be a singleton. Thus $X$ is the span $u_0\to t\leftarrow u_1$. Collapse $U$ to $U'=e$ by $p$. The category $u'\backslash U/t$ has exactly two objects and no nonidentity arrows, so the pushout $X'=X\amalg_Ue$ has two objects and \emph{two} distinct parallel arrows $u'\rightrightarrows t$. Identifying the two objects of $U$ does not identify arrows which have no relation between them. Its nerve realizes to a circle, whereas the nerve of $X$ is contractible. Here $p$ is not a weak homotopy equivalence, so this is an illustration of the formula, not a counterexample to any asserted stability under a weak equivalence.

A group action makes the quotient and its lost homotopy particularly explicit. Put $U=BG$, let $T=e$, and take a right $G$-set $S$ as the presheaf $H(-,t)$. On pushing $BG$ to $e$, the mixed hom set becomes
\[
 \Hom_{X'}(u',t)=S/G.
\]
Indeed the category of elements has objects $s\in S$ and an arrow labelled $g$ from $s$ to $s'$ when $s=s'g$. All arrows are invertible, and two objects lie in the same connected component exactly when they lie in the same orbit. If $S=K\backslash G$ is a transitive right action, the automorphism group of the object $K$ is $K$. Inclusion of that object is fully faithful and essentially surjective, hence the category of elements is equivalent to $BK$.

Thus a single crossing arrow in the quotient may stand for a factorization category with homotopy type $BK$. Choosing $K$ nontrivial shows that a singleton ordinary colimit says nothing about the contractibility of its category of presentations. If $S=G$ with the free right action, the stabilizer is trivial and the category of elements is contractible. If $S$ is the singleton with trivial action, it is $BG$. Both have the same orbit set. The distinction is exactly the one that motivates the total asphericity conditions later in this section.
\end{dcomment}


Alors
\[
i_0\in W\qquad\Longleftrightarrow\qquad i\in W
\]
(et symétriquement
\[
i'_0\in W\qquad\Longleftrightarrow\qquad i'\in W).
\]
\end{description}

\textbf{NB} Par dualité, cet axiome est équivalent à l'axiome similaire, avec des \textit{fermés} $T$, $T'$ de $X$ de réunion $X$.

\textbf{Corollaire.} \textit{Supposons \textup{W(7 bis)}. Considérons les conditions suivantes sur le carré précédent.}

\begin{itemize}
\item[\textit{a)}] \textit{$U\cap U'$ $W$-asphérique.}
\item[\textit{b)}] \textit{$U$, $U'$ $W$-asphériques.}
\item[\textit{c)}] \textit{$X$ $W$-asphérique.}
\end{itemize}

\textit{Alors (a et b) implique c, et (b et c) implique a.}

\textbf{NB} Moralement, il faudrait aussi que (c et a) implique b, i.e. si $U\cap U'$ et $X$ sont $W$-asphériques, il en est de même de $U$ et de $U'$. Mais il ne semble pas que l'axiome du carré cocartésien permette de le prouver. Par contre, l'axiome qui suit impliquera le corollaire précédent sous sa forme complète.
% foot: Version 7 décembre 2002

% ===================== page 14 =====================
\printedpage{56}
% head: A. Grothendieck : Dérivateurs, Ch. XII, Édition des Universités de Montpellier II et Paris VII

\begin{description}[leftmargin=5em, style=sameline]
\item[W(7)] (`Axiome du cube', ou `Lemme des cinq'.) Soit un diagramme commutatif cubique
\[
\begin{tikzcd}[row sep=small, column sep=small]
U_1\cap U'_1 \arrow[rr] \arrow[dr] \arrow[dd, "f_{U_1\cap U'_1}"'] & & U_1 \arrow[dr] \arrow[dd, "f_{U_1}"{swap,pos=.25}] & \\
& U'_1 \arrow[rr, crossing over] & & X_1 \arrow[dd, "f"] \\
U_2\cap U'_2 \arrow[rr] \arrow[dr] & & U_2 \arrow[dr] & \\
& U'_2 \arrow[rr] \arrow[from=uu, crossing over, "f_{U'_1}"{pos=.25}] & & X_2
\end{tikzcd}
\]
représentant un morphisme (flèches verticales) d'un

\msp{page 77}
carré cocartésien\quad $\vcenter{\hbox{\begin{tikzcd}
U_1\cap U'_1 \arrow[r] \arrow[d] & U_1 \arrow[d] \\
U'_1 \arrow[r] & X_1
\end{tikzcd}}}$\quad du type envisagé dans W(7 bis) ($X_1$ réunion de deux ouverts $U_1$, $U'_1$) dans un autre (à indices 2). Considérons les conditions suivantes.

\dgrblock{dgrb516}{Ordinary covers yield homotopy pushouts, with one direction of excision}
\begin{dcomment}{exposition}{p76}{Ordinary covers yield homotopy pushouts, with one direction of excision}
Grothendieck's new axioms try to turn a cover by two categorical opens into an excision principle. The underlying geometric fact is especially clean. Suppose $X=U\cup V$ with $U,V$ full sieves and $A=U\cap V$. A simplex of $NX$ is a composable chain $x_0\to\cdots\to x_n$. Its last vertex lies in $U$ or in $V$; since these are sieves, all earlier vertices then lie in the same one. Hence
\[
 NX=NU\cup_{NA}NV
\]
is an actual pushout of simplicial sets. The two maps out of $NA$ are monomorphisms. On geometric realization they are inclusions of CW subcomplexes, so this pushout also computes the homotopy pushout. For cosieves the same proof uses the first vertex of the chain.

This proves the direct implication of W(7 bis) for $W_\infty$: if $A\to V$ is a weak equivalence then $U\to X$ is a weak equivalence. It also proves the direct cube assertion $(a\text{ and }b)\Rightarrow c$: an objectwise weak equivalence of the two spans induces a weak equivalence of their homotopy pushouts. This invariance is the property of a derived colimit. It may be checked using the double mapping-cylinder model, or by mapping into a fibrant space, where the universal mapping space is the homotopy pullback of the three mapping spaces of the span.

The reverse implication requested in W(7 bis), and the two additional implications requested in W(7), are strictly stronger assertions. The formal corollaries drawn here from the proposed axioms are valid as conditional statements, as is Proposition~6's implication W(7)$\Rightarrow$W(7 bis). However, Theorem~2 on manuscript page~78 is false for the full printed axiom. A homotopy pushout can destroy the fundamental group of its intersection. The next block gives a finite categorical counterexample, with the ordinary and homotopy pushouts agreeing throughout. This is not a failure of the nerve to preserve the relevant pushout, and is not repaired by using a derived pushout instead.
\end{dcomment}


\begin{itemize}
\item[a)] $f_{U_1\cap U'_1}\in W$.
\item[b)] $f_{U_1}$ et $f_{U'_1}$ sont \texttt{[dans]} $W$.
\item[c)] $f\in W$.
\end{itemize}

Alors deux parmi ces conditions impliquent la troisième. En d'autres termes, quand l'une de ces conditions est satisfaite, les deux autres sont équivalentes entre elles.
\end{description}

\textbf{NB} Cet axiome est équivalent à l'axiome dual, en termes d'objets $X_1$, $X_2$ représentés chacun comme réunion de deux fermés.

\textbf{Corollaire.} \textit{Supposons \textup{W(7)} satisfait. Alors le corollaire à \textup{W(7 bis)} (axiome qu'on n'a pas à utiliser ici), sous sa forme complétée (cf. le NB qui le suit) est valable : deux parmi les conditions a), b), c) impliquent la troisième.}

Il suffit de regarder l'unique morphisme du diagramme $(*)$, p.~76, dans le diagramme trivial\quad $\vcenter{\hbox{\begin{tikzcd}[column sep=small]
& e \arrow[dl] \arrow[dr] & \\
e \arrow[dr] & & e, \arrow[dl] \\
& e &
\end{tikzcd}}}$\quad correspondant au cas où $X$ est la catégorie finale $e$, et $U$, $U'$ égaux à $X$. Alors le corollaire n'est autre que W(7) appliqué à ce cas particulier.

Cet argument fait sentir que W(7) est sans doute plus fort que W(7 bis). Incité par ce sentiment, je m'aperçois avec plaisir qu'il est fondé :
% foot: Version 7 décembre 2002

% ===================== page 15 =====================
\printedpage{57}
% head: A. Grothendieck : Dérivateurs, Ch. XII, Édition des Universités de Montpellier II et Paris VII

\msp{page 78}
\textbf{Proposition 6.} \textit{L'axiome \textup{W(7)} (lemme des cinq) implique \textup{W(7 bis)} (axiome du carré cocartésien).}

On regarde le diagramme cubique suivant
\[
\begin{tikzcd}[row sep=small, column sep=small]
U\cap U' \arrow[rr, "i_0"] \arrow[dr, "i'_0"] \arrow[dd, "f_{U\cap U'}=i_0"'] & & U' \arrow[dr, "i'"] \arrow[dd, "f_{U'}=\mathrm{id}"{swap,pos=.25}] & \\
& U \arrow[rr, crossing over, "i"] & & X \arrow[dd, "f=\mathrm{id}_X"] \\
U' \arrow[rr, "\mathrm{id}"] \arrow[dr, "i'"] & & U' \arrow[dr, "i'"] & \\
& X \arrow[rr, "\mathrm{id}"] \arrow[from=uu, crossing over, "f_U=i"{pos=.25}] & & X,
\end{tikzcd}
\]
du carré cocartésien donné dans le carré cocartésien déduit de $X$ et de ses deux ouverts $X$, $U'$, d'intersection $U'$. Ici les deux flèches verticales $f$ et $f_{U'}$ sont des isomorphismes, donc \texttt{[dans]} $W$. Avec les notations de W(7), on a donc c) ($f\in W$), donc a) et b) sont équivalentes, i.e. $i_0\in W$ et $i\in W$ sont équivalentes.

\hfill q.e.d.

\textbf{Théorème 2.} \textit{Le localiseur fondamental $W_\infty$ satisfait à l'axiome \textup{W(7)}.}

C'est sûrement une conséquence du critère cohomologique d'\textsc{Artin-Mazur}, compte tenu de la suite exacte de Mayer-Vietoris pour l'un et l'autre des deux carrés cocartésiens, et les variantes non commutatives de Mayer-Vietoris. Je n'ai pas écrit la vérification en forme, et devrais sans doute réparer cette omission.

\dgrblock{dgrb517}{A finite counterexample to the asserted reverse excision}
\begin{dcomment}{development}{p78:Je n'ai pas écrit la vérification en forme}{A finite counterexample to the asserted reverse excision}
Grothendieck hopes that a five-lemma argument will recover the homotopy type of an intersection from the two pieces and their union. The obstruction can be isolated using the familiar acyclic noncontractible complex of \citet[Example~2.38]{Hatcher2002}. We give the categorical consequence in full; the input space and its properties are classical. We checked that example, the finite simplicial replacement theorem, the subsequent distinction between the direct and reverse axioms in Chapter~XVI, and the model-independent homotopy-pushout argument. We did not find this exact finite-poset formulation attached to Theorem~2 in the sources checked; no priority is claimed.

\textbf{Proposition (finite failure of W(7 bis) and W(7)).}\ours
\emph{There is a finite poset $X$, covered by two full sieves $U,V$, such that $NU,NV,NX$ are contractible but $N(U\cap V)$ is not weakly contractible. Consequently $W_\infty$ does not satisfy either printed axiom.}

\emph{Proof.} Start with a connected CW complex $K_0$ with one vertex, two oriented edges $a,b$, and two two-cells attached along
\[
 a^5b^{-3},\qquad b^3(ab)^{-2}.
\]
Its cellular differential from dimension two to dimension one is
\[
 d_2=\begin{pmatrix}5&-2\\-3&1\end{pmatrix}.
\]
Its determinant is $-1$, so $d_2$ is an isomorphism. The differential $d_1$ is zero and there are no higher cells. It follows that $\widetilde H_*(K_0;\mathbb Z)=0$.

Here is an explicit verification that its fundamental group is nontrivial. With permutations composed from right to left, send
\[
 a\longmapsto(1\ 4\ 3\ 2\ 5),\qquad
 b\longmapsto(1\ 2\ 3).
\]
The product $ab$ is $(1\ 5)(3\ 4)$, so the images satisfy $a^5=b^3=(ab)^2=1$ and define a homomorphism from the presented group to $A_5$. The image of $a$ is nonidentity, proving nontriviality. We need neither a computation of the full group nor asphericity of the presentation complex.

By the finite case of \citet[Theorem~2C.5]{Hatcher2002}, choose a finite simplicial complex $K$ homotopy equivalent to $K_0$. Let $A$ be the poset of its nonempty simplices ordered by inclusion. Its nerve is the barycentric subdivision of $K$, so $|NA|$ is connected and acyclic, and has nontrivial fundamental group. Adjoin two incomparable elements $u,v$ above every element of $A$, and call the resulting finite poset $X$. Put
\[
 U=A\cup\{u\},\qquad V=A\cup\{v\}.
\]
Both are downsets of $X$, hence full sieves, and $U\cap V=A$. They have terminal objects $u$ and $v$, respectively. Their nerves are the two cones on $NA$. A chain in $X$ cannot contain both $u$ and $v$, and therefore $|NX|$ is the unreduced suspension of $|NA|$.

The suspension of a connected space is simply connected, by van Kampen applied to its two thickened cones. Mayer--Vietoris gives
\[
 \widetilde H_n(|NX|;\mathbb Z)
       \cong\widetilde H_{n-1}(|NA|;\mathbb Z)=0.
\]
Thus $|NX|$ is a simply connected acyclic CW complex. It is contractible by the homological Whitehead theorem \citep[Corollary~4.33]{Hatcher2002}. The inclusion $U\to X$ is consequently a weak homotopy equivalence. The inclusion $A\to V$ is not, since its source has nontrivial fundamental group and its target is contractible. This contradicts the reverse implication of W(7 bis). The map from the cover square of $X$ to the constant square on $e$ satisfies conditions $b,c$ of W(7) but not $a$. It therefore contradicts W(7) as well.\hfill$\square$

The proposed cohomological proof fails at a specific point. Local systems on the intersection need not extend to either piece or to the union. In the example, the nontrivial representation of $\pi_1(NA)$ in $A_5$ cannot extend across either cone. Mayer--Vietoris for coefficients restricted from the ambient contractible space only sees constant coefficients on the intersection, for which its acyclicity suppresses the obstruction. The nonabelian comparison required by a weak-equivalence criterion cannot be recovered from that restricted coefficient class. This is the same general difficulty as in the commentary on Chapter~I, manuscript p.~81 (``The coefficient lost in the attempted proof''), now realized in a finite categorical cover.
\end{dcomment}


\msp{page 79}
\textbf{Remarque.} La version duale (au sens de \textsc{Quillen}) de l'axiome W(7 bis) du carré cocartésien, i.e. l'axiome correspondant pour les carrés cartésiens, n'est pas valable pour $W_\infty$. L'énoncé pertinent serait le suivant : soit
\[
\begin{tikzcd}
X & X' \arrow[l, "q"'] \arrow[d, "f'"] \\
S & S' \arrow[l, "p"']
\arrow[from=1-1, to=2-1, "f"']
\end{tikzcd}
\]
un carré cartésien, avec $f$ une \textit{$W$-fibration.} Alors $f\in W\Longleftrightarrow f'\in W$. Si $W$ satisfait à W(4c) (ce qui est le cas pour $W_\infty$), alors $f\in W$ implique bien que ses fibres sont $W$-asphériques, donc que c'est une $W$-équivalence universelle, i.e. $f'\in W\ \forall\, p:S'\longrightarrow S$. Mais pour l'implication inverse $f'\in W\Longrightarrow f\in W$, il faut supposer visiblement que
\[
\pi_0(p):\pi_0(S')\longrightarrow\pi_0(S)
\]
% foot: Version 7 décembre 2002

% ---- from Der-ChXIIb-012_p016-020.tex
\printedpage{58}
% head: A. Grothendieck : Dérivateurs, Ch. XII, Édition des Universités de Montpellier II et Paris VII
% foot: Version 7 décembre 2002

est \textit{surjectif.} (P. ex. si $S'=\emptyset$, on a $X'=\emptyset$, donc $f'\in W$, sans que ça implique $f\in W$!) Le fait que dans le cas cocartésien de W(7 bis), $i_0\in W\Longrightarrow i\in W$ doit être considéré comme ‘allant de soi’ - autrement le carré ne mérite pas, du point de vue $W$-homotopique, le nom de ‘cocartésien’. Par contre, l’implication inverse $i\in W\Longrightarrow i_0\in W$ m’apparaît comme un trait très particulier du dérivatuer HOT$_W$ (quand il le possède), et notamment % [sic] dérivatuer -> dérivateur, misprint
de HOT $=$ HOT$_{W_\infty}$. Le début
\msp{page 80}
de la remarque signifie simplement que le dérivateur dual HOT$^\circ$ ne possède pas cette même propriété. Je dois renvoyer à ce sujet aux deux axiomes pour les dérivateurs correspondants à W(7 bis), W(7), ce que je ferai sans doute sous peu, de façon à donner un sens plus précis aux présents commentaires.

\dgrblock{dgrb518}{What survives of the cube principle}
\begin{dcomment}{status}{p79}{What survives of the cube principle}
The distinction between direct and reverse excision is already visible in Grothendieck's later reorganization. Chapter~XVI separates the cube conditions into Loc(6a), Loc(6b), and Loc(6c); its Commentary~3.4 distinguishes the direct theory of homotopy cofibres from the extra reversibility property \citep[\S3, manuscript pp.~16--27]{GrothendieckDerivateursXVI}. For the pages at hand the reliable assertion for ordinary homotopy types is the direct one: weak equivalences of all three entries of a span give a weak equivalence of its homotopy pushout. A reverse five lemma cannot be added without hypotheses that exclude the preceding example. Theorem~2 has no repair with its stated hypotheses; its direct part is a theorem.

The pullback remark makes a different, correct observation. Base change need not meet every connected component of the base. Take $S=\{0,1\}$ discrete, and the discrete fibration $X\to S$ with one point over $0$ and two points over $1$. Pulling back along $\{0\}\to S$ gives an isomorphism, while $X\to S$ is not even bijective on connected components. This is a nonempty version of the empty-base example printed in the source.

Under the intended fibration hypothesis, contractibility of homotopy fibres is the correct test. If base change meets every component and the pullback is a weak equivalence, its homotopy fibres are contractible. Fibre transport along paths then shows that every homotopy fibre of the original fibration is contractible, so the original map is a weak equivalence. This argument explains the appearance of surjectivity on $\pi_0$; it must be applied where the source's W-fibration axioms really identify categorical fibres with homotopy fibres. Merely calling an arbitrary functor a fibration does not supply that identification. The corresponding discussion of categorical fibres was developed in the commentary on Chapter~VI, manuscript p.~45 (``Why the new fibration axiom selects full homotopy theory'').

The final contrast with the dual derivator should therefore be read in two parts. Its warning about detection after pullback is justified. Its attribution of unrestricted reverse excision to $\mathrm{HOT}$ depends on the false Theorem~2. Both qualifications are needed before applying the proposed axioms to the less special pushouts that follow.
\end{dcomment}


\dgrblock{dgrb519}{What the symmetric cube rule can detect}
\begin{dcomment}{development}{p79}{What the symmetric cube rule can detect}
The failed implication in Theorem~2 has a structural source. A homotopy pushout can erase information in its intersection. It is useful to determine whether passing to a bounded homotopy theory removes that obstruction, and to compare the answer with the homological argument suggested in the manuscript. The following finite family gives a uniform negative answer to the first question. We did not locate this precise family of tests in the source, its Appendix, or the later papers checked; the observation is elementary and is offered without a priority claim.

\textbf{Finite truncation obstruction.}\ours\quad
For an integer $n\geq0$, let $W_{\leq n}$ consist of functors whose nerves induce a bijection on components and isomorphisms on all homotopy groups in degrees $1,\ldots,n$, at every basepoint. There is a finite poset covered by two sieves for which the reverse implication of W(7 bis) fails for $W_{\leq n}$.

\emph{Proof.} Let $A_n$ be the poset of nonempty proper subsets of a set with $n+2$ elements, ordered by inclusion. A nondegenerate simplex of $NA_n$ is a strictly increasing chain of such subsets. Thus $NA_n$ is the barycentric subdivision of the boundary of the $(n+1)$-simplex, and $|NA_n|$ is homeomorphic to $S^n$. Adjoin two distinct elements $+$ and $-$, each greater than every element of $A_n$, with no order relation between them. Call the resulting poset $D_n$, and put $U_+=A_n\cup\{+\}$ and $U_-=A_n\cup\{-\}$. These are sieves, their intersection is $A_n$, and each has a greatest element. Their nerves are cones on $NA_n$. Since a chain of $D_n$ can contain at most one of $+$ and $-$, their union is exactly $ND_n$, whose realization is the suspension of $|NA_n|$, hence $S^{n+1}$.

It follows that $U_+\hookrightarrow D_n$ belongs to $W_{\leq n}$: both spaces are connected, and all their homotopy groups in degrees at most $n$ vanish. On the other hand, $A_n\hookrightarrow U_-$ does not belong to $W_{\leq n}$. For $n=0$ it identifies two components. For $n\geq1$ it kills the nonzero group $\pi_n(S^n)\simeq\mathbb Z$. This is precisely the forbidden reverse implication in the cover square.\hfill$\square$

The example for $n=1$ has eight objects in the union: six objects forming the barycentric subdivision of a triangle boundary and the two cone points. These tests do not replace the acyclic example at p.~78, since the latter is needed for untruncated weak equivalences. Together they show that the issue is not repaired by merely stopping the homotopy groups at some finite degree.

For comparison, consider the class of functors inducing isomorphisms on nerve homology with coefficients in a fixed ring $R$. The full symmetric cube rule \emph{does} hold for that class. Here is the complete chain-level reason. If $X=U\cup U'$ is a union of two sieves and $A=U\cap U'$, every chain of composable arrows in $X$ belongs to one of the two subcategories: choose a subcategory containing its last vertex and use the sieve condition. The same statement holds for nondegenerate chains. Consequently the normalized chain complexes have the natural short exact sequence
\[
0\longrightarrow C_*(NA;R)
 \xrightarrow{(\iota,-\iota')}
 C_*(NU;R)\oplus C_*(NU';R)
 \xrightarrow{\jmath+\jmath'} C_*(NX;R)
 \longrightarrow0.
\]
Injectivity, exactness in the middle, and surjectivity can each be checked on the basis of chains: those in the intersection occur twice, all others once. For a morphism between two such covers, take the mapping cones of the three vertical chain maps. The result is again a short exact sequence of complexes. Its long exact homology sequence implies that if any two of these cones are acyclic, then so is the third. The middle cone is the direct sum of the cones for the maps on $U$ and $U'$, so its acyclicity is equivalent to the acyclicity of both summands. This proves all three implications requested in W(7). For cosieve covers, the first vertex supplies the same proof.

The manuscript's proposed Mayer--Vietoris justification therefore has an exact domain of validity: it works for fixed-coefficient homology equivalences. The passage from those equivalences to weak homotopy equivalences is the step that cannot be supplied. The acyclic intersection in the preceding counterexample makes that distinction visible without any ambiguity about the existence or computation of the pushout.
\end{dcomment}


\subsection*{2.2 Application aux carrés cocartésiens déduits d’une immersion ouverte ou fermée.}

Les axiomes W(7), W(7 bis) (du cube, et du carré cocartésien) concernent des carrés cocartésiens de nature très particulière. Nous allons voir à présent comment en déduire des énoncés du même type pour des carrés cocartésiens beaucoup plus généraux, y compris pour des carrés qui n’ont rien de cocartésien (dans Cat) - et que nous appelerons % [sic] appelerons -> appellerons, misspelling
‘homotopiquement cocartésiens’ (relativement à $W$), ou simplement $W$-cocartésiens. (Il y aura également des carrés $W$-cartésiens un jour \dots) Dans le cas où $W=W_\infty$, je les appelerai aussi Hot-cocartésiens. % [sic] appelerai -> appellerai, misspelling

Je commencerai par le cas particulier d’un carré cocartésien associé à une immersion ouverte (diagramme ($*$), p. 74) :
\[
(*)\qquad\qquad
\begin{tikzcd}
U \arrow[r, hook, "i"] \arrow[d, "p"'] & X \arrow[d, "q"'] & T \arrow[l, hook', dotted, "j"'] \arrow[d, dotted] \\
U' \arrow[r, hook, "i'"] & X' & T'\simeq T. \arrow[l, hook', dotted, "j'"']
\end{tikzcd}
\]
\msp{page 81}

Le grand traquenard ici qu’il faut arriver à exorciser, c’est qu’il peut arriver que $i\in W$, sans qu’on ait $i'\in W$, et aussi $p\in W$ sans qu’on ait $q\in W$. Je m’en étais convaincu en 1983 en écrivant Pursuing Stacks (durant les semaines de ‘scratchwork’), et n’ai pas envie de rechercher un contre-exemple ici, s’il ne se présente de lui-même (ce qu’il fera sans doute).

Je vais donner un cas important où ces ennuis disparaissent :

\dgrblock{dgrb520}{Separating the boundary from the ambient category}
\begin{dcomment}{exposition}{p81:Je vais donner un cas important où ces ennuis disparaissent}{Separating the boundary from the ambient category}
Grothendieck wants to isolate the part of a pushout which can change its homotopy type. The complement $T$ is carried isomorphically to $T'$, whereas the arrows entering it from $U$ are changed by a Kan extension. The useful intermediate category is therefore the closure $\bar U$, which consists of $U$ and the objects of $T$ reached by at least one arrow from $U$. Write $A=\bar U\cap T$. We then have a square
\[
\begin{tikzcd}
A\arrow[r]\arrow[d]&T\arrow[d]\\
\bar U\arrow[r]&X
\end{tikzcd}
\]
of cosieve inclusions. It is both Cartesian and cocartesian. Every arrow of $X$ lies in one of these two full subcategories: if its source lies in $\bar U$, its target does too; otherwise its source lies in $T$, and its target lies in $T$. The same assertion for composable strings shows that the nerve is the union of the two corresponding simplicial subsets. Its realization is consequently a homotopy pushout, since inclusions of simplicial subsets realize as cofibrations.

The formula for the pushed-out gluing presheaf shows that an object $t\in T$ receives an arrow from $U'$ after pushout exactly when it received an arrow from $U$ before pushout. Indeed, a nonempty Kan-extension value is represented by an original arrow and a morphism in $U'$; conversely, an original arrow supplies such a value at the image of its source. Thus $A\to A'$ is an isomorphism. All the movement has been confined to $\bar U\to\overline{U'}$.

This explains the architecture of Proposition~6. If $U\to\bar U$ and $U'\to\overline{U'}$ are weak equivalences, two-out-of-three replaces the question about $p$ by the question about $\bar p$. The resulting map of closed-cover squares is the identity on $A$ and $T$. The ordinary homotopy gluing lemma therefore proves, for $W_\infty$, the useful implication
\[
p\in W_\infty\quad\Longrightarrow\quad q\in W_\infty.
\]
This argument requires only the usual forward gluing statement.

The biconditionals in the printed proposition have stronger hypotheses. They explicitly invoke W(7~bis) and W(7), with their reverse implications as formulated on pp.~76--77. The counterexample discussed at p.~78 prevents us from treating those axioms as properties of $W_\infty$. Proposition~6 remains a valid formal deduction under its stated axioms; its proof is not a proof that ordinary weak homotopy equivalences can be detected backwards from every such pushout. This distinction will matter when the one-point test is proposed on p.~85.
\end{dcomment}


\textbf{Proposition 6 :} \textit{Considérons les adhérences $\bar U$ et $\overline{U'}$ de $U$ et $U'$ dans $X$, $X'$ respectivement, et supposons que les inclusions}
\[
U\xrightarrow{i_0}\bar U,\qquad U'\xrightarrow{i'_0}\overline{U'}
\]
\printedpage{59}
% head: A. Grothendieck : Dérivateurs, Ch. XII, Édition des Universités de Montpellier II et Paris VII
% foot: Version 7 décembre 2002
\textit{soient dans $W$} (\footnote[64]{cf. th. 3 p. 107 pour des conditions générales assurant que cette hypothèse est satisfaite.}). \textit{Alors}
\[
\left\{
\begin{array}{llll}
\circled{\textit{a}}\ i\in W & \iff & i'\in W & \textit{si on a }\mathrm{W}(7\text{ bis})\\
\circled{\textit{b}}\ p\in W & \iff & q\in W & \textit{si on a }\mathrm{W}(7).
\end{array}
\right.
\]

\textsc{Démonstration}. Considérons d’abord le cas \circled{a} de $i$, $i'$. Soit
\[
k:\bar U\hookrightarrow X
\]
l’inclusion. Le triangle
\begin{tikzcd}
U \arrow[r, "i_0"] \arrow[rr, bend right=8, "i"'] & \bar U \arrow[r, "k"] & X
\end{tikzcd}
, où $i_0\in W$, montre que
\[
i\in W\iff k\in W.
\]

D’autre part, $X$ est la réunion des fermés $\bar U$, $T$, donc on a un carré cocartésien dans Cat
\[
(*)\qquad\qquad
\begin{tikzcd}
\bar U\cap T=\dot T \arrow[r, "k_0"] \arrow[d, "j_0"'] & T \arrow[d, "j"] \\
\bar U \arrow[r, "k"] & X.
\end{tikzcd}
\]

L’axiome W(7 bis) (sous sa forme duale (i.e. appliqué au carré formé par les catégories opposées \dots)) implique
\msp{page 82}
que
\[
k\in W\iff k_0\in W,
\]
d’où
\[
i\in W\iff (k_0:\dot T\longrightarrow T)\in W.
\]
De même on trouve, en appliqunt ceci à $(X',U')$ au lieu de $(X,U)$ % [sic] appliqunt -> appliquant, misprint
\[
i'\in W\iff (k'_0:\dot T'\longrightarrow T')\in W.
\]
Or on a $T\xrightarrow{\sim}T'$ (isomorphisme), et par la proposition 4 cet isomorphisme transforme $\dot T$ en $\dot T'$. Donc
\[
k_0\in W\iff k'_0\in W\qquad\quad\text{(trivialement)},
\]
d’où a).

Prouvons b). Le carré commutatif induit par ($*$), p. 80,
\printedpage{60}
% head: A. Grothendieck : Dérivateurs, Ch. XII, Édition des Universités de Montpellier II et Paris VII
% foot: Version 7 décembre 2002
\[
\begin{tikzcd}
U \arrow[r, hook, "i_0"] \arrow[d, "p"'] & \bar U \arrow[d, "\bar p"] \\
U' \arrow[r, hook, "i'_0"] & \overline{U'},
\end{tikzcd}
\]
où $i_0,i'_0\in W$ par hypothèse, montre que
\[
p\in W\iff\bar p\in W,
\]
donc tout revient à prouver que
\[
\bar p\in W\overset{?}{\iff}q\in W.
\]
Mais $q$ induit un homomorphisme du carré ($*$) déduit de $(X,\bar U,T)$ (avec $X=\bar U\cup T$), dans le carré similaire pour $X'$
\[
\begin{tikzcd}
\overline{U'}\cap T'=\dot T' \arrow[r, "k'_0"] \arrow[d, "j'_0"'] & T \arrow[d, "j'"] \\ % [sic] T -> T', the square for X' should have the corner T' (although T'\simeq T)
\overline{U'} \arrow[r, "k'"] & X',
\end{tikzcd}
\]
les flèches verticales étant notées $q_{\dot T}$, $q_T$, $\bar p$, $q$.
\msp{page 83}
On sait que $q_{\dot T}$, $q_T$ sont des isomorphismes, donc dans $W$, donc le ‘lemme des cinq’ W(7) nous assure que $\bar p\in W\iff q\in W$, q.e.d.

\textbf{Commentaire.} Cette proposition nous amène à chercher des conditions sur un ouvert $U$ d’un $X$, qui assurent que l’inclusion $i_0:U\longrightarrow\bar U$ soit \texttt{[dans]} $W$, et \textit{que cette propriété soit conservée par tout cochangement de base $U\longrightarrow U'$.} De plus, nous recherchons des conditions qui soient de nature locale sur $\bar U$, plus exactement sur $\bar U\cap T=\dot U=\dot T$. Nous pouvons motiver ce souci (a posteriori) en disant que nous voudrions que la propriété $i_0\in W$ soit conservée, non seulement par cochangement de base quelconque $U\longrightarrow U'$, mais aussi, de plus, par passage de $X$ à toute sous-catégorie pleine $X_1$ contenant $U$, donnant donc naissance à une situation
\[
U\overset{i_Y}{\hookrightarrow}X_1,\qquad U\xrightarrow{(i_Y)_0}\bar U^{X_1}=\bar U\cap X_1
\]
\texttt{[}$Y=X_1$\texttt{?]}, similaire à la situation de départ. Plus généralement encore, nous voudrions qu’elle soit conservée en passant du triple $(X,U,T)$ à un triple “induit” $(X_1,U,T_1)$, par une flèche quelconque $g:T_1\longrightarrow T$ de but $T$ dans Cat, donnant naissance au couple de flèches
\[
U\xrightarrow{\mathrm{id}}U,\ g:T_1\longrightarrow T
\]
(cf. pages 60, 61). Le cas précédent est celui où
\printedpage{61}
% head: A. Grothendieck : Dérivateurs, Ch. XII, Édition des Universités de Montpellier II et Paris VII
% foot: Version 7 décembre 2002
\msp{page 84}
$T_1$ est une sous-catégorie pleine de $T$, $g$ la flèche d’inclusion, Notons à ce sujet la proposition soritale % [sic] inclusion, Notons -> inclusion. Notons, comma before a capital letter

\textbf{Proposition 7.} \textit{Conditions équivalentes sur l’ouvert $U$ d’un objet $X$ de} Cat (moyennant W(7 bis)) :

\begin{enumerate}[label={}, leftmargin=*]
\item[\textit{a)}] \textit{L’inclusion $i_0:U\longrightarrow\bar U$ est $W$-asphérique, i.e. pour tout $t\in\dot U=\bar U\backslash U$, la catégorie $U/t$ est $W$-asphérique.}
\item[\textit{b)}] \textit{L’inclusion $i_0:U\longrightarrow\bar U$ est dans $W$, et cette condition reste satisfaite, quand on remplace $(X,U)$ par $(X_1,U)$, où $(X_1,U)$ est déduit de $(X,U)$ au moyen d’une flèche $T_1\longrightarrow T(=X\backslash U)$ dans} Cat\textit{, comme rappelé ci-dessus.}
\item[\textit{b$'$)}] \textit{Comme b), dans le cas où $T_1$ est la catégorie ponctuelle $e$, et où de plus son image est dans $\dot U$.}
\end{enumerate}

\textsc{Démonstration}. Il est évident que si dans b), b$'$) on remplace la condition “$i_0\in W$” par la condition plus forte ‘$i_0$ est $W$-asphérique’, les quatre \texttt{[plutôt trois]} conditions a) à b$'$) sont équivalentes. Il reste à prouver que la condition que les $U/t$ soient $W$-asphériques ($t\in\bar U$ - mais c’est trivial pour $t\in U$, donc il suffit de l’exprimer pour $t\in\bar U\backslash U=\dot U$) est déjà impliquée par la plus faible b$'$ et on devine qu’on a ceci :

\msp{page 85}

\textbf{Corollaire.} (Moyennant W(7 bis).) \textit{Soit $X$ dans} Cat\textit{, $U$ un ouvert de $X$ tel que $X\backslash U=\overline{\{t\}}$ soit la catégorie finale. Alors l’inclusion $U\xrightarrow{i}X$ est dans $W$ si et seulement si elle est $W$-asphérique, i.e. si et seulement si $U/t$ est $W$-asphérique (où $t$ est l’unique élément de $T$).}

Nous démontrons ultérieurement ce corollaire (\footnote[65]{cf. cor. 4 de la proposition 15, page 149.}). La proposition 7 ne nous servira pas d’ici là - elle n’aura qu’un rôle purement esthétique à jouer, pour justifier notre propos délibéré de chercher des conditions sur $i:U\longrightarrow X$ (pour assurer que $U\longrightarrow\bar U$ est dans $W$ et le reste, quand on remplace $(X,U)$ par d’autres couples $(X_1,U_1)$ qui s’en déduisent par certaines opérations \dots) qui
\msp{page 86}
soient de \textit{nature locale} sur $\bar U$ (en fait, en les points de $\dot U$), et qui en tout état de cause entrainent la \textit{$W$-asphéricité} de $U\longrightarrow\bar U$.

\dgrblock{dgrb521}{The one-point test and the direction of Theorem A}
\begin{dcomment}{exposition}{p85}{The one-point test and the direction of Theorem A}
The purpose of Proposition~7 is to turn a condition on the whole inclusion into a condition at each new point. For a sieve $i:U\hookrightarrow X$, the comma category $U/u$ has final object $(u,\mathrm{id}_u)$ whenever $u\in U$. All possible failures of asphericity are therefore concentrated at the complementary objects. Replacing $T$ by the point selecting $t$ preserves precisely the presheaf $u\mapsto\Hom_X(u,t)$, and hence precisely the category $U/t$. This proves the equivalence of the three versions in which every occurrence of ``belongs to $W$'' is replaced by ``is $W$-aspherical.'' The remaining implication is exactly the deferred corollary, not a consequence of the definition.

For ordinary weak homotopy equivalences, the corollary's reverse implication is false. To see its geometry without altering the source, take a connected acyclic category $A$ whose nerve is not weakly contractible, as in the counterexample at p.~78. Form $U$ by adjoining a final object $u_+$ to $A$. Form $X$ by adjoining instead two distinct final cone points $u_+$ and $t$, with no morphisms between them. Then $U\hookrightarrow X$ is a sieve and its complement is the one-object category $\{t\}$. Every object of $A$ has its unique arrow to $t$, while $u_+$ has none, so
\[
U/t\simeq A.
\]
On the other hand, $BU$ is contractible and $BX\simeq\Sigma BA$. The suspension of a connected acyclic CW complex is simply connected and acyclic, hence weakly contractible by the Hurewicz and Whitehead theorems. Thus $U\to X$ is a weak homotopy equivalence although it is not aspherical.

The missing distinction is between the homotopy type of an attachment and the homotopy type of its attaching category. A cone attachment can erase a nontrivial perfect fundamental group without leaving any homology behind. Theorem~A supplies the implication from contractible comma categories to weak equivalence; it does not supply the converse. In modern language this is also the distinction between a weak equivalence of classifying spaces and a homotopy initial functor: the comma categories here are $U/t=(i\downarrow t)$. Compare the dual of \citet[Theorem~4.1.3.1 and Corollary~4.1.3.3]{Lurie2009HTT}, applied to $i^{\mathrm{op}}$. The source has retained W(7~bis) in the corollary, and the audit retains it too. Its subsequent use for $W_\infty$ needs an independent argument in each application.
\end{dcomment}


J’avais espéré que cette condition d’asphéricité (suggérée par la notion ‘d’immersion de Dwyer’) serait stable par cochangement de base quelconque $U\longrightarrow U'$, ce qui aurait assuré en même temps qu’après une telle opération, l’immersion ouverte $i':U'\longrightarrow X'$ obtenue satisfait encore à la condition que $i'_0:U'\longrightarrow\overline{U'}$ est dans $W$. Nous verrons que cette attente n’est pas fondée - la condition de $W$-asphéricité n’est \textit{pas} stable par cochangement de base, et elle ne suffit \textit{pas} même à assurer qu’après un cochangement de base, $i'_0:U'\longrightarrow\overline{U'}$ soit
\printedpage{62}
% head: A. Grothendieck : Dérivateurs, Ch. XII, Édition des Universités de Montpellier II et Paris VII
% foot: Version 7 décembre 2002
dans $W$. Il faudra une notion plus forte, qui est suggérée par le cas où $T=\{t\}$ est la catégorie finale et contenue dans $\bar U$ (cas auquel nous avons essayé de nous ramener dans la proposition 7). Dans ce cas, on l’a dit (corollaire précédent), la condition que $i:U\longrightarrow\bar U=X$ soit dans $W$, équivaut à la $W$-asphéricité de $i$, i.e. à la $W$-asphéricité de $U/t$, et de même pour $i':U'\longrightarrow\overline{U'}=X'$ : elle est dans $W$ si et seulement si $U'/t'$ est $W$-asphérique, où
\msp{page 87}
$t'$ est l’unique élément de $T'(\simeq T)=X'\backslash U'$. Donc il faut des conditions sur $U/t$ qui assurent que non seulement $U/t$, mais aussi toutes les catégories similaires $U'/t$ \texttt{[plutôt} $U'/t'$\texttt{]}, soient $W$-asphériques. Dans le cas général où $T$ est quelconque, c’est cette condition-là que nous allons imposer aux $U/t$ ($t\in\Ob T$), et il s’avérera que cette fois-ci, elles sont stables \texttt{[plutôt elle est stable]} par tout cochangement de base $U\xrightarrow{p}U'$, i.e. que les $U'/t'$ ($t'\in\Ob T'$) satisferont alors à la même condition. C’est cette condition que nous allons à présent reformuler.

\dgrblock{dgrb522}{A boundary presheaf and all its direct images}
\begin{dcomment}{exposition}{p87:C’est cette condition que nous allons à présent reformuler.}{A boundary presheaf and all its direct images}
Grothendieck now strengthens the boundary condition in the way dictated by the actual pushout formula. The extra quantifier is the essential feature. Given $H:U^{\mathrm{op}}\to\mathbf{Set}$, let $\int_U H=U/H$ denote its category of elements. Its objects are pairs $(u,h)$ with $h\in H(u)$, and a morphism $(u,h)\to(v,k)$ is an arrow $a:u\to v$ satisfying $H(a)(k)=h$. If $X$ is obtained by adjoining a complementary point $t$ with $\Hom_X(u,t)=H(u)$, then $U/t=\int_U H$ literally records all ways to enter that point.

For $p:U\to U'$, the new boundary is $p_!H$, where $p_!$ is the ordinary left Kan extension along $p^{\mathrm{op}}$. It is useful to calculate it through the category of elements, rather than to lose the variance in an informal colimit. With $\alpha:\int_UH\to U$ the projection,
\[
(p_!H)(u')\cong\pi_0\bigl(u'\downarrow p\alpha\bigr).
\]
An object on the right is $(u,h,\beta:u'\to p(u))$. Its connected component is the generator of the Kan extension represented by that triple. A map $a:u\to v$ with $H(a)(k)=h$ identifies the generator $(u,h,\beta)$ with $(v,k,p(a)\beta)$; these are exactly the relations in the pointwise colimit. This gives both the formula and its functoriality in $u'$.

The ordinary asphericity of $H$ only tests $p=\mathrm{id}_U$. Total asphericity requires the entire category of elements of every $p_!H$ to be aspherical. The operation takes connected components before the resulting presheaf is assembled over $U'$. It can therefore destroy homotopy information inside each comma category while creating nontrivial topology in the assembled category. The finite example at p.~100 will make this effect visible in degree two.

The definition has a deliberate stability property. If $H$ is totally $W$-aspherical and $p:U\to U'$ is any functor, then $p_!H$ is totally $W$-aspherical: for every $q:U'\to U''$ there is a canonical isomorphism $q_!p_!H\cong(qp)_!H$. Thus the strengthened condition is preserved by the very operation against which it was designed to protect the boundary. This is the formal reason the later definition of a $W$-Dwyer immersion uses these presheaves.

The editors' warning about terminology is necessary. The earlier test-category usage of ``totally aspherical'' concerns products or diagonal functors. Here the quantifier runs over all outgoing functors and their set-valued left Kan extensions. The commentary on Chapter~II, manuscript p.~36 (``The quantifier in total asphericity'') already treats this quantifier; the present pages apply it to a concrete gluing problem. Likewise the filtered examples should be read together with the commentary on Chapter~II, manuscript pp.~32--35, rather than as a new assertion that contractibility alone controls all direct images.
\end{dcomment}


\subsection*{2.3 Objets totalement $W$-asphériques de \textnormal{Cat}, et morphismes totalement $W$-asphériques.}

Reprenons la situation
\[
U\overset{i}{\hookrightarrow}X\overset{j}{\hookleftarrow}T,\qquad\qquad\text{avec }T=\{t\}\text{, catégorie finale, }t\in\bar U\text{, i.e. }X=\bar U,
\]
et celle d’un cochangement de base $U\longrightarrow U'$, donnant naissance à un carré cocartésien dans Cat
\[
(*)\qquad\qquad
\begin{tikzcd}
U \arrow[r, hook, "i"] \arrow[d, "p"'] & X \arrow[d, "q"'] & T=\{t\} \arrow[l, hook', dotted, "j"'] \arrow[d, dotted, "\mathrm{iso}"] \\
U' \arrow[r, hook, "i'"] & X' & T'=\{t'\}. \arrow[l, hook', dotted, "j'"']
\end{tikzcd}
\]
La donnée de recollement $H$ entre $U$ et $T$, $H:\underbrace{U^\circ\times T}_{\simeq U^\circ}\longrightarrow\mathrm{Ens}$, peut s’interpréter comme un préfaisceau sur $U$, $H\in\Ob U^\wedge:U^\circ\longrightarrow\mathrm{Ens}$,
\msp{page 88}
et on a
\[
U/t\simeq U/H,
\]
où $U/H$ est interprété en termes du plongement canonique
\[
U\hookrightarrow U^\wedge.
\]
De même, le recollement entre $U'$ et $T'$ sera donné par un préfaisceau
\[
H'\in\Ob U^{\prime\wedge}:U^{\prime\circ}\longrightarrow\mathrm{Ens},
\]

% ---- from Der-ChXIIb-012_p021-025.tex
\printedpage{63}
% head: A. Grothendieck : Dérivateurs, Ch. XII, Édition des Universités de Montpellier II et Paris VII
et on aura
\[ U'/t' \simeq U/H'. \]
La relation entre $H'$ et $H$ a été explicitée au §1, comme
\[ H' = p_!H, \qquad \text{où } p_! : U^{\wedge} \longrightarrow U'^{\wedge} \]
est l'adjoint à gauche de $p^* : U'^{\wedge} \longrightarrow U^{\wedge}$, donné par
\[ p^*F' = F' \circ p. \]
On a
\[ p_!(H)(u') \simeq \pi_0(u'\backslash U/H). \]

\textbf{Définition 1.} \cercle{a} Soit $U$ un objet de Cat, $H \in \Ob U^{\wedge}$ un préfaisceau sur $U$. On dit que $H$ est \textit{totalement $W$-asphérique} si pour toute flèche $p : U \longrightarrow U'$ de Cat, $p_!(H)$ dans $U'^{\wedge}$ est un objet $W$-asphérique, i.e. $U'/p_!H$ est $W$-asphérique. \cercle{b} On dit qu'un objet $V$ de \msp{page 89} Cat est \textit{totalement $W$-asphérique,} si le foncteur final $e_V$ sur $V$ ($v \longmapsto \{\emptyset\}$) est totalement $W$-asphérique, i.e. si pour toute flèche $f : V \longrightarrow V'$ dans Cat, le préfaisceau
\[ \pi_0(f) \defeq f_!(e_V) \defeq (v' \longmapsto \pi_0(v'\backslash V)) \]
sur $V'$ est $W$-asphérique, i.e. si la catégorie
\[ \upi_0(f) \ \text{ ou } \ \upi_0(V/V') \defeq V'/\pi_0(f) \]
est $W$-asphérique. \texttt{[Cette notion est différente de la notion de totale asphéricité de Pursuing Stacks.]}

Notons que la définition des foncteurs $p_!$ (ou $f_!$) comme adjoints, et leur `calcul' en termes de $\pi_0$ de certaines catégories, ne dépend pas du localiseur $W$. Donnons quelques précisions à ce sujet. On a un morphisme d'adjonction dans $V^{\wedge}$
\[ f^*f_!(F) \longleftarrow F, \]
d'où
\[ V/F \longrightarrow V/f^*f_!(F) \simeq V \times_{V'} (V'/f_!(F)), \tag*{$(*)$} \]
le dernier isomorphisme provenant du fait que l'on a commutativité, à isomorphisme canonique près, dans le diagramme de foncteurs
\[
\begin{tikzcd}[column sep=14em, row sep=large]
V'^{\wedge} \arrow[r, "f^*=(F'\longmapsto F'\circ f)"] \arrow[d, "F'\longmapsto V'/F'"'] & V^{\wedge} \arrow[d, "F\longmapsto F/V"] \\
\Cat/V' \arrow[r, "f^*_{\Cat}=(W'\longmapsto W'\times_{V'}V)"] & \Cat/V,
\end{tikzcd}
\]
% foot: Version 7 décembre 2002
\printedpage{64}
% head: A. Grothendieck : Dérivateurs, Ch. XII, Édition des Universités de Montpellier II et Paris VII
\msp{page 90}
en d'autres termes, l'opérateur $f^*$ sur les préfaisceaux s'interprète, en termes géométriques, comme l'opération $W' \longmapsto W' \times_{V'} V$ de changement de base par $V \longrightarrow V'$, sur les catégories au dessus de $V'$ qui sont fibrées à fibres discrètes (propriété qui se conserve par changement de base).

Ceci dit, la donnée d'une flèche $(*)$ équivaut, par la propriété universelle du changement de base, à celle d'une flèche (canonique en l'occurrence)
\[ f/F : V/F \longrightarrow V'/f_!(F) \]
rendant commutatif le carré
\[
\begin{tikzcd}[column sep=large, row sep=huge]
V/F \arrow[r, "\text{oubli de }F", "\substack{\text{ou foncteur}\\ \text{source}}"'] \arrow[d, "f/F"'] & V \arrow[d, "f"] \\
V'/(F' = f_!(F)) \arrow[r, "\text{oubli de }F'"] & V'.
\end{tikzcd}
\tag*{$(*)$}
\]
Dans le cas qui nous intéresse surtout ici, $F = e_V$, $V/F \longrightarrow V$ est un isomorphisme, et on obtient une factorisation canonique de foncteurs \msp{page 91}
\[
\begin{tikzcd}[row sep=large]
 & V \arrow[dl, "\tilde{f}=f/e_V"'] \arrow[dd, "f"] \\
\begin{array}{r} \upi_0(V/V') \\ = V'/f_!(e_V) \end{array} \arrow[dr, "\text{can.}"'] & \\
 & V'.
\end{tikzcd}
\tag*{$(**)$}
\]
Explicitons le $V'$-morphisme $\tilde{f}$ sur les fibres,
\[ \tilde{f}_{v'} : V_{v'} \longrightarrow \upi_0(V/V')_{v'} \defeq \pi_0(v'\backslash V), \tag*{$(***)$} \]
où dans le dernier membre, on a identifié l'ensemble $\pi_0(v'\backslash V)$ à la catégorie discrète qu'il définit. Pour expliciter cette flèche $(***)$, notons qu'on a une inclusion canonique
\[ V_{v'} \longrightarrow v'\backslash V \]
et pour toute catégorie $W$, un morphisme canonique (d'adjonction) $W \longrightarrow \pi_0(W)$, d'où un composé
\[ V_{v'} \longrightarrow v'\backslash V \longrightarrow \pi_0(v'\backslash V), \]
qui n'est autre que $\tilde{f}_{v'}$ dans $(***)$.

Il doit être facile de montrer que la factorisation $(**)$ de $f$ est caractérisée, à isomorphisme unique près, en tant que
% foot: Version 7 décembre 2002
\printedpage{65}
% head: A. Grothendieck : Dérivateurs, Ch. XII, Édition des Universités de Montpellier II et Paris VII
\msp{page 92}
factorisation dans $\Cat/V'$
\[
\begin{tikzcd}[column sep=small]
 & \tilde{V} \arrow[dr, "\pi"] & \\
V \arrow[ur, "\tilde{f}"] \arrow[rr, "f"'] & & V',
\end{tikzcd}
\]
par les propriétés suivantes :
\begin{enumerate}[label={}, leftmargin=*]
\item[a)] $\pi$ est fibrée à fibres discrètes, i.e. un isomorphisme local (i.e. $\tilde{V}$ est isomorphe dans $\Cat/V'$ à une catégorie $V/F'$, avec $F'$ dans $V'^{\wedge}$). % [sic] V/F' -> V'/F', une catégorie au-dessus de V' définie par F' dans V'^∧ est de la forme V'/F'
\item[b)] $\tilde{f}_!(e_V) = e_{\tilde{V}}$, i.e. pout tout $\tilde{v} \in \Ob \tilde{V}$, la catégorie $\tilde{v}\backslash V$ est 0-connexe : les colocalisés de $V$ au dessus de $\tilde{V}$ sont 0-connexes. % [sic] pout -> pour, faute de frappe
\end{enumerate}

Notons ici, en passant, que pour \textit{tout} morphisme $\pi : W' \longrightarrow W$ dans Cat qui soit fibré à fibres discrètes, i.e. tel que $W'$ soit isomorphe dans $\Cat/W$ à une catégorie de la forme $W/F$ ($F \in \Ob W^{\wedge}$), on a
\[ \pi_!(e_{W'}) = F, \qquad \text{i.e. } W' \underbrace{\ \simeq\ }_{\text{dans } \Cat/W} W/\pi_!(e_{W'}). \tag*{$(*)$} \]
Il faut vérifier que le foncteur
\[ w \longmapsto \pi_0(\underbrace{w\backslash(W/F)}_{\defeq w\backslash W/F}) \]
est canoniquement isomorphe à $F$. Mais l'inclusion
\[ \underbrace{W'_w}_{= F(w) \text{ (catégorie discrète)}} \qquad \longrightarrow w\backslash W' \]
\msp{page 93}
définit bien comme tantôt
\[ F(w) = W'_w \longrightarrow \pi_0(w\backslash W'), \]
il faut prouver que c'est bijectif. Or, $W'$ étant $W_\infty$-lisse sur $W$, on sait que l'inclusion $W'_w \hookrightarrow w\backslash W'$ est dans $W_\infty$, a fortiori induit un isomorphisme pour les $\pi_0$.

Ceci posé, si on a une factorisation de $f : V \longrightarrow V'$ en $f = \pi\tilde{f}$ ayant les propriétés a) b) de la page précédente, on voit déjà que l'on a
\[ f_!(e_V) \simeq \pi_!(\, \underbrace{\tilde{f}_!(e_V)}_{\simeq e_{\tilde{V}} \text{ par b)}} \,) \simeq \pi_!(e_{\tilde{V}}), \]
ce qui déjà détermine $\tilde{V}$ à isomorphisme canonique près (grâce à a) comme
\[ \tilde{V} \simeq V'/f_!(e_V) \defeq \upi_0(V'/V), \] % [sic] \underline{\pi}_0(V'/V) -> \underline{\pi}_0(V/V'), notation employée partout ailleurs (Définition 1, (**))
% foot: Version 7 décembre 2002
\printedpage{66}
% head: A. Grothendieck : Dérivateurs, Ch. XII, Édition des Universités de Montpellier II et Paris VII
et il ne doit pas être bien dur de se convaincre que $\tilde{f}$ n'est autre que la flèche canonique de la page 91.

\textbf{Proposition 8.} \textit{Pour qu'un objet $H$ de $U^{\wedge}$ soit totalement $W$-asphérique, il suffit que l'objet $U/H$ de} Cat \textit{soit totalement $W$-asphérique.}

\msp{page 94}
Supposons en effet cette dernière condition satisfaite et prouvons que $H$ est totalement $W$-asphérique, i.e. que pour tout $p : U \longrightarrow U'$, $U'/p_!(H)$ soit $W$-asphérique.
\[
\begin{tikzcd}[row sep=large]
V \defeq U/H \arrow[r, "\alpha"] \arrow[dr, "f"'] & U \arrow[d, "p"] \\
 & U'
\end{tikzcd}
\]
Soit $f$ le foncteur composé du diagramme ci-contre \texttt{[ci-dessus]}, on a $H \simeq \alpha_!(e_V)$ (par $(*)$ p. 92), d'où
\[ p_!(H) \simeq p_!\alpha_!(e_V) \simeq (p\alpha)_!(e_V) \simeq f_!(e_V), \]
qui est $W$-asphérique par hypothèse sur $V = U/H$.

\dgrblock{dgrb523}{The connected-component factorization, with its proof}
\begin{dcomment}{exposition}{p93}{The connected-component factorization, with its proof}
Grothendieck's relative $\underline\pi_0$ turns a map of categories into a genuine family of sets over its target. This is the right construction for understanding what ordinary Kan extension forgets. Let $f:A\to B$, put
\[
F(b)=\pi_0(b\downarrow f),\qquad E=B/F,
\]
and let $\pi:E\to B$ be the projection. Precomposition with a morphism $b'\to b$ gives the restriction map $F(b)\to F(b')$, so $F$ is a presheaf. The factorization map
\[
j:A\longrightarrow E,\qquad
j(a)=\bigl(f(a),[(a,\mathrm{id}_{f(a)})]\bigr)
\]
is functorial: a morphism $a\to a'$ identifies the evident two objects in $f(a)\downarrow f$. We have $f=\pi j$.

To prove the missing connectedness assertion, fix $e=(b,c)\in E$. An object of $e\downarrow j$ is an $a\in A$ together with $\beta:b\to f(a)$ such that the component of $(a,\beta)$ is $c$. Its morphisms are exactly the morphisms of $b\downarrow f$ between those objects. Thus
\[
e\downarrow j\ \cong\ (b\downarrow f)_c,
\]
the component indexed by $c$. It is nonempty and connected by definition. In the usual ordinary-category terminology, $j$ is a \emph{final functor}. No higher contractibility is being asserted.

We also need to recover a discrete fibration from this procedure. Suppose $\pi:E\to B$ is the category of elements of a presheaf $F$. For $(e,\beta:b\to\pi e)$ in $b\downarrow\pi$, the unique lift of $\beta$ into $e$ has source $(b,F(\beta)(e))$. This assignment gives a functor to the discrete fibre $E_b$. Its composite with the fibre inclusion is naturally connected to the identity by the lifted arrows, and its restriction to the fibre is the identity. Consequently
\[
\pi_0(b\downarrow\pi)\cong F(b),\qquad
\pi_!1_E\cong F.
\]
The same argument proves a homotopy equivalence between the comma category and the discrete fibre, which explains the appeal to smoothness on p.~93 without requiring that theory here.

Finally, suppose $f=\rho k$ is another factorization with $\rho:D\to B$ discrete and $k:A\to D$ final. By the Kan-extension formula $k_!1_A=1_D$, and hence
\[
f_!1_A\cong\rho_!k_!1_A\cong\rho_!1_D.
\]
The preceding recovery calculation identifies $D$ over $B$ with $B/F=E$. This identification also respects the first factor. Indeed, for any presheaf $G$ on $B$, a functor $A\to B/G$ over $B$ is the same as a natural transformation $1_A\to f^*G$, and adjunction identifies these with natural transformations $f_!1_A\to G$. For $G=F$, the canonical first factor corresponds to $\mathrm{id}_F$; the identification obtained from $\rho_!1_D\cong F$ carries $k$ to that same map. This proves uniqueness with its compatibility, rather than only determining the middle category as an abstract object.

The result is the discrete-fibration form of the comprehensive factorization of \citet{StreetWalters1973}. A complete direct proof in the opposite variance is in \citet[\S4.7, pp.~73--75, especially (4.77)--(4.79)]{Kelly1982}; Kelly also identifies an invalid Kan-extension inference in the original proof and replaces it by the component calculation used above. The theorem predates these manuscript pages. Its role here is to distinguish the final first factor from a homotopy final map, whose comma categories would have to be weakly contractible.
\end{dcomment}


\dgrblock{dgrb524}{From comprehensive factorization to homotopy families}
\begin{dcomment}{later}{p93}{From comprehensive factorization to homotopy families}
The geometric interpretation of $\underline\pi_0$ has continued to be useful independently of derivators. \citet[\S1.3, Lemma~1.6, and Theorem~1.7]{BergerKaufmann2017} place the two ordinary comprehensive factorizations in a general theory of comprehension schemes: presheaves supply discrete fibrations, and the unit of comprehension supplies their universal approximation. Their Proposition~1.10 and Remark~1.11 explain the associated base-change and Frobenius properties. These are structural statements about the factorization, not claims that its first factor is a weak homotopy equivalence.

For the comparison with higher categories, the decisive change is to keep the homotopy type of each comma category instead of only its set of components. The homotopy-finality criterion requires $B(b\downarrow f)$ to be contractible for every $b$, as in \citet[Theorem~4.1.3.1]{Lurie2009HTT}. Here $j$ only makes these spaces connected. Thus comprehensive factorization is a useful discrete approximation, but it cannot replace a construction of a homotopy Kan extension. The examples below measure precisely what is lost when this approximation is used before forming a total category.
\end{dcomment}


La réciproque de \texttt{[la]} prop. 8 est-elle valable, i.e. si $H$ est totalement $W$-asphérique, en est-il de même de $V = U/H$? La difficulté provient du fait que cela signifie que pour \textit{toute} flèche $f : V \longrightarrow V'$ dans Cat de source $V$, on ait $f_!(e_V)$ $W$-asphérique, alors que l'hypothèse équivaut à le postuler seulement pour les flèches qui se factorisent par $U$. Notons que si on ne fait pas cette hypothèse sur $f$, mais qu'on suppose au lieu de cela que $f$ s'insère dans un carré commutatif \msp{page 95}
\[
\begin{tikzcd}[row sep=large]
V \arrow[r, "\alpha"] \arrow[d, "f"'] & U \arrow[d, "p"] \\
V' \arrow[r, "\alpha'"] & U',
\end{tikzcd}
\]
avec $\alpha'$ lui aussi un isomorphisme local, alors on peut conclure que $\upi_0(f) = \upi_0(V/V')$ est bien asphérique. En effet, on a un isomorphisme canonique dans $\Cat/U'$,
\[ \upi_0(\alpha'f) \simeq \upi_0(f), \qquad \text{i.e. } \ \upi_0(V/V') \simeq \upi_0(V/U'), \]
du fait p. ex. que la factorisation canonique de $f$
% foot: Version 7 décembre 2002
\printedpage{67}
% head: A. Grothendieck : Dérivateurs, Ch. XII, Édition des Universités de Montpellier II et Paris VII
\[
\begin{tikzcd}[row sep=normal]
 & V \arrow[dl, "\tilde{f}"'] \arrow[dd, "f"] & \\
\upi_0(V/V') \arrow[dr, "\pi"'] & & \\
 & V' \arrow[dr, dotted, "\alpha'"'] & \\
 & & U'
\end{tikzcd}
\]
fournit une factorisation du même type pour $\alpha'f$, en
\[ \alpha'f = (\alpha'\pi) \circ \tilde{f}. \]
Ceci dit, on peut appliquer l'hypothèse à $p$ et on trouve bien la $W$-asphéricité de $\upi_0(f) \simeq \upi_0(\alpha'f) = \upi_0(p\alpha)$.

\msp{page 96}
L'idée naturelle, pour une réciproque à \texttt{[la]} prop. 8, serait de prendre pour $U'$ la somme amalgamée de $U$ et $V'$ sur $V$, ce qui marchera pourvu qu'on sache que $\alpha'$ est bel et bien un isomorphisme local tout comme $\alpha$. C'est bien le cas p. ex. si $V \longrightarrow U$ est une immersion ouverte, i.e. $H$ un sous-préfaisceau du préfaisceau final. Mais je crains que ce ne soit pas vrai en général - ça a l'air un peu trop beau pour être vrai! C'est en effet déjà faux dans le cas
\[
\begin{tikzcd}[row sep=huge]
V = E_G \arrow[r, "\alpha\ \text{can}"] \arrow[d, "f"'] & B_G = U \\
V' = B_{G'}, &
\end{tikzcd}
\]
dans le cas p. ex. où $G = \{\pm 1\}$, groupe à deux éléments, et où l'élément $\alpha$ de $G'$, image par $f$ de l'une des deux flèches non identiques de $E_G$ est $\neq 1$ \texttt{[plutôt de carré} $\neq 1$\texttt{]}. Si $H'$ est le sous-groupe invariant engendré par $\alpha$ \texttt{[plutôt par} $\alpha^2$\texttt{]} dans $G'$, la somme amalgamée s'identifie à $B_{G'/H'}$, or $B_{G'} \longrightarrow B_{G'/H'}$ n'est \textit{pas} un isomorphisme local (car $(B_{G'})/t_0 \longrightarrow (B_{G'/H'})/t'_0$ n'est pas bijectif sur les objets, s'identifiant à $G' \overset{q}{\longrightarrow} G'/H' \ \ldots$).

Mais d'autre part la condition de totale $W$-asphéricité sur un objet $H$ de $U^{\wedge}$ paraît déjà à tel point exigeante, que je me vois mal con\msp{page 97}struire un exemple où elle est satisfaite, sans que $U/H$ elle-même ne soit totalement $W$-asphérique.

\dgrblock{dgrb525}{The converse for subterminal presheaves and its obstruction}
\begin{dcomment}{exposition}{p96}{The converse for subterminal presheaves and its obstruction}
Proposition~8 compares two quantifiers. Write $V=U/H$ and $\alpha:V\to U$. Total asphericity of $V$ tests every outgoing functor $f:V\to C$. Total asphericity of $H$ tests the special family $p\alpha$, with $p:U\to B$. The proposition follows because $\alpha_!1_V=H$, but reversing it requires a way of placing an arbitrary $f$ in that special family without changing its comprehensive middle category.

The square criterion on p.~95 does this whenever one can complete
\[
\begin{tikzcd}
V\arrow[r,"\alpha"]\arrow[d,"f"']&U\arrow[d,"p"]\\
C\arrow[r,"\beta"']&B
\end{tikzcd}
\]
with $\beta$ a discrete fibration. To justify the assertion completely, factor $f$ as a final functor $V\to E$ followed by a discrete fibration $E\to C$. Since discrete fibrations compose, the same $E$ is the comprehensive middle category of $\beta f=p\alpha$. Its category is therefore $W$-aspherical when $H$ is totally $W$-aspherical.

If $H$ is subterminal, $\alpha$ is a sieve inclusion. Form the ordinary categorical pushout $B=U\amalg_V C$. Theorem~1 preserves the sieve inclusion on the other side, so $\beta:C\to B$ is again a sieve, in particular a discrete fibration. The preceding argument applies to every $f:V\to C$. This supplies the full proof of the source's subterminal case: for such $H$, the total asphericity of $H$ and of $U/H$ are equivalent.

The groupoid example identifies why the same proof cannot be extended merely by replacing ``sieve'' with ``discrete fibration.'' Let $G=C_2$ and let $E_G$ be the indiscrete groupoid on two objects. Its projection to $BG$ sends both nonidentity arrows to the unique nontrivial element $g$. A functor $E_G\to BG'$ sends one of these arrows to an element $a\in G'$ and the inverse arrow to $a^{-1}$. In a cocone on this span, the two objects must become equal, and compatibility with the projection to $BG$ forces
\[
a=a^{-1},\quad\text{equivalently }a^2=1.
\]
There are no additional relations. The pushout is therefore
\[
B\bigl(G'/\langle\!\langle a^2\rangle\!\rangle\bigr).
\]
For example, with $G'=\mathbb Z$ and $a$ its generator, the induced map is $B\mathbb Z\to BC_2$. A discrete fibration between one-object groupoids must lift every arrow uniquely into the sole object, so its map of morphism sets must be bijective. This quotient map is not. The bracketed corrections in the edition, replacing $a$ by $a^2$, are thus mathematically necessary.

This counterexample defeats the proposed pushout method; it does not by itself refute the general converse of Proposition~8. That converse is false: a presheaf with three elements on a monoid with two elements is given in the commentary at ms.~97. The subterminal case is proved, and the obstruction to the broader argument is explicit.
\end{dcomment}


L'existence d'objets totalement $W$-asphériques dans Cat n'est pas tellement évidente a priori. Qu'il y en ait un bon paquet est assuré par la
% foot: Version 7 décembre 2002

% ---- from Der-ChXIIb-012_p026-030.tex
\setcounter{footnote}{65}

\printedpage{68}
% head: A. Grothendieck : Dérivateurs, Ch. XII, Édition des Universités de Montpellier II et Paris VII

\noindent\textbf{Proposition 9.}

\begin{enumerate}[label=\textit{\alph*)}]
\item \textit{Tout objet $V$ de} Cat \textit{ayant un objet final est totalement $W$-asphérique - donc tout préfaisceau d'ensembles $H$} représentable \textit{sur une catégorie $U$ (la représentabilité signifiant justement que $U/H$ a un objet final) est totalement $W$-asphérique.}

\item (\footnote{!! Sans doute \textit{faux} tel quel, cf. rectification page suivante.}) \textit{(Supposons} W(7 bis) \textit{satisfait.) Soient $V$ dans} Cat, \textit{$V'$, $V''$ deux ouverts de réunion $V$. Si $V'$, $V''$, $V'\cap V''$ sont totalement $W$-asphériques, il en est de même de $V$.}
\end{enumerate}

\noindent\textsc{Démonstration} de a). Soit $V$ avec objet final et $f : V \longrightarrow V'$, d'où la factorisation
\[
\begin{tikzcd}[column sep=small, row sep=small]
 & \tilde{V} = \upi_0(V/V') \arrow[dr, "\pi"'] & \\
V \arrow[ur, "\tilde{f}"] \arrow[rr, "f"'] & & V'
\end{tikzcd}
\]
de $f$. On démontre sans mal (\footnote{Voici une démonstration simple : On a pour toute $V$ dans Cat $e_V = \varinjlim_{v \text{ dans } V} v$ (limite inductive du foncteur d'inclusion de $V$ dans $V^\wedge$), car $\forall\, x$ dans $V$
\[
\varinjlim_{v \text{ dans } V} \Hom(x,v) \simeq \pi_0(x\backslash V) = e.
\]
Comme $f_!$ commute aux $\varinjlim$ , on a donc $f_!(e_V) = \varinjlim_{v \text{ dans } V} f(v)$. Si $v_0$ est un objet final de $V$, cela montre que $f_!(e_V) = f(v_0)$, q.e.d.}) (histoire de cofinalité) ceci :

\noindent\textbf{Corollaire.} \textit{L'image par $\tilde{f}$ d'un objet final de $V$ est un objet final de $\tilde{V} = \upi_0(V/V')$.}

Comme un objet de Cat avec objet final est $W$-asphérique par W(3), cela établit a).

Prouvons b). J'ai envie de dire que si on a $f : V \longrightarrow W$ dans Cat, alors $\upi_0(V/W)$ est réunion de deux ouverts isomorphes à $\upi_0(V'/W)$, $\upi_0(V''/W)$, dont l'intersection est isomorphe à $\upi_0(V'\cap V''/W)$. Cela signifierait ni plus ni
\msp{page 98}
moins que ceci : $f_!$ transforme les monomorphismes $V' \longrightarrow V$, $V'' \longrightarrow V$ (interprétés comme flèches dans $V^\wedge$, de but $e_V$) en monomorphismes, et le produit fibré $V'\cap V'' = V' \times_V V''$ en le produit fibré correspondant dans $W^\wedge$. Mais le foncteur $f_!$ commute aux $\varinjlim$ quelconques (étant un adjoint à gauche), et je ne vois aucune raison a priori pour qu'il commute (disons) aux produits fibrés (ce qui impliquerait aussi qu'il transforme monomorphismes en itou). Si je me rappelle bien, c'est ultrafaux en général. Donc je doute fort à présent que l'énoncé b) soit vrai, sans hypothèse supplémentaire. Pour assurer la monomorphie, on songe à imposer l'existence de \textit{rétractions} de $V$ sur $V'$, $V''$ et de $V'$ ou $V''$ sur $V'\cap V''$.

\dgrblock{dgrb526}{Total asphericity tested on final images}
\begin{dcomment}{development}{p97}{Total asphericity tested on final images}
Grothendieck's search becomes more transparent if we express the quantifier through the first half of comprehensive factorization. The following elementary consequence is our formulation of what the construction already permits; it is not a claim of priority. We checked the factorization accounts of Street--Walters, Kelly, and Berger--Kaufmann, which establish the factorization but do not state this characterization of the manuscript's total-asphericity class.

\noindent\textbf{Final-image criterion.}\ours\quad
For any localizer $W$ for which asphericity is invariant under isomorphism, a small category $A$ is totally $W$-aspherical if and only if every ordinary final functor $A\to C$ has $W$-aspherical target. Moreover, the class of totally $W$-aspherical categories is closed under ordinary final functors.

\emph{Proof.} If $u:A\to C$ is final, then $u_!1_A=1_C$. Hence the category of elements used in the definition of total asphericity is $C/1_C\cong C$, proving the necessity. Conversely, factor an arbitrary $f:A\to B$ comprehensively as
\[
A\xrightarrow{j}E=B/(f_!1_A)\xrightarrow{\pi}B.
\]
The first factor is final, so the assumed condition says that $E$ is $W$-aspherical. This is exactly the required test for $f$.

For the closure assertion, suppose $A$ is totally $W$-aspherical and $u:A\to C$ is final. For every $g:C\to D$,
\[
g_!1_C\cong g_!u_!1_A\cong(gu)_!1_A.
\]
Its category of elements is $W$-aspherical by the hypothesis on $A$. Thus $C$ is totally $W$-aspherical.\hfill$\square$

There is an equivalent extremal description which is sometimes useful for constructing examples. Among subclasses of the $W$-aspherical categories which are closed under taking targets of ordinary final functors, the totally $W$-aspherical categories form the largest one. They themselves have that closure property, and they are aspherical by taking $f=\mathrm{id}$. Conversely, if $A$ belongs to any such subclass, every comprehensive middle category of an outgoing map from $A$ belongs to it, hence is aspherical. The criterion just proved applies.

The result also separates two kinds of positive examples. If $A$ has final object $a_0$, then for every $f:A\to B$,
\[
f_!1_A\cong B(-,f(a_0)),\qquad
B/(f_!1_A)\cong B/f(a_0),
\]
whose final object is $(f(a_0),\mathrm{id})$. This is Proposition~9(a) in an explicit form. Filtered categories provide the further examples already explained in the commentary on Chapter~II, manuscript pp.~32--35, through filtered colimits of representables. For instance $(\mathbb N,\leq)$ has no final object, so the conjecture in footnote~73 cannot survive even in ordinary homotopy theory. The new criterion explains what both examples must control: all final images, not just their own classifying spaces.
\end{dcomment}


\dgrblock{dgrb527}{A total presheaf whose category of elements is not total}
\begin{dcomment}{development}{p97}{A total presheaf whose category of elements is not total}
Grothendieck expects that Proposition~8 has no converse outside the subterminal case, and proposes to look for a totally $W$-aspherical presheaf $H$ whose category of elements is not totally $W$-aspherical. Such a presheaf exists on the smallest category that is not a preorder and not a group. We found no source for the example; it is given without a claim of priority.

\textbf{Proposition (the converse of Proposition~8 fails).}\ours
\textit{Let $U$ be the category with one object $*$ and $\operatorname{End}(*)=\{1,e\}$, $e^2=e$. Let $H$ be the presheaf on $U$ with $H(*)=\{i,h_1,h_2\}$, on which $e$ acts by the constant map with value $i$. Then $H$ is totally $W$-aspherical for every fundamental localizer $W$. The category $U/H$ is not totally $W$-aspherical for $W=W_\infty$, for the homology equivalences $W_R$ with coefficients in a nonzero commutative ring $R$, and for the $n$-equivalences with $n\geq2$.}

\textit{Proof.}
Let $R=U(-,*)$ be the representable presheaf. Its elements are $1$ and $e$, and $e$ acts on them by $x\mapsto xe$, the constant map with value $e$. The endomorphism $\varepsilon:x\mapsto ex$ of $R$ is idempotent with image the subpresheaf $Q=\{e\}$. The presheaf $H$ is the pushout $R\amalg_QR$ of two copies of $R$ along $Q$: the element $e$ is $i$, and the two copies of $1$ are $h_1,h_2$.

Let $p:U\to C$ be any functor to a small category, and put $c=p(*)$, $d=p(e)$, so that $d^2=d$. The functor $p_!$ preserves colimits, sends $R$ to $C(-,c)$, and sends $\varepsilon$ to the endomorphism $\alpha\mapsto d\alpha$. Every functor preserves the splitting of an idempotent. Hence $p_!Q$ is the subpresheaf
\[
 Q_d(v)=\{\alpha:v\to c\mid d\alpha=\alpha\}\subseteq C(v,c),
\]
with its inclusion, and
\[
 p_!H\cong C(-,c)\amalg_{Q_d}C(-,c).
\]
This pushout is formed objectwise along injections. The category of elements $E'=C/p_!H$ is therefore the union of two full sieves $S_1,S_2$, each isomorphic to $C/c$, with $S_1\cap S_2=C/Q_d$. A chain of arrows of $E'$ lies in the sieve that contains its last object, so $NE'=NS_1\cup NS_2$ and $NS_1\cap NS_2=N(C/Q_d)$. Each $S_j$ has the terminal object $(c,\mathrm{id}_c)$. The category $C/Q_d$ has the object $(c,d)$, and the arrows $\alpha:(v,\alpha)\to(c,d)$ are the components of a natural transformation from the identity functor to the constant functor at $(c,d)$: for an arrow $g:(v,\alpha)\to(w,\beta)$ one has $\beta g=\alpha$. Thus $BS_1$, $BS_2$ and their intersection are contractible, and so is $BE'$. Since $p$ was arbitrary, $H$ is totally $W_\infty$-aspherical. Every fundamental localizer contains $W_\infty$ \citep[Th\'eor\`eme~2.2.11]{Cisinski2004Minimal}, which gives the first assertion.

The category $E=U/H$ has the three objects $i,h_1,h_2$ and, besides identities, the arrows
\[
 e_i:i\to i,\qquad u_1:i\to h_1,\qquad u_2:i\to h_2,
 \qquad e_i^2=e_i,\quad u_je_i=u_j.
\]
Let $q:E\to\Psi$ be the functor to the outgoing fork $\Psi=(a\leftarrow z\to b)$ of the commentary at ms.~100 which sends $i,h_1,h_2$ to $z,a,b$ and $e_i$ to the identity. It is bijective on objects and full. For an object $y$ of $\Psi$ with preimage $\tilde y$, every object $(x,\,y\to q(x))$ of $y\backslash E$ receives an arrow from $(\tilde y,\mathrm{id}_y)$, obtained by lifting $y\to q(x)$. So $q$ is final. The inclusion $f:\Psi\to C$ into the four-object category of the same commentary is final and $BC\simeq S^2$. The composite $fq:E\to C$ is a final functor whose target is not aspherical for any of the localizers named. By the final-image criterion above, $U/H$ is not totally aspherical for them.\hfill$\square$

The objects $h_1,h_2$ of $U/H$ have no common successor; this is the property that fails, and it cannot be seen by functors that factor through $U$. For the localizers $W_0$ and $W_1$ (bijections on $\pi_0$; bijections on $\pi_0$ and isomorphisms on all $\pi_1$) the converse of Proposition~8 does hold. A final functor induces a bijection on connected components, and a final image of a $1$-connected category is $1$-connected (the argument with coverings at ms.~101). If $H$ is totally $W_0$-aspherical, the test $p=\mathrm{id}$ shows that $U/H$ is nonempty and connected, and then so is every final image of $U/H$; for $W_1$ the same reasoning applies with $1$-connectedness.
\end{dcomment}


Je viens de vérifier la chose, et donne le
% foot: Version 7 décembre 2002
\printedpage{69}
% head: A. Grothendieck : Dérivateurs, Ch. XII, Édition des Universités de Montpellier II et Paris VII
\textbf{Corollaire 1} (rectification de \texttt{[la]} prop. 9, partie b) (\footnote{NB Il y a un énoncé légèrement plus général, en partant d'une catégorie $U$ et d'un diagramme
$\begin{tikzcd}[ampersand replacement=\&, row sep=small]
V'\cap V'' \arrow[r] \arrow[d] \& V' \arrow[d] \\
V'' \arrow[r] \& V
\end{tikzcd}$
\textit{dans} $U^\wedge$ : si $V'$, $V''$, $V'\cap V''$ sont totalement $W$-asphériques, il en est de même de $V$.}) : \textit{L'énoncé est valable si on suppose qu'il existe une} rétraction \textit{$r'_0$ de $V''$ sur $V'\cap V''$, et de même une rétraction $r''_0$ de $V'$ sur $V'\cap V''$.} (\footnote{Corollaire rectifié faux également - les rétractions ne sont pas des flèches de $V^\wedge$, et ne passent dans $p_!$.}) % [resolved: see DER_XIIb-a_corrections.csv]
\[
\begin{tikzcd}
V'\cap V'' \arrow[r, "i''_0"'] \arrow[d, "i'_0"] & V' \arrow[d, "i'"] \arrow[l, dotted, bend right=50, "r''_0"'] \\
V'' \arrow[r, "i''"] \arrow[u, dotted, bend left=50, "r'_0"] & V \arrow[u, dotted, bend right=50, "r'"'] \arrow[l, dotted, bend left=50, "r''"]
\end{tikzcd}
\]

(\textbf{NB} Nous verrons plus tard qu'il suffit de l'existence d'\textit{une} de ces rétractions pour assurer la conclusion.)

\textsc{Démonstration}. Comme $V$ est somme amalgamée de $V'$ et $V''$, il existe une rétraction $r'$ de $V$ sur $V'$ (donnée par $\mathrm{id}_{V'}$ sur $V'$, par $i''_0r'_0$ sur $V''$) qui induit $r'_0$ sur $V''$, et de même il existe
\msp{page 99}
une rétraction $r''$ de $V$ sur $V''$, induisant $r''_0$ sur $V'$. L'existence de ces rétractions assure déjà que pour tout foncteur $f : V \longrightarrow W$, le diagramme transformé de
\[
\begin{tikzcd}[row sep=large]
V'\cap V'' \arrow[r] \arrow[d] & V' \arrow[d] \\
V'' \arrow[r] & V
\end{tikzcd}
\]
(intérprété comme diagramme dans $V^\wedge$) est formé de monomorphismes (\footnote{\textbf{faux,} déjà dans le cas où $V = \Phi$, $V'$, $V''$ ses deux ouverts canoniques isomorphes à $\bDelta^1$, et $V'{\cap}V'' = e$.}) % [sic] intérprété -> interprété, misspelt accent
\[
\begin{tikzcd}[row sep=large]
f_!(V'\cap V'') \arrow[r, hook] \arrow[d, hook] & f_!(V') \arrow[d, hook] \\
f_!(V'') \arrow[r, hook] & f_!(V),
\end{tikzcd}
\]
et de plus que $f_!(V') \cap f_!(V'')$ (qui contient $f_!(V'\cap V'')$) lui est en fait \textit{égal.} Il en résulte que l'on a un diagramme à la fois cartésien et cocartésien d'immersions ouvertes dans Cat
% foot: Version 7 décembre 2002
\printedpage{70}
% head: A. Grothendieck : Dérivateurs, Ch. XII, Édition des Universités de Montpellier II et Paris VII
\[
\begin{tikzcd}[row sep=large]
W/f_!(V'\cap V'') \arrow[r, hook] \arrow[d, hook] & W/f_!(V') \arrow[d, hook] \\
W/f_!(V'') \arrow[r, hook] & W/\Bigl(f_!(V) \overset{\text{déf}}{=} f_!(e_V)\Bigr),
\end{tikzcd}
\]
lequel s'identifie aussi à
\[
\begin{tikzcd}[row sep=large]
\upi_0(V'\cap V''/W) \arrow[r, hook] \arrow[d, hook] & \upi_0(V'/W) \arrow[d, hook] \\
\upi_0(V''/W) \arrow[r, hook] & \underline{\pi_0}(V/W).
\end{tikzcd}
\]

\dgrblock{dgrb528}{Why a categorical retraction cannot be sent through this Kan extension}
\begin{dcomment}{exposition}{p98}{Why a categorical retraction cannot be sent through this Kan extension}
Grothendieck tries to protect a union from collapsing under $f_!$ by splitting its inclusions. Splitting is indeed preserved by every functor, but the retraction must be a morphism in the category on which that functor acts. Here $f_!$ acts on $V^{\wedge}$, whereas the retractions just constructed are functors between full subcategories of $V$.

For a sieve $S\subseteq V$, the associated subterminal presheaf $1_S\subseteq1_V$ takes the value one on objects of $S$ and the empty set elsewhere. If $S$ omits an object of $V$, there cannot be a map $1_V\to1_S$: at that object it would have to be a function from a singleton to the empty set. Thus a proper sieve, even if it admits a categorical retraction, does not define a split subobject of the terminal presheaf. A functor $r:V\to S$ moves the base object; a morphism in $V^{\wedge}$ must stay over it. This is exactly the defect identified by the source's corrective footnote.

What would suffice is a genuinely different hypothesis. Given a union $H=H_1\cup H_2$ in a presheaf category, suppose that for every $p$ the images of its three inclusion maps remain monomorphisms and the square remains Cartesian. Since $p_!$ preserves the pushout, the resulting categories of elements then form a union of two sieves with the expected intersection. The ordinary forward homotopy gluing lemma proves total $W_\infty$-asphericity of $H$ when the other three presheaves are totally $W_\infty$-aspherical. This is the precise exactness condition the proposed proof needs. Categorical retractions provide none of it.
\end{dcomment}


\msp{page 100}
L'hypothèse sur $V'$, $V''$, $V'\cap V''$ implique que les trois sommets du carré autres que $\upi_0(V/W)$ sont $W$-asphériques, donc par W(7 bis) il en est de même de ce dernier, q.e.d.

\textbf{Corollaire 2} (\footnote{Plus généralement, si $I$ est une catégorie \textit{ordonnée finie,} ayant un objet initial, alors}). \textit{Moyennant} W(7 bis)\textit{, la catégorie} $\Psi = \begin{tikzcd}[ampersand replacement=\&, row sep=small, column sep=small, baseline=(B.base)]
 \& a \\
|[alias=B]| 0 \arrow[ur] \arrow[dr] \& \\
 \& b
\end{tikzcd}$ \textit{(ou $0 \longleftarrow 1 \longrightarrow 2$) est totalement $W$-asphérique} (\footnote{\textbf{Faux,} $\Psi$ n'est \textit{pas} totalement $W$-asphérique.}).

On applique le corollaire 1 à $\Psi$ et à ses deux sous-catégories ouvertes, $\Psi/a$, $\Psi/b$ (isomorphes à $\bDelta^1$), et à leur intersection $\{0\} \simeq \bDelta^0 = e$. Comme ces trois catégories ont un objet final (contrairement à $\Psi$ elle-même, qui n'a qu'un objet initial \dots), elles sont totalement $W$-asphériques. Le corollaire s'applique, l'existence des rétractions $r'_0$, $r''_0$ étant ici évidente.

C'est surtout la partie b) (rectifiée) de la page 9 qui assure l'existence de catégories totalement $W$-asphériques autres que celles ayant un objet final (\footnote{Je conjecture à présent que ce sont les seules! (\footnotemark)}\footnotetext{Faux, il y a les catégories filtrantes (si $W$ satisfait l'axiome des limites).}). Il est vrai qu'en sens inverse il faudrait quand même s'assurer qu'un objet $X$ $W$-asphérique de Cat n'est pas \textit{toujours} totalement $W$-asphérique! Le contre-exemple n'est pas tellement évident finalement -
\msp{page 101}
j'ai passé des heures avant d'en trouver un et d'y voir un peu clair! Il suffit en somme d'exhiber un exemple d'une flèche dans Cat
\[
f : X \longrightarrow Y
\]
telle que
% foot: Version 7 décembre 2002
\printedpage{71}
% head: A. Grothendieck : Dérivateurs, Ch. XII, Édition des Universités de Montpellier II et Paris VII
\begin{enumerate}[label=\alph*)]
\item $\underline{\pi_0}(f) \overset{\text{déf}}{=} \upi_0(X/Y) \xrightarrow{\ \sim\ } Y$, i.e. les $\pi_0(y\backslash X)$ \texttt{[plutôt les} $y\backslash X$\texttt{]} sont 0-connexes (pour $y \in \Ob Y$),
\end{enumerate}

\dgrblock{dgrb529}{The finite fork counterexample, computed completely}
\begin{dcomment}{examples}{p100}{The finite fork counterexample, computed completely}
Grothendieck eventually finds a very small witness to the failure of the preceding gluing claims. The construction is in his Appendix to Chapter~XII, manuscript pp.~10--12, where the final corrective footnote identifies the homotopy type as $S^2$ \citep{GrothendieckDerivateursXIIAppendix}. We work it out here because it explains at once the failure of total asphericity and the failed use of retractions.

Let $C$ have four objects $x,z,a,b$ and the nonidentity morphisms generated by
\[
\begin{tikzcd}[column sep=large]
x\arrow[r,shift left=.5ex,"\alpha"]\arrow[r,shift right=.5ex,"\beta"']&z\arrow[r,"r"]\arrow[d,"s"']&a\\
&b&
\end{tikzcd}
\qquad r\alpha=r\beta,\quad s\alpha=s\beta.
\]
Besides the four displayed arrows there are just the two composites $x\to a$ and $x\to b$. Let $A$ be its full subcategory on $z,a,b$, the outgoing fork $\Psi$, and let $f:A\hookrightarrow C$ be the inclusion. The category $A$ has initial object $z$, so $BA$ is contractible.

We first compute $f_!1_A$. At $a$ and $b$ the outgoing comma category is a point. At $z$ it has initial object $(z,\mathrm{id})$. At $x$ it has four objects
\[
(z,\alpha),\ (z,\beta),\ (a,r\alpha),\ (b,s\alpha),
\]
and one arrow from each of the first two to each of the last two. Its nerve is a four-edge circle. In particular it is nonempty and connected. Thus every $c\downarrow f$ is connected, and
\[
f_!1_A=1_C,\qquad \underline\pi_0(A/C)=C.
\]
The map $f$ is final in the ordinary sense, although its comma category at $x$ is not contractible.

The homotopy type of $C$ can be read without a spectral sequence. Let $K$ be the full subcategory on $x,z$, with its two parallel arrows. Its nerve realizes as a circle. The full subcategories $C_a$ on $x,z,a$ and $C_b$ on $x,z,b$ each adjoin a final object to $K$. Their nerves are therefore cones on $NK$. They cover $NC$, with intersection $NK$, because there are no arrows between $a$ and $b$. Consequently
\[
BC=BC_a\cup_{BK}BC_b\cong\Sigma BK\simeq S^2.
\]
Here the union description is literal at the level of nerves. Hence $A$ is not totally $W_\infty$-aspherical: its direct image along $f$ has category of elements $C$, with $H_2(BC;\mathbb Z)\cong\mathbb Z$.

Now return to the proposed retraction repair. Inside $A$, put $A_a=\{z\to a\}$ and $A_b=\{z\to b\}$. Both have final objects, their intersection is $\{z\}$, and both admit retractions to that intersection. Proposition~9(a) therefore gives total asphericity of all three pieces. Nevertheless their union is $A$, which we have just shown is not total. The false step can be inspected at a single value. Write $y_c=C(-,c)$. Since left Kan extension carries a representable to the corresponding representable,
\[
f_!(1_{\{z\}})=y_z,\quad
f_!(1_{A_a})=y_a,\quad
f_!(1_{A_b})=y_b.
\]
The equality $1_{A_a}=A(-,a)$ follows because $a$ receives exactly one arrow from $z$ and none from $b$, and similarly for $b$. Evaluating the first inclusion after Kan extension at $x$ yields
\[
\{\alpha,\beta\}=C(x,z)\longrightarrow C(x,a)=\{r\alpha=r\beta\}.
\]
It is not injective. The other leg has the same defect. Their pushout is the singleton $1_C(x)$, exactly as preservation of colimits predicts. The putative intersection square has failed before any homotopy argument is applied.

The same finite construction also explains the two kinds of pushout failure announced on p.~81. Adjoin a final object to $A$, giving $A^{\triangleright}$. Pushing out $A\subset A^{\triangleright}$ along $f$ adjoins a final object to $C$, because $f_!1_A=1_C$. The upper inclusion is an aspherical weak equivalence; the lower inclusion $C\to C^{\triangleright}$ is not, since $BC\simeq S^2$.

To make the other leg a weak equivalence as well, replace the target of $f$ by $C^{\triangleright}$, with final point $c_+$. The composite $p:A\to C^{\triangleright}$ is a weak equivalence between contractible categories. Its direct image of $1_A$ is the subterminal presheaf supported on $C$: it is one on every object of $C$ and empty at $c_+$. The resulting pushout $P$ is $C$ with two distinct final cone points, $c_+$ and the newly attached point. Hence
\[
BP\simeq\Sigma BC\simeq S^3.
\]
Both $A\to A^{\triangleright}$ and $p$ are weak equivalences, while neither map from $C^{\triangleright}$ or $A^{\triangleright}$ to $P$ is. This proof uses the explicit nerve, so it does not depend on the invalid reverse implication of the one-point corollary on p.~85.
\end{dcomment}


\dgrblock{dgrb530}{Testing the retracted assertions under their stated axiom}
\begin{dcomment}{development}{p100}{Testing the retracted assertions under their stated axiom}
The fork example proves that ordinary weak contractibility is insufficient for total asphericity. There is a logical point still to check: several of the source's retracted conclusions explicitly assume W(7 bis), whereas $W_\infty$ does not satisfy that axiom. To refute those conditional conclusions themselves, rather than only their proposed proof, we need an example in a weak equivalence class for which the axiom holds. Fixed-coefficient homology provides such a test. The homological tools are standard; the application below is our explicit check of these conditional assertions, with no claim of priority.

\textbf{A homological counterexample with W(7).}\ours\quad
Let $R$ be a nonzero commutative ring, and let $W_R$ consist of functors inducing isomorphisms on all homology groups of their nerves with coefficients in $R$. Then $W_R$ is a fundamental localizer satisfying the full W(7). Nevertheless a union of two totally $W_R$-aspherical sieves with totally $W_R$-aspherical intersection need not be totally $W_R$-aspherical, even when both pieces retract onto their intersection.

\emph{Verification of the test class.} Identities, two-out-of-three, and closure under retracts follow by applying homology. A category with a terminal object has contractible nerve and hence maps to the terminal category by an arrow in $W_R$. We spell out the relative Theorem~A axiom, since that is the remaining condition needed here. Suppose $A\xrightarrow{u}B\to C$ is a triangle and every induced map $A/c\to B/c$ is an $R$-homology equivalence. For the composite $f:A\to C$, form the bisimplicial set with $(p,q)$-simplices
\[
 \bigl(a_0\to\cdots\to a_p,\quad
       f(a_p)\to c_0\to\cdots\to c_q\bigr).
\]
Deleting the last $a$-vertex composes the crossing arrow with the image of that last arrow under $f$; deleting the first $c$-vertex composes on its other side. These rules, with the ordinary remaining faces and identity degeneracies, make the two simplicial directions commute. Take free $R$-modules, normalize in each direction, and use the usual sign in the total differential.

For each fixed $a$-chain, the second direction is the nerve of $f(a_p)\backslash C$, a category with initial object $\id_{f(a_p)}$. Its positive homology vanishes and its degree-zero homology is $R$. The spectral sequence that first takes this homology therefore identifies the homology of the total complex with $H_*(NA;R)$. In the other direction, for each fixed chain $c_0\to\cdots\to c_q$, the first simplicial direction is $N(A/c_0)$. Consequently the map induced by $u$ gives an isomorphism on that first homology page by the hypothesis on the comma categories. The comparison of the first-quadrant spectral sequences gives an isomorphism of total homology, and thus $H_*(NA;R)\cong H_*(NB;R)$. Convergence is sufficient here without any finiteness assumption: there are only finitely many bidegrees in any given total degree. This proves the relative localizer axiom. Finally the normalized-chain argument after ms.~79 proves W(7) for precisely this class $W_R$.

\emph{The counterexample.} Use the finite functor $f:\Psi\to C$ computed in the preceding block, with $\Psi=(z\to a,z\to b)$. It is ordinary final, so $f_!e_\Psi=e_C$. But $BC\simeq S^2$, and
\[
 H_2(BC;R)=R\ne0.
\]
Hence $\Psi$ is not totally $W_R$-aspherical. Its two full sieves $\{z\to a\}$ and $\{z\to b\}$, and their intersection $\{z\}$, each have a terminal object. Proposition~9(a), whose proof does not use W(7), makes all three totally $W_R$-aspherical. The constant functors from the two arrows to $\{z\}$ are retractions of the inclusions of the intersection. Thus both the union assertion and its proposed retraction repair fail, and the conclusion asserting totality of the outgoing fork fails as well. All this occurs in a fundamental localizer satisfying W(7), so the conditional qualifications in those statements do not save them.\hfill$\square$

The same calculation keeps two different defects apart. The proof's use of left Kan extension to preserve a monomorphism is invalid before any weak equivalence class is chosen. The example above additionally shows that the resulting total-asphericity conclusion can be false after all its stated homotopical axioms have been imposed.
\end{dcomment}


et telle que de plus

\begin{enumerate}[label=\alph*), start=2]
\item $X$ est $W$-asphérique, et $Y$ non.
\end{enumerate}

% a vertical bar is printed in the left margin beside the next paragraph (lines "La difficulté ..." to "... le foncteur")
La difficulté est la suivante, quand on cherche des contre-exemples `idiots' : c'est que l'hypothèse a) implique déjà, quand $X$ est $W_\infty$-asphérique (ou seulement 1-connexe) que $Y$ est 1-connexe : le foncteur
\[
Y' \longmapsto X' = X \times_Y Y'
\]
de la catégorie des revêtements étales $Y'$ de $Y$ vers celle des revêtements étales $X'$ de $X$ est pleinement fidèle (vérification immédiate), d'où le résultat annoncé. Donc un contre-exemple doit faire appel à l'existence d'un $\pi_i(Y,y)$ non nul, avec $i \geq 2$. Cela implique
\msp{page 102}
qu'il n'y a pas de contre-exemple avec des catégories de la forme $Y = Z/F$ (déduites d'un foncteur $q : X \longrightarrow Z$ en prenant $F = q_!(e_X) = \pi_0(X/Z)$), où $Z$ est une catégorie telle que $\bDelta^1$ (?), $\Phi$, $\Psi$, ou $\begin{tikzcd}[ampersand replacement=\&, sep=small, baseline=(M.base)]
 \& {} \& \\
{} \arrow[ur] \arrow[dr] \& |[alias=M]| \phantom{x} \& {} \arrow[ul] \arrow[dl] \\
 \& {} \&
\end{tikzcd}$ (ayant le type d'homotopie du cercle) etc., car les composantes connexes de telles catégories $Y$ sont des $K(\pi,1)$ - les $\pi_i$ supérieures ($i \geq 2$) sont nuls.

\dgrblock{dgrb531}{Why the first obstruction occurs above the fundamental group}
\begin{dcomment}{exposition}{p101}{Why the first obstruction occurs above the fundamental group}
Grothendieck's assertion about coverings explains why a counterexample with a circle as target cannot work. Let $f:A\to B$ be final, so every $b\downarrow f$ is nonempty and connected. Let $L,M:B\to\mathbf{Set}$ be local systems, meaning that every arrow is sent to a bijection. Given a natural transformation $\theta:Lf\to Mf$, choose any object $(a,\beta:b\to f(a))$ of $b\downarrow f$ and define
\[
\bar\theta_b=M(\beta)^{-1}\theta_a L(\beta).
\]
If $h:a\to a'$ is a morphism in that comma category, then $f(h)\beta=\beta'$ and the naturality of $\theta$ identifies the expressions obtained from $(a,\beta)$ and $(a',\beta')$. Connectedness makes the expression independent of all choices. For $\gamma:b'\to b$, the object $(a,\beta\gamma)$ computes $\bar\theta_{b'}$, proving naturality in $b$. Its restriction to $A$ is $\theta$, using $(a,\mathrm{id})$, and the same formula proves uniqueness. Thus restriction along $f$ is fully faithful on local systems.

Finality also preserves the set of connected components: it preserves the colimit of the constant singleton diagram, which is $\pi_0$. Suppose now that $BA$ is simply connected. Then $BB$ is connected, and every pulled-back local system on $A$ is constant. Fix a fibre of a local system $L$ on $B$. Each element of that fibre determines a section of the pulled-back system; full faithfulness, applied to maps from the singleton local system, makes it descend uniquely to a global section of $L$. Every monodromy action is therefore trivial. Applying this to the regular action of $\pi_1(BB)$ shows that the group itself is trivial. This supplies the omitted verification on p.~101.

The four-object example above is consistent with this restriction: $BC\simeq S^2$ is simply connected but not contractible. Its critical comma category is a circle, and taking $\pi_0$ discards exactly that loop before rebuilding the total category. By contrast, if the target category is a discrete fibration over one of the height-one categories drawn on p.~102, it has no composable nonidentity arrows. Its nerve is a graph. Each connected component is a $K(\pi,1)$, so simple connectivity would force it to be contractible. The source's exclusion of these examples is thus correct for the displayed bases; it is not a statement about arbitrary categories merely sharing their homotopy type.

In the modern cofinality theorem, connectedness of the comma categories is strengthened to their weak contractibility \citep[Theorem~4.1.3.1]{Lurie2009HTT}. That stronger condition controls all homotopy colimits and implies weak equivalence of classifying spaces. The ordinary finality used here controls ordinary colimits and already constrains coverings, but it leaves the higher homotopy obstruction exposed.
\end{dcomment}


Cette difficulté vue, le contre-exemple devient évident. Je pars d'un objet 1-connexe \textit{quelconque} $Y$ de $X$ \texttt{[plutôt de} Cat\texttt{]}, et je prends $y_0 \in \Ob Y$,
\[
X = \uCh(Y; y_0, -) \xrightarrow{\ f\ } Y.
\]
On sait que $X$ est $W_\infty$-asphérique, et que $f$ est $W_\infty$-parfait, donc les inclusions
\[
X_y \longrightarrow y\backslash X
\]
sont des $W_\infty$-équivalences. L'hypothèse a) que les $y\backslash X$ soient 0-connexes équivaut donc à la même hypothèse sur les fibres $X_y$, ou
\msp{page 103}
encore sur une seule fibre ($f$ étant parfait, donc une $W_\infty$-fibration, et $Y$ 0-connexe), soit
\[
X_{y_0} \simeq \Omega(Y, y_0).
\]
Sa 0-connexité équivaut à l'hypothèse de 1-connexité de $Y$ (on n'a utilisé encore que celle de 0-connexité). Donc pour un contre-exemple, il suffit de prendre $Y$ 1-connexe et non $\infty$-connexe. Le cas le plus simple est celui où $Y$ a le type d'homotopie de la sphère $S^2$, réalisé p. ex. par la catégorie ordonnée ci-contre \texttt{[ci-dessous]} à 6 sommets.
% foot: Version 7 décembre 2002

\printedpage{72}
% head: A. Grothendieck : Dérivateurs, Ch. XII, Édition des Universités de Montpellier II et Paris VII
\[
\begin{tikzcd}[row sep=small]
e_2^+ \bullet & \bullet\, e_2^- \\
e_1^+ \bullet \arrow[ur] & \bullet\, e_1^- \arrow[ul] \\
e_0^+ \bullet \arrow[ur] & \bullet\, e_0^- \arrow[ul]
\end{tikzcd}
\]
\texttt{[de plus,} $e_0^+ \leq e_1^+ \leq e_2^+$\texttt{,} $e_0^- \leq e_1^- \leq e_2^-$\texttt{]} (\footnote{C'est en somme cet exemple (à homotopie topologique près) que je réalise (sans d'abord m'en apercevoir dans l'Appendice p. 10-12, mais avec des catégories finies).}). Je résume le résultat obtenu.

\textbf{Proposition 10.} \textit{Soit $Y$ une petite catégorie} 1\textit{-connexe ($y_0 \in \Ob Y$). Alors} $\uCh(Y; y_0, -)$ \textit{est $W_\infty$-asphérique, mais n'est pas totalement $W_\infty$-asphérique si $Y$ elle-même n'est $\infty$-connexe, i.e. $W_\infty$-asphérique.} (\footnote{Mais un exemple plus percutant est $\Psi = \begin{tikzcd}[ampersand replacement=\&, row sep=small, column sep=small, baseline=(A.base)]
 \& b \\
|[alias=A]| a \arrow[ur] \arrow[dr] \& \\
 \& c
\end{tikzcd}$ , qui est $W_\infty$-asphérique, mais pas totalement!}).

On peut se demander si $Y$ $W_\infty$-asphérique implique que $\uCh(Y; y_0, -)$ est totalement $W_\infty$-asphérique. J'en doute fort!

\msp{page 104}
Revenons maintenant à la situation du début de la section, diagramme $(*)$, p.~87,
\[
\begin{tikzcd}
U \arrow[r, hook, "i"] \arrow[d, "p"'] & X \arrow[d, "q"'] & T = \{t\} \arrow[l, hook', dotted, "j"'] \arrow[d, dotted, "\mathrm{iso}"] \\
U' \arrow[r, hook, "i'"] & X' & T' = \{t'\}. \arrow[l, hook', dotted, "j'"']
\end{tikzcd}
\tag*{$(*)$}
\]
Nous pouvons donner maintenant un exemple où $i$ est dans $W$ (et même $W$-asphérique) sans que $i'$ le soit :
\[
i \in W, \qquad i' \notin W. \tag*{$(**)$}
\]
On peut même supposer que $t$ soit objet final de $X$ (ce qui signifie que l'objet $H$ de $U^\wedge$ est isomorphe à $e_U$, l'objet final), et $t'$ objet final de $X'$ (i.e. $p_!(e_U) = e_{U'}$, i.e. les $u'\backslash U$ sont 0-asphériques). Dans ce cas, comme $X$, $X'$ sont $W$-asphériques, les relations $(**)$ équivalent à
\[
U \text{ est } W\text{-asphérique, et } U' \text{ non,} \tag*{$(***)$}
\]
ce qui est justement la situation que nous venons de construire (p. 101-103).

\dgrblock{dgrb532}{The information lost by taking components}
\begin{dcomment}{exposition}{p103}{The information lost by taking components}
Grothendieck has found the obstruction that the preceding attempts were missing: an ordinary direct image of a set-valued presheaf remembers connected components of comma categories, whereas a weak equivalence of nerves must preserve all homotopy groups. The path construction makes the difference particularly visible. Put $P=\underline{\mathrm{Ch}}(Y;y_0,-)$ and let $f:P\to Y$ be its endpoint functor. The earlier path theory supplies the two facts used here: $BP$ is contractible and the homotopy fibre of $Bf$ over $y$ is a based path space from $y_0$ to $y$. We write $BA=|NA|$ for the classifying space of a category. The comparison of strict fibres with comma categories is the additional consequence of the perfectness hypothesis; contractibility of $BP$ alone would not supply that comparison.

For connected $Y$, the components of each such path space form a torsor for $\pi_1(BY,y_0)$. Thus simple connectivity gives
\[
 (f_!e_P)(y)=\pi_0(y\backslash P)=\{*\}.
\]
Consequently $f_!e_P=e_Y$, and the category of elements of this direct image is $Y$ itself. If $BY$ is not contractible, this single functor $f$ disproves total $W_\infty$-asphericity of $P$. This argument uses precisely the universal quantifier in Definition~1: the homotopy type of $P$ cannot determine how all its ordinary direct images behave.

The six-element ordered set displayed in the text is a useful finite model of the required target. Its elements occur in three levels $\{e_k^+,e_k^-\}$, $k=0,1,2$, and every element of a lower level precedes every element of a higher level. Its order complex is therefore
\[
 S^0*S^0*S^0\cong S^2.
\]
There are six vertices, twelve edges, and eight triangular faces. Each simplex chooses at most one vertex in each level, so these are exactly the faces of the boundary of an octahedron. The two chains supplied in the editors' brackets are essential: the diagonal arrows alone would not specify this ordered set.

The finite fork construction in the commentary at ms.~100 supplies the same failure without a path category: a contractible $\Psi$ admits a functor $\Psi\to C$ with connected undercategories and $BC\simeq S^2$. We shall use it below. The question about $\underline{\mathrm{Ch}}(Y;y_0,-)$ when $BY$ is contractible is a different question. Proposition~10 gives no implication in that case, and a replacement of the path category by some equivalent contractible category would not answer it. The next commentary answers it, negatively, with this same fork.
\end{dcomment}


\dgrblock{dgrb533}{The path category over a contractible base need not be total}
\begin{dcomment}{development}{p103}{The path category over a contractible base need not be total}
The question that closes this page, whether $\underline{\mathrm{Ch}}(Y;y_0,-)$ is totally $W_\infty$-aspherical when $Y$ is $W_\infty$-aspherical, has a negative answer, and the witness is the category that Grothendieck names in the footnote to Proposition~10.

\textbf{Proposition.}\ours
\textit{Let $Y=\Psi$ be the outgoing fork $a\leftarrow z\to b$ and $y_0=z$. Then $BY$ is contractible, and $P=\underline{\mathrm{Ch}}(Y;z,-)$ is $W_\infty$-aspherical and not totally $W_\infty$-aspherical. It is not totally $W_R$-aspherical for any nonzero commutative ring $R$ either.}

\textit{Proof.}
The object $z$ is initial, so $BY$ is contractible; in particular $Y$ is $1$-connected. Proposition~10 applies: $P$ is $W_\infty$-aspherical, and, as recalled in the preceding commentary, the endpoint functor $f:P\to Y$ satisfies $f_!e_P=e_Y$, that is, every $y\backslash P$ is nonempty and connected. So $f$ is final. The inclusion $g:\Psi\to C$ of the commentary at ms.~100 is final, and $BC\simeq S^2$. Final functors compose, because $(gf)_!e_P=g_!f_!e_P=g_!e_Y=e_C$. The category of elements of $(gf)_!e_P$ is $C$, which is neither $W_\infty$-aspherical nor $W_R$-aspherical.\hfill$\square$

The proof uses only two facts: the endpoint functor of the path category over a $1$-connected base is final, and total asphericity passes to final images. Hence $\underline{\mathrm{Ch}}(Y;y_0,-)$ is totally $W_\infty$-aspherical only if $Y$ is. Since total asphericity of $Y$ is far stronger than asphericity (for a category with finitely many objects it means that $Y$ is filtered; see the commentary on Chapter~II, ms.~36), the path construction does not produce totally aspherical categories from aspherical ones.
\end{dcomment}


On aimerait aussi trouver un exemple où dans $(*)$ on ait $p \in W$, $q \notin W$.
% foot: Version 7 décembre 2002

% ---- from Der-ChXIIb-012_p031-035.tex
\setcounter{footnote}{76}

\printedpage{73}
% head: A. Grothendieck : Dérivateurs, Ch. XII, Édition des Universités de Montpellier II et Paris VII
% foot: Version 7 décembre 2002

\msp{page 105}

Je prendrai encore $(X,U,\{t\})$ tel que $t$ soit objet final de $X$, mais si on veut $q\notin W$, on ne peut supposer $t'$ élément final de $X'$. Au lieu de cela, je vais supposer

\begin{enumerate}[label=\alph*)]
\item $U$, $U'$ (et aussi $X$, ayant objet final $t$) $W$-asphériques, donc $i\in W$ (et même $i$ $W$-asphérique), $p\in W$, mais
\item $X'$ non $W$-asphérique, donc $i',q\notin W$.
\end{enumerate}

Pour ceci, notons que
\[ U'/t'\simeq U'/p_!(e_U)\overset{\text{déf}}{=}\underline{\pi}_0(U/U'). \]
Le corollaire de la proposition 7 (cf. p. 85), non encore démontré, nous dit que $i'\in W$ (i.e. ici, $X'$ $W$-asphérique) si et seulement si $i'$ est $W$-asphérique, i.e. $U'/t'$ est $W$-asphérique. Nous allons l'admettre ici. Tout revient donc à trouver une flèche
\[ p:U\longrightarrow U' \]
dans Cat avec les propriétés suivantes.

\begin{enumerate}[label=\alph*)]
\item $U$ et $U'$ sont $W$-asphériques,
\item $\underline{\pi}_0(U/U')$ n'est pas $W$-asphérique.
\end{enumerate}

Nous allons reprendre l'exemple précédent
\[ f:X\longrightarrow Y \]
(p. 101-103), en prenant
\[ U=X,\qquad \underline{\pi}_0(U/U')=Y, \]

\msp{page 106}

donc tout revient, à partir de l'exemple donné $X\xrightarrow{f}Y$, à construire une flèche
\[ \pi:Y\longrightarrow U' \tag*{$(*)$} \]
dans Cat qui satisfasse les conditions

\begin{enumerate}[label=\alph*)]
\item $\pi$ est un isomorphisme local,
\item $U'$ est $W$-asphérique.
\end{enumerate}

On prendra alors
\[ p=\pi f. \]
Mais par la façon dont $X$ a été construit à partir d'un $Y$ 1-connexe et non $\infty$-connexe, par ailleurs quelconque, on est ramené finalement à trouver $(*)$ satisfaisant a), b), et de plus

\begin{enumerate}[label=\alph*), start=3]
\item \printedpage{74}
% head: A. Grothendieck : Dérivateurs, Ch. XII, Édition des Universités de Montpellier II et Paris VII
% foot: Version 7 décembre 2002
$Y$ est 1-connexe, mais non $\infty$-connexe.
\end{enumerate}

(NB On suppose maintenant $W=W_\infty$.) Il revient au même de trouver

% printed with a vertical rule at the left of this block
\begin{quote}
$U'$ objet \textit{asphérique} de Cat, $F\in\operatorname{Ob}U'^\wedge$ tels que $U'/F$ soit 1-connexe et non $\infty$-connexe.
\end{quote}

On songe par exemple à prendre simplement une sous-catégorie ouverte $U'_0$ de $U'$, qui soit 1-connexe et non $\infty$-connexe. Partons d'une telle catégorie $U'_0$ - par exemple une réalisation homotopique dans Cat, de la sphère $S^2$. On l'immerge comme ouvert dans un $U'$ asphérique, en lui ajoutant un objet final. On a gagné.

\msp{page 107}

Faisons un énoncé en forme.

\textbf{Proposition 11.} \textit{On peut trouver une situation} $(*)$ (p. 124 \texttt{[plutôt 104]}) \textit{avec}

\begin{enumerate}[label=\textit{\alph*)}]
\item \textit{$t$ objet final de $X$.}
\item \textit{$U$, $U'$ (et $X$) $W_\infty$-asphériques, donc $i,p\in W_\infty$ (et même $i$ $W_\infty$-asphérique).}
\item \textit{$X'$ non $W_\infty$-asphérique, donc $i',q\notin W_\infty$.}
\end{enumerate}

Il suffit pour cela de partir d'un objet 1-connexe $Y$ de Cat, de lui adjoindre un objet final $u'_0$, ce qui donne $U'$, de prendre $U=\underline{\mathrm{Ch}}(Y;y_0,-)$, $p:U\longrightarrow U'$ le composé $U\xrightarrow{\text{but}}Y\underset{\text{can.}}{\overset{\text{incl.}}{\hra}}U'$, et $X$ déduit de $U$ en lui ajoutant un objet final $t$. On a alors $U'/t'\simeq Y$, d'où
\[ X'\ W_\infty\text{-asphérique (i.e. }i'\in W_\infty)\Longleftrightarrow Y\ W_\infty\text{-asphérique.} \]
On a le contre-exemple cherché en prenant $Y$ (1-connexe mais) non $\infty$-connexe.

\dgrblock{dgrb534}{Two cones and the two failures of pushout invariance}
\begin{dcomment}{examples}{p107:On a le contre-exemple cherché en prenant}{Two cones and the two failures of pushout invariance}
Grothendieck now converts the failure of total asphericity into failures of ordinary categorical pushouts. The two constructions deserve to be separated because they test different edges of the square. We can make both finite by using the functor $f:\Psi\to C$ constructed at ms.~100. Here $\Psi=(z\to a,z\to b)$ has an initial object, $BC\simeq S^2$, and every category $c\backslash\Psi$ is connected and nonempty. Hence
\[
 f_!e_\Psi=e_C.
\]
The crucial point is that this is an equality of ordinary presheaves; $f$ is not a weak equivalence.

For the first square take $U=\Psi$, $U'=C$, and let $X=\Psi^\triangleright$ be obtained by adjoining a new terminal object $t$. The gluing weight for $U\subset X$ is $e_\Psi$. The pushout formula proved in \S1 says that $X'=X\amalg_\Psi C$ has gluing weight $f_!e_\Psi=e_C$. Thus $X'=C^\triangleright$, with its new terminal object $t'$. Both $X$ and $X'$ are contractible. The inclusion $i:\Psi\to\Psi^\triangleright$ is a weak equivalence and is even $W_\infty$-aspherical: its slices over old objects have terminal objects and its slice over $t$ is $\Psi$. But $i':C\to C^\triangleright$ is not a weak equivalence, because its source has the homotopy type of $S^2$. This proves the assertion at ms.~104 without any converse to the gluing lemma.

For the second square first adjoin a terminal object $s$ to $C$, and write $U'=C^\triangleright$. Let $j:C\hookrightarrow U'$ be the resulting sieve and put $p=jf$. Now $U$, $U'$ and $X=U^\triangleright$ are all contractible, so both $i$ and $p$ are weak equivalences. Since $j$ is a sieve, extension by $j_!$ sends $e_C$ to the subterminal presheaf which is a singleton on $C$ and empty at $s$. Indeed, $c\backslash C$ has an initial object for $c\in C$, whereas $s\backslash C$ is empty. Therefore
\[
 p_!e_U=j_!e_C,
\qquad
 (p_!e_U)(c)=\{*\},\quad (p_!e_U)(s)=\varnothing.
\]
The pushout $X'=X\amalg_U U'$ is consequently obtained from $C$ by adjoining two distinct objects $s,t'$ above every object of $C$, with no arrow between $s$ and $t'$. It is the union of two terminal cones on $C$.

This description computes the homotopy type directly. Every composable string in $X'$ lies in one of the two cones, because a string cannot contain both of their incomparable apexes. Hence
\[
 NX'=N(C^\triangleright_s)
       \amalg_{NC}N(C^\triangleright_{t'}).
\]
The two inclusions of $NC$ are simplicial monomorphisms, so this ordinary union is a homotopy pushout. Its realization is the unreduced suspension of $BC$, and therefore
\[
 BX'\simeq\Sigma BC\simeq S^3.
\]
Neither $i':U'\to X'$ nor $q:X\to X'$ is a weak equivalence. This proves Proposition~11 at ms.~107 and shows exactly how an arbitrary sieve can destroy pushout invariance even when the other three categories are contractible.

The text invokes the converse in Proposition~7 to justify the second construction. That converse depends on the stronger gluing axiom discussed at ms.~78--85 and cannot be used as an unconditional theorem for $W_\infty$. The suspension computation repairs this use of it. More generally, for the simply connected $Y$ of the manuscript, $\Sigma BY$ is contractible exactly when $BY$ is contractible: suspension shifts reduced integral homology, and a simply connected acyclic CW complex is contractible by the Hurewicz and Whitehead theorems. The hypothesis of simple connectivity is doing real work in this last equivalence.
\end{dcomment}


\subsection*{2.4\quad Application aux immersions ouvertes dwyériennes.}

Donnons quand même l'enoncé \textit{positif} auquel nous en avions. % [sic] l'enoncé -> l'énoncé, missing accent

\textbf{Théorème 3.} \textit{Soient $(X,U,T)$ un triple $\in\mathcal{T}$ ($U$ ouvert de $X$ dans} Cat\textit{, $T=X\backslash U$), et pour toute flèche $p:U\longrightarrow U'$ flèche de} Cat\textit{, le carré cocartésien dans} Cat

\[
\begin{tikzcd}
U \arrow[r, hook, "i"] \arrow[d, "p"'] & X \arrow[d, "q"'] & T \arrow[l, dotted, hook'] \arrow[d, dotted, "\text{iso}"] \\
U' \arrow[r, hook, "i'"] & X' & T'. \arrow[l, dotted, hook']
\end{tikzcd}
\]

\printedpage{75}
% head: A. Grothendieck : Dérivateurs, Ch. XII, Édition des Universités de Montpellier II et Paris VII
% foot: Version 7 décembre 2002

\msp{page 108}

\textit{Supposons que pour tout $t\in\dot{T}=\dot{U}$, le préfaisceau d'ensembles}
\[ i^*(t)=(u\longmapsto\operatorname{Hom}(u,t))\in U^\wedge \]
\textit{soit totalement $W$-asphérique (déf. 1, p. 88)} (\footnote{\textbf{NB} On dit sous ces conditions que $i_0:U\longrightarrow\bar{U}$ est \textit{strictement $W$-asphérique} (cf. déf. 2 page suivante).}) \textit{- p. ex. que la catégorie $U/t$ soit totalement $W$-asphérique (prop. 8, page 93). (NB le cas où les $U/t$ ont un objet final, cf. prop. 9 a), i.e. celui où l'inclusion $U\longrightarrow\bar{U}$ admet un adjoint à droite, est celui envisagé par Dwyer, celui des immersions ouvertes de Dwyer} (\footnote{Je ne connais pas d'autre cas où th. 3 s'applique - faute de connaître de catégorie totalement $W$-asphérique qui n'ait un objet final, ni d'objet $F$ d'une catégorie $U^\wedge$ qui soit totalement $W$-asphérique sans être représentable, i.e. $U/F$ a un objet final. À présent, je sais que les catégories \textit{filtrantes} sont totalement $W$-asphériques si $W$ satisfait l'axiome des limites inductives.})\textit{.) Alors pour toute $p:U\longrightarrow U'$, l'immersion ouverte}
\[ i'_0:U'\longrightarrow\overline{U'} \]
\textit{est $W$-asphérique (et même, tout comme $i_0$, strictement $W$-asphérique), et a fortiori elle est} \texttt{[dans]} \textit{$W$. Par suite (prop. 6} \texttt{[p. 81]}\textit{) on a ceci.}

\[
\begin{aligned}
&\text{\textcircled{a}}\qquad && i\in W && \Longleftrightarrow\quad && i'\in W\quad && \textit{moyennant } \mathrm{W}(7\text{ bis}).\\
&\text{\textcircled{b}}\qquad && p\in W && \Longleftrightarrow\quad && q\in W\quad && \textit{moyennant } \mathrm{W}(7).\ (\footnotemark)
\end{aligned}
\]
\footnotetext{Démonstration page 112.}

\dgrblock{dgrb535}{A condition on weights that survives every direct image}
\begin{dcomment}{exposition}{p108:Démonstration page 112.}{A condition on weights that survives every direct image}
After the counterexamples, Grothendieck seeks a hypothesis expressed in the actual gluing data. For a sieve $i:U\hookrightarrow X$, write $\overline U$ for the cosieve generated by $U$, and put $\dot T=\overline U\setminus U$. An object $t$ belongs to $\dot T$ precisely when at least one arrow $u\to t$ exists with $u\in U$. Its gluing weight is
\[
 H_t=\Hom_X(i(-),t)\in\widehat U.
\]
Only these objects need testing. At $u\in U$ the weight is representable and therefore totally aspherical; outside $\overline U$ it is empty and no assertion of asphericity is intended.

The hypothesis of Theorem~3 is that every $H_t$, $t\in\dot T$, remains aspherical after every ordinary left Kan extension. A pushout along $p:U\to U'$ replaces this weight by $p_!H_t$. Its total asphericity follows without a gluing axiom: for a second functor $r:U'\to V$,
\[
 r_!(p_!H_t)\cong(rp)_!H_t,
\]
and the right-hand side is aspherical by the original hypothesis. This calculation explains both the terminology ``strictly aspherical'' for $U\to\overline U$ and its stability under arbitrary pushout.

The support is preserved as well. If $H_t$ is nonempty, a representative $h\in H_t(u)$ produces a class at $p(u)$ in $p_!H_t$; a colimit of sets cannot delete that class. If $H_t$ is empty, its direct image is empty. Thus the boundary objects of the two pushouts identify with the same $\dot T$, and the preceding test covers exactly the new boundary. It follows that
\[
 i'_0:U'\longrightarrow\overline{U'}
\]
is strictly $W$-aspherical and hence $W$-aspherical. The claim at the start of the proof on ms.~112 that $i'$ itself is aspherical must read $i'_0$. A $W$-dwyérienne inclusion need not be a weak equivalence: the empty sieve into a nonempty category already demonstrates the distinction.

The two equivalences printed as (a) and (b) have a different logical status from this weight calculation. They use respectively W(7 bis) and W(7), exactly as stated. Those additional axioms must remain visible; they are not automatic properties of weak homotopy equivalences of nerves. The modern Dwyer pushout theorem gives the useful forward implications in its own hypotheses, as discussed in the later-development block. The intrinsic statement here, valid before imposing either gluing axiom, is preservation of the condition on all the weights.
\end{dcomment}


\dgrblock{dgrb536}{A boundary with no last object}
\begin{dcomment}{examples}{p108:Démonstration page 112.}{A boundary with no last object}
The footnote about filtered categories points beyond the original adjoint criterion. The commentary on Chapter~II, manuscript pp.~35--36 (``Left Kan extension carries the ind-presentation'' and ``The quantifier in total asphericity''), explains why filtered presentations make this universal condition tractable. Here their role can be seen in one concrete inclusion.

Let $U$ be the ordered set $\mathbb N$ with its usual order, and let $X=\mathbb N\cup\{\infty\}$ with $n<\infty$ for every $n$. The inclusion $i:U\hookrightarrow X$ is a sieve and $\overline U=X$. For an old object $n$, the slice $U/n$ has terminal object $n$. For the new object $\infty$, the slice is $U$ itself. For any $p:\mathbb N\to B$, one has
\[
 F:=p_!e_{\mathbb N}\cong\varinjlim_n\Hom_B(-,p(n)).
\]
Write $p_{nm}:p(n)\to p(m)$ for the transition map when $n\leq m$, and let $E=B/F$. This category is nonempty, since it contains $(p(0),[\id_{p(0)}])$. Given finitely many objects $(b_j,x_j)$ of $E$, choose representatives $g_j:b_j\to p(n_j)$ of their elements. For $m\geq n_j$ for every $j$, the arrows $p_{n_jm}g_j$ give morphisms from all these objects to $(p(m),[\id_{p(m)}])$.

For parallel morphisms $h,k:(b,x)\rightrightarrows(c,y)$ in $E$, choose a representative $g:c\to p(n)$ of $y$. The defining condition on these morphisms gives $[gh]=x=[gk]$ in $F(b)$. Equality in this filtered colimit means that, for some $m\geq n$,
\[
 p_{nm}gh=p_{nm}gk.
\]
Thus $p_{nm}g:(c,y)\to(p(m),[\id_{p(m)}])$ equalizes $h$ and $k$. These are exactly the nonemptiness, finite-cocone, and parallel-arrow tests for a filtered category. A filtered category has weakly contractible nerve: a cocone on the finite data of a simplicial sphere map contracts that map. Thus every such $E$ is $W_\infty$-aspherical.

It follows that $i$ is totally $W_\infty$-dwyérienne. It is also a couniversal weak equivalence, without using either additional gluing axiom. After any pushout its slice over the new object is the filtered category of elements just computed, while its slices over the old objects have terminal objects. Quillen's Theorem~A therefore applies to every pushout inclusion. It is not a Dwyer immersion in the original adjoint sense. A right adjoint $r:X\to U$ to $i$ would require
\[
 n\leq r(\infty)\quad\text{for every }n\in\mathbb N,
\]
which is impossible. This distinguishes the two classes even though both $\mathbb N$ and $\mathbb N^\triangleright$ are contractible.

The negative answer also holds within the full gluing framework assumed in the source's question. Let $W_{\mathbb Z}$ denote integral-homology equivalences of nerves. Every pushout just considered is a weak homotopy equivalence and therefore an integral-homology equivalence, so $i$ is couniversally a $W_{\mathbb Z}$-equivalence. Every tested category of elements is already weakly contractible, so $i$ is also totally $W_{\mathbb Z}$-dwyérienne. The development at ms.~100, \emph{Testing the retracted assertions under their stated axiom}, verifies that $W_{\mathbb Z}$ is a fundamental localizer satisfying full W(7). The failure of the right-adjoint condition is independent of the choice of weak equivalences. Thus the proposed recovery of an original Dwyer immersion fails even with W(7) imposed. The distinct converse from an arbitrary homotopical cofibration to Grothendieck's condition on every boundary weight is not established by this example.
\end{dcomment}


\textbf{Corollaire 1.} \textit{Supposons} W(7)\textit{. Si $i:U\hra X$ est une immersion ouverte telle que $i_0:U\hra\bar{U}$ soit totalement} \texttt{[il suffit strictement]} \textit{$W$-asphérique, alors $i:U\longrightarrow X$ est une $W$-cofibration (`ultraforte'), par quoi j'entends ceci : pour toute flèche $U'\xrightarrow{f}U''$ dans $U\backslash$}Cat\textit{,}

\msp{page 109}

\textit{telle que $f\in W$, la flèche $g:X'\longrightarrow X''$ qui s'en déduit par cochangement de base à partir de l'objet $X$ de $U\backslash$}Cat\textit{, est aussi dans $W$.}

\textsc{Démonstration du corollaire} : on applique le théorème 3 à l'immersion ouverte $U'\overset{i'}{\hra}X'$ déduite de $U\overset{i}{\hra}X$ par le cochangement de base $U\xrightarrow{f}U'$, en notant de plus que l'hypothèse du théorème 3 sur $i$ est préservée en passant de $i$ à $i'$.

Il faut quand même introduire une terminologie adéquate :

\textbf{Définition 2 :}

\begin{enumerate}[label=\alph*)]
\item Une flèche $f:X\longrightarrow Y$ est dite \textit{strictement $W$-asphérique} si pour toute $y\in\operatorname{Ob}Y$, $f^*(y)=(x\longmapsto\operatorname{Hom}_Y(f(x),y))$ est un objet totalement $W$-asphérique de $X^\wedge$ (cf. déf. 1, page 88). Elle est dite \textit{totalement $W$-asphérique} si pour tout $y\in\operatorname{Ob}Y$, $X/y$ ($\simeq X/f^*(y)$) est un objet totalement $W$-asphérique de Cat.

\item \printedpage{76}
% head: A. Grothendieck : Dérivateurs, Ch. XII, Édition des Universités de Montpellier II et Paris VII
% foot: Version 7 décembre 2002
Une immersion ouverte $i:U\longrightarrow X$ est dite \textit{$W$-dwyérienne} (resp. \textit{totalement} \texttt{[$W$-]} \textit{dwyérienne}) si l'immersion induite $i_0:U\longrightarrow\bar{U}$ est strictement (resp. totalement) $W$-asphérique. Une immersion fermée $j:T\longrightarrow X$ est dite \textit{$W$-dwyérienne} (resp. \textit{totalement $W$-dwyérienne}) si et seulement si $j^\circ:T^\circ\longrightarrow X^\circ$ est une immersion ouverte $W$-dwyérienne (resp totalement $W$-dwyérienne).

\item Une flèche $i:U\longrightarrow X$ de Cat (pas nécessairement une immersion)

\msp{page 110}

est appelée une \textit{$W$-cofibration,} si pour toute flèche $f:U'\longrightarrow U''$ dans $U\backslash$Cat, la flèche $g:X'\longrightarrow X''$ qui s'en déduit par le cochangement de base par $X$ (objet de $U\backslash$Cat) est dans $W$ quand $f$ l'est.
\end{enumerate}

\textbf{NB} Cette notion de $W$-cofibration est stable par composition et par cochangement de base, et les isomorphismes de Cat sont des $W$-cofibrations. Si $f$ est une $W$-cofibration, $f^\circ$ l'est et inversement.

(Le corollaire s'énonce en disant qu'une immersion ouverte ou fermée qui est $W$-dwyérienne est une $W$-cofibration.)

\begin{enumerate}[label=\alph*), start=4]
\item Une flèche $i:U\longrightarrow X$ dans Cat est appelée une \textit{$W$-cofibration triviale}, ou une \textit{$W$-équivalence couniverselle,} si pour toute $U'$ dans $U\backslash$Cat, i.e. toute flèche $p:U\longrightarrow U'$ dans Cat, la flèche déduite de $i$,
\[ i':U'\longrightarrow X'=U'\amalg_U X'\quad\texttt{[plutôt }U'\amalg_U X\texttt{]}\ \ , \]
est dans $W$. (NB Il est immédiat que cela implique bien que $i$ est une $W$-cofibration.)
\end{enumerate}

(Cette notion de $W$-cofibration triviale est elle aussi stable par composition et par cochangement de base, et les isomorphismes dans Cat en sont.)

\msp{page 111}

\dgrblock{dgrb537}{What the cofibration condition quantifies over}
\begin{dcomment}{exposition}{p110}{What the cofibration condition quantifies over}
Grothendieck's definition concerns the entire undercategory, and this is why he calls the condition ``ultraforte''. For $i:U\to X$, let
\[
 i_!:U\backslash\Cat\longrightarrow X\backslash\Cat,
 \qquad (U\to A)\longmapsto(X\to X\amalg_U A).
\]
The condition is that this functor preserve the weak equivalences of underlying categories. In particular it tests a weak equivalence $A\to B$ after an arbitrary preliminary change of base $U\to A$. Merely requiring preservation of weak equivalences whose source is $U$ would be weaker. The modern terminology and its relation to model-category cofibrations are treated in the later-development block.

The elementary closure assertions follow from the universal properties of pushouts. If $U\to X\to Z$ are two such maps, pushout along the composite is the composite of their pushout functors, up to the canonical isomorphism, so it preserves weak equivalences. If $U\to V$ is arbitrary and $i':V\to X\amalg_U V$ is the pushout of $i$, then for $A\to B$ under $V$ its two pushouts along $i'$ are canonically the pushouts along $i$ of the same map considered under $U$. This proves stability under change of base. Isomorphisms give equivalences of the undercategories. Passage to opposite categories proves the duality assertion provided that $W$ is invariant under opposites, as in the standing localizer framework.

The adjective ``trivial'' is defined by a separate universal property: every pushout $A\to X\amalg_U A$ of $i$ belongs to $W$. This implies the preceding preservation condition by two-out-of-three. Indeed the natural square
\[
\begin{tikzcd}
A\arrow[r]\arrow[d,"f"']&X\amalg_U A\arrow[d,"g"]\\
B\arrow[r]&X\amalg_U B
\end{tikzcd}
\]
has both horizontal maps in $W$; when $f\in W$, commutativity and two-out-of-three give $g\in W$. Stability of these universal weak equivalences under composition and pushout follows by factoring each new pushout into two pushouts of the original maps.

Without an additional hypothesis one should not replace this definition by ``a $W$-cofibration which belongs to $W$''. Corollary~2 establishes that implication for the $W$-dwyérienne sieves considered there, using Theorem~3 and the stated gluing assumptions. This is the reason for retaining the full definition even though its name resembles the terminology of a model category.
\end{dcomment}


\dgrblock{dgrb538}{The later theory of homotopy cofibrations}
\begin{dcomment}{later}{p110}{The later theory of homotopy cofibrations}
Grothendieck wants a property of an attachment that is visible from weak equivalences alone. This is exactly the property now called an \emph{$h$-cofibration}: the functor of cobase change on undercategories preserves weak equivalences. The terminology and a general model-categorical treatment appear in \citet[Definition~1.1]{BataninBerger2017}, who explicitly acknowledge the antecedent in Grothendieck's manuscripts. In a left proper model category their Proposition~1.6 and Corollary~1.9 identify this condition with the requirement that every pushout along the map compute a homotopy pushout. Their Lemma~1.10 also proves that an $h$-cofibration which is a weak equivalence is couniversal. These are precise later replacements for the useful direction of the present gluing programme; none requires the reverse implication of W(7 bis). The same undercategory definition is retained as ``$W$-cofibration'' by \citet[Definition~2.1 and Lemma~2.4]{ManettiMeazzini2019}.

The definition here must be kept distinct from the cofibrations of a chosen Quillen model structure. For example, for every small category $C$, the map $\varnothing\to C$ is a $W_\infty$-cofibration: its cobase-change functor is $A\mapsto A\amalg C$, and the nerve of a disjoint union is a disjoint union of nerves. A weak equivalence on the first summand and the identity on the second induce a weak equivalence on the union, component by component. Thus every category is cofibrant for Grothendieck's \emph{relative} notion. This statement does not assert that every category is cofibrant in Thomason's model structure. The distinction explains why a satisfactory description of these attachments can be broader than the particular generating cofibrations used to construct a model category.
\end{dcomment}


\textbf{Corollaire 2.} (Moyennant W(7).) \textit{Soit $i:U\longrightarrow X$ une immersion ouverte $W$-dwyérienne. Pour que $i$ soit une $W$-cofibration triviale, il faut et il suffit que $i\in W$, ou ce qui revient au même} (cf. page 82), \textit{que $\dot{T}\hra T$ soit dans $W$.}

C'est évidemment nécessaire, et la suffisance résulte du théorème 3 a).

\textbf{Questions ouvertes.}

\begin{enumerate}[label=\alph*)]
\item J'ignore si une immersion ouverte $i$ qui est une $W$-cofibration, voire une $W$-cofibration triviale, est forcément $W$-dwyérienne (\footnote{et peut-être même une immersion de Dwyer?}). Mais par le corollaire de la proposition 7 (p. 85), je sais que c'est le cas si on suppose $T$ isomorphe à la catégorie finale $e$. Il s'ensuit toute au moins ceci : Pour que non seulement $i:U\longrightarrow X$, mais aussi toutes les immersions ouvertes $i_1:U\longrightarrow X_1$ qui s'en déduisent par le procédé de prop. 7 b) au moyen de flèches quelconques $T_1\longrightarrow T$ de but $T$ dans Cat, soient des $W$-cofibrations, il faut et il suffit que $i$ soit dwyérienne. % [sic] toute au moins -> tout au moins, agreement error

\item \printedpage{77}
% head: A. Grothendieck : Dérivateurs, Ch. XII, Édition des Universités de Montpellier II et Paris VII
% foot: Version 7 décembre 2002
Y-a-t-il une bonne notion d' immersion dwyérienne, applicable à toutes les immersions

\msp{page 112}

\dgrblock{dgrb539}{The open questions of this page}
\begin{dcomment}{development}{p111}{The open questions of this page}
Question~a) asks whether a sieve that is a $W$-cofibration, or even a trivial $W$-cofibration, must be $W$-dwy\'erienne, and asserts this when the complement $T$ is the final category. Three small examples separate the notions. In each of them $U=\{u\}$ is a single object, so a functor $p:U\to A$ is the choice of an object $c=p(u)$ of $A$.

\textbf{Proposition (a trivial cofibration that is not dwy\'erienne).}\ours
\textit{Let $X$ have objects $u,t,s$ and nonidentity arrows $\alpha,\beta:u\to t$, $\gamma:t\to s$, $\delta:u\to s$, with $\gamma\alpha=\gamma\beta=\delta$. The sieve $U=\{u\}\hookrightarrow X$ is a trivial $W$-cofibration for every fundamental localizer $W$. It is not $W$-dwy\'erienne when $W$ is $W_\infty$, a homology localizer $W_R$ with $R\neq0$, or the class $W_n$ of $n$-equivalences for some $n\geq0$.}

\textit{Proof.}
The presheaf $v\mapsto\operatorname{Hom}_X(v,t)$ on $U$ has two elements, so $U/t$ is a discrete category with two objects; the test $p=\mathrm{id}$ shows that this presheaf is not totally $W$-aspherical for the localizers named. Let $p:U\to A$ be arbitrary, $c=p(u)$, and $P=X\amalg_UA$. Its objects are those of $A$ and $t,s$; for $v\in A$,
\[
 P(v,t)=A(v,c)\times\{\alpha,\beta\},\qquad P(v,s)=A(v,c),\qquad P(t,s)=\{\gamma\},
\]
composition with $\gamma$ forgets the second coordinate, and there are no arrows from $t$ or $s$ to $A$. Let $j:A\to P$ be the inclusion and $r:P\to A$ the functor that is the identity on $A$, sends $t,s$ to $c$, projects the two crossing sets to $A(v,c)$ and sends $\gamma$ to $\mathrm{id}_c$. Let $F:P\to P$ be the identity on $A$, send $t$ and $s$ to $s$, each crossing arrow to its composite into $s$, and $\gamma$ to $\mathrm{id}_s$. There are natural transformations
\[
 \mathrm{id}_P\Longrightarrow F\Longleftarrow jr,
\]
both the identity on $A$; the first has components $\gamma$ at $t$ and $\mathrm{id}_s$ at $s$, the second has at $t$ and at $s$ the arrow $c\to s$ corresponding to $\mathrm{id}_c$. Naturality at a crossing arrow $(g,\alpha)$ or $(g,\beta)$ is the relation $\gamma\alpha=\gamma\beta=\delta$. Since $rj=\mathrm{id}_A$, the map $Bj$ is a homotopy equivalence. So every cobase change of $U\to X$ is in $W_\infty\subseteq W$.\hfill$\square$

This answers the first sentence of question~a) negatively; the footnote's stronger hope fails with it, since a Dwyer immersion is dwy\'erienne.

\textbf{Proposition (the case $T\simeq e$).}\ours
\textit{Let $X_0$ have objects $u,t$ and exactly two nonidentity arrows $\alpha,\beta:u\to t$, and let $U=\{u\}$, $T=\{t\}$. For $W=W_\infty$ and for $W=W_R$, the sieve $U\hookrightarrow X_0$ and every sieve $U\hookrightarrow X_1$ deduced from it by a change of complement $T_1\to T$ (Proposition~7~b)) is a $W$-cofibration, and none of them with $T_1\neq\varnothing$ is $W$-dwy\'erienne.}

\textit{Proof.}
The category $X_1$ has the objects $u$ and those of $T_1$, the arrows of $T_1$, and two arrows $\alpha_x,\beta_x:u\to x$ for every $x\in T_1$, compatible with the arrows of $T_1$. Each crossing presheaf has two elements, so no $U/x$ is connected. For $p:U\to A$ with $c=p(u)$, the pushout $P_A=X_1\amalg_UA$ has $P_A(v,x)=A(v,c)\times\{\alpha,\beta\}$ for $v\in A$, $x\in T_1$. A chain of arrows of $P_A$ either lies in $A$, or lies in $T_1$, or consists of a chain in $A$, a crossing arrow $(g,\lambda)$ with $\lambda\in\{\alpha,\beta\}$, and a chain in $T_1$; the chain in $A$ together with $g$ is a chain of $A/c$. Write $J=N(A/c)*NT_1$ for the join, and $J_\alpha\cup J_\beta$ for two copies of $J$ glued along $NT_1$. Then $NP_A$ is the pushout of simplicial sets
\[
 NA\longleftarrow N(A/c)\amalg N(A/c)\hookrightarrow J_\alpha\cup J_\beta ,
\]
in which the left map forgets the arrow to $c$ and the right map is a monomorphism; it is therefore a homotopy pushout. A functor $h:A\to A'$ under $U$ induces a map of these diagrams. The categories $A/c$ and $A'/h(c)$ have terminal objects, so their nerves and the joins $J$ are contractible, and the maps on $N(A/c)\amalg N(A/c)$ and on $J_\alpha\cup J_\beta$ are weak equivalences whatever $h$ is. If $h\in W_\infty$, the map $NP_A\to NP_{A'}$ is then a weak equivalence. For $W_R$ one uses the Mayer--Vietoris sequences of the same pushouts.\hfill$\square$

So the assertion ``je sais que c'est le cas si on suppose $T$ isomorphe \`a la cat\'egorie finale $e$'' is false for $W$-cofibrations, and with it the equivalence stated at the end of question~a), also under the axiom W(7), which the localizers $W_R$ satisfy. The corollary of Proposition~7 concerns membership in $W$ of one inclusion, not the cofibration property. What it does give is the statement for \emph{trivial} cofibrations and homology:

\textbf{Proposition.}\ours
\textit{Let $U\hookrightarrow X$ be a sieve whose complement is one object $t$ with only its identity, and suppose that $H=\operatorname{Hom}_X(-,t)|_U$ is not empty. For a nonzero commutative ring $R$, the inclusion is a trivial $W_R$-cofibration if and only if $H$ is totally $W_R$-aspherical, that is, if and only if it is $W_R$-dwy\'erienne.}

\textit{Proof.}
After a cobase change $p:U\to A$ the complement is still $\{t\}$ and the crossing presheaf is $p_!H$. A chain of arrows of $X\amalg_UA$ that does not lie in $A$ is a chain of $A/p_!H$ followed by $t$, so $N(X\amalg_UA)$ is $NA$ with the cone on $N(A/p_!H)$ attached along the forgetful map, and $B(X\amalg_UA)$ is the mapping cone of $B(A/p_!H)\to BA$. The inclusion of $BA$ is an $R$-homology isomorphism exactly when the reduced $R$-homology of $A/p_!H$ vanishes, that is, when $A/p_!H$ has the $R$-homology of a point. Quantifying over $p$ gives the statement.\hfill$\square$

For $W_\infty$ the same computation shows that a one-point sieve is a trivial $W_\infty$-cofibration if and only if every $A/p_!H$ is nonempty and integrally acyclic and the cone attachment does not change fundamental groups. Whether this forces $H$ to be totally $W_\infty$-aspherical is not decided here: a counterexample would be a presheaf $H$ all of whose direct images have acyclic categories of elements, one of them with a nontrivial (necessarily perfect) fundamental group.
\end{dcomment}


et incluant le cas des immersions ouvertes et des immersions fermées?
\end{enumerate}

\textsc{Démonstration} du théorème 3 (p. 107). Il suffit de prouver que si $i$ est $W$-dwyérienne (i.e. $i_0$ strictement $W$-asphérique). alors $i'$ l'est également - car il en résultera que $i'$ est $W$-asphérique et a fortiori $\in W$. Il faut donc vérifier que pour tout $t'\in\dot{U}'=\dot{T}$, l'élément $i^*(t')=(u'\longmapsto\operatorname{Hom}_{X'}(u',t'))$ de $U'^\wedge$ est totalement $W$-asphérique. Or $t'$ provient d'un unique élément $t$ de $\dot{T}$, et on sait qu'on a % [sic] asphérique). alors -> asphérique), alors, period before lower-case word

\[ i'^*(t)=p_!(i^*(t))\quad\texttt{[plutôt }i'^*(t')=p_!(i^*(t))\texttt{]},\qquad\text{où }\ p_!:U^\wedge\longrightarrow U'^\wedge. \]

Par hypothèse, $F=i^*(t)$ est un élément totalement $W$-asphérique de $U^\wedge$, ce qui signifie que pour toute flèche $p:U\longrightarrow U'$ dans Cat de souce $U$, $p_!(F)$ est un élément $W$-asphérique de $U'$. En fait, on a le lemme immédiat (par transitivité $(p'p)_!\simeq p'_!p_!$). % [sic] souce -> source, missing letter

\textbf{Lemme.} \textit{Soit $U$ dans} Cat\textit{, et $F$ dans $U^\wedge$. Si $F$ est totalement $W$-asphérique (resp. si $U/F$ est totalement $W$-asphérique), alors pour toute flèche $p:U\longrightarrow U'$ dans} Cat \textit{de source $U$, $F'=p_!(F)$ est un objet totalement $W$-asphérique de $U'^\wedge$ (resp.}

\msp{page 113}

\textit{$U'/F'$ est un objet totalement $W$-asphérique de} Cat\textit{.)}

\textbf{Corollaire 3} (du théorème 3). \textit{Si $i:U\longrightarrow X$ est une immersion ouverte ou fermée $W$-dwyérienne (resp. totalement $W$-dwyérienne), alors pour toute flèche $p:U\longrightarrow U'$ dans} Cat\textit{, $i':U'\longrightarrow X'=U'\amalg_U X'$} \texttt{[plutôt $U'\amalg_U X$]} \textit{l'est aussi.}

(NB Le cas non respé est essentiellement le théorème 3. Il ne dépend d'aucun axiome supplémentaire sur $W$, mais n'a d'intérêt qu'en vue des équivalences \textcircled{a} et \textcircled{b} (p. 108) qui en résultent, et qui sont valables moyennant W(7 bis) resp. W(7) seulement.)

\dgrblock{dgrb540}{The stronger clause uses comprehensive factorization}
\begin{dcomment}{exposition}{p113:Le cas non respé est essentiellement le théorème 3.}{The stronger clause uses comprehensive factorization}
The lemma proves more than stability of a condition on a presheaf. Its parenthesized assertion says that the category of elements after direct image is itself totally aspherical, so that one may test it by arbitrary functors leaving that category. This does not follow merely by repeating the formula for a second direct image on the original base. The canonical comprehensive factorization supplies the missing step.

Write $V=U/F$, let $a:V\to U$ be the projection, and set $b=pa:V\to U'$. The identity $a_!e_V=F$ gives
\[
 b_!e_V=p_!F=F'.
\]
Let $E=U'/F'$. The canonical functor $v:V\to E$ appearing in the comprehensive factorization of $b$ is final in the ordinary sense: for every $e\in E$, the category $e\backslash V$ is connected and nonempty. Equivalently,
\[
 v_!e_V=e_E.
\]
This is the factorization established at ms.~91--93; it involves connectedness, not contractibility.

Now suppose that $V$ is totally $W$-aspherical, and let $r:E\to Z$ be any functor. Transitivity of ordinary left Kan extension gives
\[
 r_!e_E\cong r_!v_!e_V\cong(rv)_!e_V.
\]
The last presheaf is $W$-aspherical by totality of $V$, so $E$ is totally $W$-aspherical. This proves the parenthesized clause with its full quantifier. The unparenthesized clause is simpler: for $s:U'\to Z$, one has $s_!F'\cong(sp)_!F$, to which the definition applies directly.

The same argument yields a useful formal principle. If $v:V\to E$ is an ordinary final functor and $V$ is totally $W$-aspherical, then $E$ is totally $W$-aspherical. The converse is not being asserted. In particular, an arbitrary final functor need not be a weak homotopy equivalence; the finite example at ms.~100 makes that failure explicit. Here totality of the source is precisely what permits the final functor to be used in every subsequent direct-image calculation.

Corollary~3 follows by applying the two clauses to the gluing weights of each boundary object, with the support calculation already explained after Theorem~3. This proof establishes stability under pushout. It supplies no composition theorem for arbitrary strictly or totally $W$-aspherical functors: in a composite one pulls back a weight of the second functor, whereas the lemma controls left Kan extension. The distinction identifies the mathematical content of the doubt expressed on ms.~113--114. Composition is treated in the commentary at ms.~114: both classes of functors are stable under it, by a different argument.
\end{dcomment}


\textbf{Commentaire de motivation :} les immersions ouvertes ou fermées $W$-dwyériennes, ou totalement $W$-dwyériennes, sont stables par cochangement de base (d'après le corollaire précédent), mais n'ont pas du tout l'air de l'être par composition. À cet égard, les `immersions de Dwyer' (où on suppose l'existence de l'adjoint à droite de $i_0:U\longrightarrow\bar{U}$, dans le cas d'une immersion ouverte) se comportent mieux, étant stables par composition. (\textbf{NB} Il ne semble pas

\msp{page 114}

non plus clair que les flèches strictement $W$-asphériques, ou les totalement $W$-asphériques, soient stables par composition - sorry!) Ma motivation pour m'échiner sur la notion $W$-dwyérienne est double.

\begin{enumerate}[label=\alph*)]
\item J'aimerais avoir des conditions à la fois compréhensibles, et aussi générales que possible, sur une immersion (ouverte p. ex), pour que ce soit une $W$-cofibration, propriété visiblement importante (toute comme la notion duale de $W$-fibration). % [sic] toute comme -> tout comme, agreement error
\end{enumerate}

% ---- from Der-ChXIIb-012_p036-040.tex
\printedpage{78}
% head: A. Grothendieck : Dérivateurs, Ch. XII, Édition des Universités de Montpellier II et Paris VII
\begin{enumerate}
\item[b)] J'ai envie, pour des applications ultérieures, de disposer d'une classe d'immersions ouvertes ayant les propriétés suivantes :
\begin{enumerate}
\item[1°)] Ce sont des $W$-cofibrations, et
\item[2°)] Si $U \overset{i}{\hookrightarrow} X$ est dans la classe, et si $\varepsilon \xrightarrow{j} I$ est l'inclusion ‘source-but’ de la catégorie discrète $\{0,1\}$ dans l'intervalle
\[
\Phi_n = 0 \longrightarrow 1 \longleftarrow 2 \longrightarrow \dots 2n-1 \longleftarrow 2n
\]
$(j(0)=0$, $j(1)=2n)$, $n\geq 2$, l'inclusion $\alpha$ de $(U\times I)\cup(X\times\varepsilon)$ dans $X\times I$ est dans la classe.
\end{enumerate}
Or cette dernière propriété de stabilité n'est \textit{pas} vérifiée par les immersions de Dwyer. (Prendre $i : \{0\} \longrightarrow \boldsymbol{\Delta}^1 = \{0,1\}$, l'inclusion $\alpha$ n'est pas immersion de Dwyer.)
\end{enumerate}

\msp{page 115}

\dgrblock{dgrb541}{The interval test fails because a slice has two maxima}
\begin{dcomment}{examples}{p114}{The interval test fails because a slice has two maxima}
The interval condition asks for a class of admissible inclusions that can also attach both ends of a homotopy. Its categorical form is a pushout product. The original Dwyer condition is too rigid for this purpose, and the example in the text can be checked at a single object.

Take $X=[1]=(0<1)$, $U=\{0\}$, and
\[
 I=\Phi_n=(0\to1\leftarrow2\to\cdots\leftarrow2n),
 \qquad n\geq2.
\]
Let $A=(\{0\}\times I)\cup([1]\times\{0,2n\})$, considered as a full subposet of $[1]\times I$. The inclusion $A\hookrightarrow[1]\times I$ is a sieve: both end vertices of $I$ are sources, and the entire lower row is already in $A$. Its generated cosieve is the whole product, since every upper-row object $(1,k)$ receives an arrow from $(0,k)$.

If this were an original Dwyer immersion, the inclusion into that cosieve would admit a right adjoint. For an inclusion of ordered sets, this means that for every object $x$ of the larger poset the lower set $A/x$ has a greatest element. Test $x=(1,1)$. Its lower set inside $A$ contains exactly
\[
 (0,0),\quad(0,1),\quad(0,2),\quad(1,0).
\]
The strict relations are
\[
 (0,0)<(0,1),\qquad
 (0,2)<(0,1),\qquad
 (0,0)<(1,0).
\]
Thus $(0,1)$ and $(1,0)$ are distinct incomparable maximal elements, and neither is above the other. There is no greatest element. This proves the failure of the adjoint condition for every $n\geq2$.

The calculation also explains why ordinary contractibility of this slice is insufficient to recover the missing adjoint: its nerve is an interval with three edges, hence contractible. An adjoint selects a universal object, a much stronger requirement. Whether another class provides the desired interval stability is a structural question about the interaction of pushout and homotopy, rather than a reason to weaken \emph{terminal object} to \emph{contractible} in this particular criterion. The subsequent use of integration replaces a difficult ordinary pushout by a category designed to retain the homotopy of the entire diagram.
\end{dcomment}


\dgrblock{dgrb542}{A terminal fibre controls all the ordinary direct images}
\begin{dcomment}{exposition}{p114}{A terminal fibre controls all the ordinary direct images}
The last subsection returns to a situation in which the apparent loss of higher information can be controlled. For a $W$-proper functor over a category with a terminal object, the terminal-fibre inclusion is $W$-coaspherical. Under W(5) it is ordinary final, which is the property needed for total asphericity. Making this explicit shortens the proof and explains the editors' repeated suggestion to replace $W$-properness by $W_0$-properness.

Let $f:X\to I$ be $W$-proper, let $i_0$ be terminal in $I$, and write $k:X_{i_0}\hookrightarrow X$. The properness criterion says that $k$ is $W$-coaspherical: every undercategory $x\backslash X_{i_0}$ is $W$-aspherical. If W(5) holds, a $W$-aspherical category is nonempty and connected. Hence $k$ is an ordinary final functor, and the pointwise formula gives
\[
 (k_!e_{X_{i_0}})(x)=\pi_0(x\backslash X_{i_0})=\{*\}.
\]
For every $p:X\to Y$ it follows that
\[
 (pk)_!e_{X_{i_0}}
 \cong p_!k_!e_{X_{i_0}}
 \cong p_!e_X.                                      \tag{C1}
\]
If $X_{i_0}$ is totally $W$-aspherical, the left side of (C1) is an aspherical presheaf on $Y$ for every $p$, proving Proposition~11 of ms.~115. This proposition has the same number as the counterexample on ms.~107; the manuscript page is needed to distinguish them. The proof uses no W(7) or W(7 bis).

The general component formula in Proposition~12 explains how (C1) fits into a family. $W$-properness together with W(5) identifies $\pi_0(X_i)$ with $\pi_0(X/i)$, naturally in $i$, and therefore supplies the transition maps even when $f$ is not an opfibration. For an arbitrary functor $f$, there is a canonical bijection
\[
 \varinjlim_{i\in I}\pi_0(X/i)\ \cong\ \pi_0(X).        \tag{C2}
\]
One can see it without choosing a cleavage. A representative $(x,\beta:f(x)\to i)$ in $X/i$ has the same image in the colimit as $(x,\id_{f(x)})$, by the transition along $\beta$. A morphism $x\to x'$ identifies the latter representatives through the induced object of $X/f(x')$. Thus the colimit imposes exactly the zigzag relation defining components of $X$. Conversely each morphism in a comma category, and every transition between comma categories, preserves that zigzag component. These two observations construct mutually inverse maps in (C2).

Replacing $\pi_0(X/i)$ by $\pi_0(X_i)$ proves Proposition~12. The dual proof printed in the manuscript is equally valid, but after passing to presheaves its colimit is indexed by $I^{\mathrm{op}}$; the printed $I$ at ms.~117 suppresses this reversal. The precise proper/smooth criteria, their duality and their stability under composition and base change are given in \citet[Theorem~1.2 and Proposition~4.7]{Maltsiniotis2005Asphericite}. The result here concerns ordinary components and ordinary Kan extensions, so the connectivity hypothesis cannot be omitted just because the discussion takes place in a homotopical setting.
\end{dcomment}


\dgrblock{dgrb543}{Dwyer attachments and higher categorical pushouts}
\begin{dcomment}{later}{p114}{Dwyer attachments and higher categorical pushouts}
The adjoint condition belongs to the Dwyer maps of \citet[Definition~4.1 and Proposition~4.3]{Thomason1980}. One subsequent correction matters when comparing them with model-category cofibrations. Classical Dwyer maps are not closed under retracts, and some Thomason cofibrations are not Dwyer maps. \citet[Definition~1, Lemma~4 and Proposition~2]{Cisinski1999Dwyer} repairs the relevant argument with pseudo-Dwyer maps: these are retract-stable and include every Thomason cofibration. Thus the manuscript's search for a broader attachment condition addresses a distinction that cannot be removed by simply identifying the original Dwyer maps with all cofibrations.

The desired stability properties concern whether an ordinary attachment retains the information needed after passage to a homotopy theory. For classical Dwyer maps, a stronger result is now available. \citet[Theorem~1.5]{HackneyOzornovaRiehlRovelli2024} prove that, for a pushout $D=B\amalg_A C$ along a Dwyer inclusion $A\hookrightarrow B$, the comparison
\[
 NB\amalg_{NA}NC\longrightarrow ND
\]
is a weak \emph{categorical} equivalence in the Joyal model structure. Their general Theorem~1.6, in its expanded statement on p.~2179, treats discretely flat cofibrations. These are defined by requiring the corresponding functor of discrete simplicial categories to preserve Bergner weak equivalences under cobase change, together with injectivity on objects. Proposition~3.3 supplies this property for Dwyer maps; Example~3.5 includes co-Dwyer maps and pseudo-Dwyer maps. This gives one modern class that accommodates both orientations requested on p.~111.

This higher categorical assertion concerns the comparison between the ordinary categorical pushout and a pushout in the category of $(\infty,1)$-categories. The weak homotopy type of a category discards all noninvertible-arrow information, whereas its image among $(\infty,1)$-categories retains that information. Consequently the two ambient weak equivalence classes must be specified before comparing definitions of flatness. The theorem just cited does not identify discretely flat inclusions with the manuscript's $W$-Dwyer immersions, nor does it by itself establish the particular interval-product stability requested here. Those further identifications would require additional arguments. The concrete calculations in the preceding blocks keep the original question at the level at which it is posed.
\end{dcomment}


\dgrblock{dgrb544}{Composition of aspherical functors and of dwyerienne immersions}
\begin{dcomment}{development}{p114}{Composition of aspherical functors and of dwy\'erienne immersions}
Grothendieck doubts that strictly or totally $W$-aspherical functors, and the two classes of dwy\'erienne immersions, are stable under composition. For functors the doubt is unfounded; for immersions it is justified for some localizers and not for others. The statements below were not found in the sources consulted; no priority is claimed.

\textbf{Theorem (functors compose).}\ours
\textit{Let $W$ be a fundamental localizer.}
\begin{enumerate}[label=(\alph*)]
\item \textit{If $f:X\to Y$ is strictly $W$-aspherical and $H$ is a totally $W$-aspherical presheaf on $Y$, then $f^*H$ is a totally $W$-aspherical presheaf on $X$. Consequently strictly $W$-aspherical functors are stable under composition.}
\item \textit{Totally $W$-aspherical functors are stable under composition.}
\item \textit{If $B$ is a totally $W$-aspherical category and $D:B\to\mathrm{Cat}$ a functor with totally $W$-aspherical values, the cofibred category $\int_BD$ is totally $W$-aspherical.}
\end{enumerate}

\textit{Proof.}
(a) Let $p:X\to X'$ be arbitrary. For $y\in Y$ put $Q_y=p_!\operatorname{Hom}_Y(f(-),y)$, a presheaf on $X'$, covariant in $y$. Each $Q_y$ is totally $W$-aspherical, by the hypothesis on $f$ and the lemma of ms.~112. Let $Z$ be the category with full subcategories $X'$ and $Y$, with $\operatorname{Hom}_Z(x',y)=Q_y(x')$ and no arrows from $Y$ to $X'$; let $j:Y\hookrightarrow Z$ be the inclusion and $F=j_!H$. Then $F|_Y=H$ and
\[
 F|_{X'}=\int^{y\in Y}Q_y\times H(y)\cong p_!\Bigl(\int^{y}\operatorname{Hom}_Y(f(-),y)\times H(y)\Bigr)=p_!f^*H=:K .
\]
Since $H$ is total, $Z/F$ is $W$-aspherical. The inclusion $l:X'/K\hookrightarrow Z/F$ is a sieve. Its comma category over an object $(x',k)$ of $X'/K$ has a terminal object. Its comma category over $(y,h)$, $h\in H(y)$, has as objects the pairs $(x',q)$ with $q\in Q_y(x')$ (the element of $K(x')$ is then the class of $(q,h)$) and as arrows the arrows of $X'$ that respect $q$; it is $X'/Q_y$, which is $W$-aspherical. By the relative form of Theorem~A that defines a fundamental localizer, $l\in W$, and so $X'/K$ is $W$-aspherical. As $p$ was arbitrary, $f^*H$ is total. If $g:Y\to Z_0$ is also strictly aspherical, apply this to $H=\operatorname{Hom}_{Z_0}(g(-),z)$, whose restriction along $f$ is $\operatorname{Hom}_{Z_0}(gf(-),z)$.

(b) Let $f:X\to Y$ and $g:Y\to Z_0$ be totally aspherical and $z\in Z_0$. The functor $f_z:X/z\to Y/z$, $(x,\gamma)\mapsto(f(x),\gamma)$, has comma category over $(y,\beta)$ canonically isomorphic to $X/y$, because the arrow to $z$ is determined as $\beta g(h)$ by $h:f(x)\to y$. So $f_z$ is totally aspherical, hence strictly aspherical by Proposition~8. The final presheaf of $Y/z$ is total because $g$ is totally aspherical; by (a) its restriction, the final presheaf of $X/z$, is total. That is, $X/z$ is a totally $W$-aspherical category.

(c) Put $E=\int_BD$, with objects $(b,x)$, $x\in D(b)$, and projection $\pi:E\to B$. For $b\in B$ the functor $D(b)\to E/b$, $x\mapsto(b,x,\mathrm{id}_b)$, has the left adjoint $(b',x,\beta)\mapsto D(\beta)(x)$. A right adjoint is final. So $E/b$ is a final image of $D(b)$ and is total; $\pi$ is a totally aspherical functor, as is $B\to e$, and (b) applies.\hfill$\square$

\textbf{Proposition (immersions: two counterexamples).}\ours
\textit{For $W=W_0$ and for $W=W_1$ neither the $W$-dwy\'erienne nor the totally $W$-dwy\'erienne sieves are stable under composition.}

\textit{Proof.}
Let $V$ be the discrete category on two objects for $W_0$, and the category with two objects and two parallel arrows for $W_1$; $BV$ is $S^0$, respectively $S^1$. Let $B=V^{\triangleleft}$ be $V$ with an initial object $o$ adjoined. Let $X$ consist of $B$ and of a second copy $V_0$ of $V$, with
\[
 X(v_0,b)=B(v,b),\qquad X(b,v_0)=\varnothing\qquad(v\in V,\ b\in B),
\]
and composition induced by $B$. The sieve $V_0\hookrightarrow X$ generates the cosieve $V_0\cup V$, on which the inclusion of $V_0$ has the right adjoint $v\mapsto v_0$; it is a Dwyer immersion. The functor $X\to B$ forgetting the copy is left adjoint to the inclusion $B\subset X$, so $BX\simeq BB$ is contractible. A final functor is bijective on $\pi_0$ and surjective on $\pi_1$; hence $X$ is totally $W_0$- and totally $W_1$-aspherical. Let $Y=X^{\triangleright}$ be $X$ with a terminal object $t$ adjoined. The sieve $X\hookrightarrow Y$ is totally $W$-dwy\'erienne, its only boundary category being $X/t=X$. For the composite sieve $V_0\hookrightarrow Y$ the boundary category at $t$ is $V_0/t\simeq V$ and the boundary presheaf is the final presheaf of $V_0$; neither is $W$-aspherical.\hfill$\square$

The targets $Y$ have six objects and $15$, respectively $21$, arrows.

\textbf{Theorem (total immersions for $W_\infty$).}\ours
\textit{For $W=W_\infty$ the totally $W$-dwy\'erienne sieves are stable under composition.}

\textit{Proof.}
We use the criterion of the commentary on Chapter~II, ms.~36: a category $A$ is totally $W_\infty$-aspherical if and only if it is nonempty and connected and every connected component of the category of elements of every functor $A\to\mathrm{Set}$ is weakly contractible; and then those components are themselves totally $W_\infty$-aspherical and any two objects of $A$ have a common successor. A nonempty cosieve $S$ of such an $A$ is the category of elements of the functor with value a point on $S$ and empty elsewhere; it is connected, because two of its objects have a common successor in $A$, which lies in $S$. So $S$ is totally $W_\infty$-aspherical.

Let $U\subset X\subset Y$ be totally $W_\infty$-dwy\'erienne sieves and $y\in Y$ with $U/y\neq\varnothing$. Then $X/y$ is total (it has a terminal object if $y\in X$). The objects $(x,\beta)$ of $X/y$ with $U/x\neq\varnothing$ form a nonempty cosieve $S_y$, total by the preceding paragraph, and the functor $(x,\beta)\mapsto U/x$ on $S_y$ has total values. By part~(c) of the theorem its cofibred category $G$ is total. The functor $G\to U/y$ that composes $u\to x\to y$ has the left adjoint $(u,\gamma)\mapsto((u,\gamma),(u,\mathrm{id}_u))$, so it is final, and $U/y$ is total.\hfill$\square$

The same proof works for the homology localizers $W_R$, with the homological form of the criterion ($R$-acyclic in place of weakly contractible). For the strict class, the $W_R$-dwy\'erienne sieves are also stable under composition: the restriction of a totally $W_R$-aspherical presheaf to a cosieve on which it is nonempty is totally $W_R$-aspherical, after which part~(a) applies; the proof of the restriction statement is not reproduced here. For $W_\infty$-dwy\'erienne sieves the question is undecided. It comes down to the question left open in the commentary at ms.~111: does a presheaf all of whose direct images have integrally acyclic categories of elements have only weakly contractible ones?
\end{dcomment}


\dgrblock{dgrb545}{The endpoint product is not a cofibration}
\begin{dcomment}{development}{p114}{The endpoint product is not a cofibration}
The class asked for in b) does not exist for $W_\infty$: the inclusion $\alpha$ built from the simplest sieve is not a $W_\infty$-cofibration at all. This is stronger than the observation that $\alpha$ is not a Dwyer immersion.

\textbf{Theorem.}\ours
\textit{Let $n\geq2$, $I=\Phi_n$, $\varepsilon=\{0,2n\}\subset I$, and let $\alpha:A\hookrightarrow P$ be the inclusion of $A=(\{0\}\times I)\cup([1]\times\varepsilon)$ in $P=[1]\times I$, deduced from $\{0\}\hookrightarrow[1]$. Let $C$ be the four-object category of the commentary at ms.~100, of the homotopy type of $S^2$.}
\begin{enumerate}[label=(\alph*)]
\item \textit{If $W$ is a fundamental localizer for which a category of the homotopy type of $S^3$ is not $W$-aspherical (for instance $W_\infty$, $W_R$ with $R\neq0$, or the $m$-equivalences with $m\geq3$), then $\alpha$ is not a $W$-cofibration.}
\item \textit{If $C$ is not $W$-aspherical (which includes the $2$-equivalences), then $\alpha$, which belongs to $W$, is not a trivial $W$-cofibration.}
\end{enumerate}

\textit{Proof.}
Write $l_k=(0,k)$ and $u_k=(1,k)$. The objects of $A$ are $l_0,\dots,l_{2n},u_0,u_{2n}$; its nonidentity arrows are those of the lower copy of $I$ and $l_0\to u_0$, $l_{2n}\to u_{2n}$, and no two of them are composable. Both $A$ and $P$ have contractible classifying spaces, so $\alpha\in W_\infty\subseteq W$.

Recall that $C$ has objects $x,z,a,b$ and arrows $\alpha',\beta':x\to z$, $r:z\to a$, $s:z\to b$ with $r\alpha'=r\beta'$, $s\alpha'=s\beta'$. Let $D$ be either $C$ or the category $C^{\triangleright}$ obtained by adjoining a terminal object $\omega$. Define $p:A\to D$ by
\[
 p(l_0)=z,\qquad p(u_0)=b,\qquad p(l_k)=a\ (1\leq k\leq2n),\qquad p(u_{2n})=a,
\]
sending $l_0\to l_1$ to $r$, $l_0\to u_0$ to $s$ and all other arrows to $\mathrm{id}_a$. Let $Q=P\amalg_AD$. Since $A$ is a sieve in $P$, the category $Q$ consists of $D$, of the objects $u_1,\dots,u_{2n-1}$ with the arrows of the upper copy of $I$ between them, and of the crossing sets $Q(v,u_k)=(p_!H_k)(v)$, where $H_k=\operatorname{Hom}_P(-,u_k)|_A$ (Theorem~1).

The presheaf $H_1$ is the final presheaf of the sieve $A_1=\{l_0,l_1,l_2,u_0\}$ of $A$ and is empty elsewhere. So $(p_!H_1)(v)=\pi_0(v\backslash A_1)$, the comma category being formed with the restriction $A_1\to D$ of $p$. For $v=z$ it is $A_1$ itself, a tree; for $v=a$ it is $l_2\to l_1$; for $v=b$ it is the point $u_0$; for $v=x$ it has the five objects $(l_0,\alpha'),(l_0,\beta'),(l_1,r\alpha'),(l_2,r\alpha'),(u_0,s\alpha')$ and is connected; for $v=\omega$ it is empty. Hence the object $u_1$ receives exactly one arrow from each object of $C$ and none from $\omega$. For $2\leq k\leq2n-1$ the category $A/u_k$ is nonempty and connected and $p$ is constant on it with value $a$, so $p_!H_k=D(-,a)$; the arrows between $u_2,\dots,u_{2n-1}$ induce identities of $D(-,a)$, and $u_2\to u_1$ induces the unique map $D(-,a)\to p_!H_1$.

Let $R\subset Q$ be the full subcategory on $D$ and $u_1$. The functor $\rho:Q\to R$ that is the identity on $R$, sends $u_k$ to $a$ for $k\geq2$, is the identity on the crossing sets $D(v,a)$, sends the arrows between these $u_k$ to $\mathrm{id}_a$ and $u_2\to u_1$ to the arrow $a\to u_1$, satisfies $\rho|_R=\mathrm{id}_R$; and the arrows $a\to u_k$ corresponding to $\mathrm{id}_a$ are the components of a natural transformation from $\rho$, followed by the inclusion of $R$, to $\mathrm{id}_Q$. So $BQ\simeq BR$.

(a) Take $D=C^{\triangleright}$. Then $p$ is a functor between categories with contractible classifying spaces, so $p\in W_\infty\subseteq W$; it is a morphism of $A\backslash\mathrm{Cat}$ from $\mathrm{id}_A$ to $p$, and its image under cobase change along $\alpha$ is $P\to Q$. The category $R$ is $C$ with the two terminal objects $\omega$ and $u_1$ adjoined, with no arrow between them: its nerve is the union of two cones on $NC$, so $BR\simeq\Sigma BC\simeq S^3$. Thus $P$ is $W$-aspherical and $Q$ is not, and $P\to Q$ is not in $W$.

(b) Take $D=C$. Then $R$ is $C$ with one terminal object adjoined and $BQ\simeq BR$ is contractible. The cobase change $C\to Q$ of $\alpha$ along $p$ goes from a category that is not $W$-aspherical to one that is, so it is not in $W$.\hfill$\square$

For $n=2$ the category $Q$ of (a) has $8$ objects and $30$ arrows. Part~(b) also excludes the cofibration property whenever a $W$-cofibration that belongs to $W$ is a trivial $W$-cofibration. For $W_\infty$ this is \citet[Lemma~1.10]{BataninBerger2017}, applied to Thomason's model structure, which is left proper; there it adds nothing to (a).

\textbf{Corollary (no class of the kind requested).}\ours
\textit{Let $W$ be as in (a) and $n\geq2$. If a class $\mathcal K$ of sieves consists of $W$-cofibrations and contains, with $i:U\hookrightarrow X$, the inclusion of $(U\times I)\cup(X\times\varepsilon)$ in $X\times I$, then every member of $\mathcal K$ is an isomorphism. The same holds for $W$ as in (b) and classes of trivial $W$-cofibrations.}

\textit{Proof.}
$W$-cofibrations, and trivial $W$-cofibrations, are stable under retracts in the category of arrows, because a retract of a functor in $W$ is in $W$. Let $i:U\hookrightarrow X$ be in $\mathcal K$. If there is an arrow $u\to x$ with $u\in U$, $x\notin U$, then $\{0\}\hookrightarrow[1]$ is a retract of $i$ (by the functor $X\to[1]$ with fibre $U$ over $0$, and the functor $[1]\to X$ choosing that arrow), the endpoint product $\alpha$ of $\{0\}\hookrightarrow[1]$ is a retract of that of $i$, and the theorem is contradicted. If there is no such arrow, then $X=U\amalg V$; if $V\neq\varnothing$, the endpoint product of $i$ is the inclusion of $(U\times I)\amalg(V\times\varepsilon)$ in $(U\amalg V)\times I$, which is in $\mathcal K$ and has the crossing arrow $(v,0)\to(v,1)$, so the first case applies to it.\hfill$\square$

At the other extreme, for $W_0$ every functor is a $W_0$-cofibration, because $\pi_0:\mathrm{Cat}\to\mathrm{Set}$ preserves pushouts, and all sieves form a class with the two properties. The same is true for $W_1$ and for every fundamental localizer containing $W_1$: the pushout along a functor injective on objects has the same fundamental groupoid as the cofibred category of the span; this proof is not reproduced here. By \citet[Proposition~9.2.17]{Cisinski2006} a fundamental localizer contains $W_1$ exactly when a category of the homotopy type of $S^2$ is aspherical for it. So for the localizers of homotopy theory proper the endpoint product has to be replaced by a homotopy-invariant construction, which is what the integration of the following pages provides.
\end{dcomment}


\subsection*{2.5\quad Complément : critère d'asphéricité totale d'une catégorie $W$-propre sur une autre.}

\dgrblock{dgrb546}{Why the interval test exceeds Dwyer maps}
\begin{dcomment}{examples}{p115:Complément : critère d'asphéricité totale}{Why the interval test exceeds Dwyer maps}
The preceding paragraph asks for enough cofibrations to build homotopies with prescribed endpoints.  Its product with the zigzag interval is therefore a structural test, rather than an incidental counterexample.  We can make the obstruction completely visible in a finite ordered set.

Write $P=[1]\times\Phi_n$ and
\[
 A=(\{0\}\times\Phi_n)\cup([1]\times\{0,2n\}).
\]
Here $[1]=\{0<1\}$, and the even vertices of $\Phi_n$ map to their adjacent odd vertices.  The inclusion $A\subset P$ is a sieve: a predecessor of an object in its bottom row is again in that row, and a predecessor of either of its two top endpoints is already in $A$.  Every object $(1,k)$ receives a map from $(0,k)$, so the cosieve generated by $A$ is all of $P$.

For a full inclusion of ordered sets, the existence of a right adjoint is particularly concrete.  If $r$ is right adjoint to $A\hookrightarrow P$, then $r(p)$ must be the greatest element of $\{a\in A:a\leq p\}$.  At $p=(1,1)$, this set has the two incomparable maximal elements $(1,0)$ and $(0,1)$; for $n\geq2$ the other top endpoint cannot lie below $p$.  There is consequently no greatest element, hence no right adjoint.  This is exactly the obstruction in the definition of a Dwyer map, since the relevant cosieve is already $P$.  The original inclusion $\{0\}\hookrightarrow[1]$ is Dwyer, with right adjoint the constant functor to $0$.  Thus the endpoint product operation genuinely leaves that class.

The definition and generated-cosieve criterion are those of \citet[Definition~4.1 and Lemma~5.2]{Thomason1980}.  There is a second, distinct limitation in the history of this class: retracts of Dwyer maps need not be Dwyer.  \citet[pp.~228--231, especially Proposition~2]{Cisinski1999Dwyer} gives the counterexample and repairs the properness argument for Thomason's model structure using pseudo-Dwyer maps.  The interval calculation above concerns a different closure operation.  Both facts explain why Grothendieck seeks a larger homotopically defined class rather than treating the adjoint criterion as the final definition of a cofibration.
\end{dcomment}


Soit
\[
f : X \longrightarrow I
\]
une catégorie propre sur une catégorie d'indices $I$. Si $I$ a un objet final $i_0$, alors $X/i_0$ s'identifie à $X$, et par suite l'inclusion
\[
X_{i_0} \longrightarrow X
\]
de la fibre de $X$ en $i_0$ est $W$-asphérique \texttt{[plutôt }$W$\texttt{-coasphérique]}, et a fortiori est dans $W$. Donc $X$ est asphérique si et seulement si $X_{i_0}$ l'est. Nous voulons étendre ce résultat en un critère de \textit{totale} $W$-asphéricité, de façon précise prouver la

\textbf{Proposition 11 :} \textit{Soit $f : X \longrightarrow I$ $W$-propre, et soit $i_0$ un objet final de $I$. Si $X_{i_0}$ est totalement $W$-asphérique} (déf. 1, page 88), \textit{il en est de même de $X$.} (Sous réserve que $W$ satisfasse l'axiome de connexité W(5).)

La démonstration sera donnée plus bas, après quelques préliminaires ayant un intérêt indépendant. Tout d'abord, rappelons que si $X$ est $W$-propre sur $I$, alors il donne naissance à un foncteur
\[
\begin{aligned}
I \quad &\longrightarrow \quad \mathrm{Hot}_W\\
i \quad &\longmapsto \quad \mathrm{hot}_W(X_i)(\simeq \mathrm{hot}_W(X/i));
\end{aligned}
\]
en le composant avec le foncteur
\[
\pi_0 : \mathrm{Hot}_W \longrightarrow \mathrm{Ens}
\]
(sous réserve que l'axiome de connexité W(5) (p. 10) soit satisfait)
% foot: Version 7 décembre 2002
\printedpage{79}
% head: A. Grothendieck : Dérivateurs, Ch. XII, Édition des Universités de Montpellier II et Paris VII
\msp{page 116}
on trouve un foncteur
\[
i \longmapsto \underbrace{\pi_0(X_i)}_{\simeq\, \pi_0(X/i)} : I \longrightarrow \mathrm{Ens}.
\]
Ceci posé, on a le résultat sorital

\dgrblock{dgrb547}{Components of a proper family}
\begin{dcomment}{exposition}{p115:La démonstration sera donnée plus bas}{Components of a proper family}
Grothendieck wants to transfer a condition on every ordinary left Kan extension from the last fibre to the whole category.  A weak equivalence of categories alone would not accomplish this: the functor $p_!$ here is the underived left Kan extension between categories of set-valued presheaves.  The useful consequence of properness is that it controls the components of all the relevant comma categories.

We use $\widehat Y=\operatorname{Fun}(Y^{\mathrm{op}},\mathrm{Set})$ and write $1_X$ for the terminal presheaf.  For $p:X\to Y$ one has
\[
 (p_!1_X)(y)=\pi_0(y\backslash X),
\]
where an object of $y\backslash X$ is $(x,a:y\to p(x))$.  In this chapter, ``totally $W$-aspherical'' means that $Y/(p_!1_X)$ is $W$-aspherical for every $Y$ and every $p:X\to Y$.  This is the definition on manuscript pp.~88--89; the editors explicitly distinguish it from the earlier use of total asphericity in \emph{Pursuing Stacks} \citep[Definition~1]{Grothendieck1990XII}.  It is this presheaf condition that Proposition~11 transfers.

The calculation of components has a direct proof that also fixes its variance.  For any $f:X\to I$, without a properness hypothesis, the categories $X/i$ form a covariant diagram by postcomposition.  There is a canonical bijection
\[
 \operatorname*{colim}_{i\in I}\pi_0(X/i)\ \longrightarrow\ \pi_0(X).
 \tag{A1}
\]
Indeed, every object $x$ occurs as $(x,\mathrm{id}_{f(x)})$ in $X/f(x)$, so the map is surjective.  The class of an arbitrary $(x,a:f(x)\to i)$ is identified in the colimit with the class of $(x,\mathrm{id})$, by the transition functor associated to $a$.  If $v:x\to x'$ is a morphism of $X$, the two resulting classes are equal: move the first object along $f(v)$ into $X/f(x')$, where $v$ becomes a morphism to $(x',\mathrm{id})$.  These are precisely the relations defining $\pi_0(X)$.  This proves injectivity as well.

If $f$ is $W_0$-proper, the fibre inclusion $X_i\to X/i$ induces a bijection on components.  Transporting the covariant structure through these bijections supplies the functor $i\mapsto\pi_0(X_i)$ and turns (A1) into Proposition~12.  There is no choice of transport functors on the fibres in this argument.  Under W(5), the inclusion $W\subseteq W_0$ makes $W$-properness sufficient.  The two footnotes recommending $W_0$-properness identify the exact hypothesis needed for this calculation.

For the relative formula, apply the same argument to $y\backslash X\to I$.  Its fibre over $i$ is $y\backslash X_i$.  The projection $y\backslash Y\to Y$ is a discrete opfibration, hence proper; properness is preserved by base change and composition.  Consequently
\[
 (p_!1_X)(y)\cong
 \operatorname*{colim}_{i\in I}(p_{i!}1_{X_i})(y).
 \tag{A2}
\]
Precomposition by a map $y'\to y$ commutes with every functor used to construct this bijection.  Thus (A2) is an isomorphism of presheaves, not merely a collection of unrelated bijections.  This supplies both the naturality and the universal property asserted in Corollary~1.
\end{dcomment}


\textbf{Proposition 12 :} \textit{Soit $f : X \longrightarrow I$ $W$-propre, et supposons que $W$ satisfasse l'axiome de connexité} W(5) (sinon, la $W$-propreté ne nous sera d'aucune utilité \dots) (\footnote[81]{Il vaut mieux dire : Soit $f$ $W_0$-propre, où $W_0$ est le localiseur de la 0-connexité \dots}). \textit{On a alors une bijection canonique}
\[
\pi_0(X) \xleftarrow{\ \sim\ } \varinjlim_{I} \pi_0(X_i), \tag*{$(*)$}
\]
\textit{où la flèche est induite par les applications canoniques $\pi_0(X_i) \longrightarrow \pi_0(X)$.}

\textsc{Démonstration.} Passant à $f^\circ : X^\circ \longrightarrow I^\circ$, on est ramené à prouver le même résultat pour $X$ \textit{$W$-lisse} sur $I$, avec la seule différence que le foncteur
\[
i \longmapsto \pi_0(X_i)
\]
sur $I$ sera contravariant, i.e. un préfaisceau, ce qui offre l'avantage de l'intuition attachée au formalisme des préfaisceaux. Soit
\[
g : I \longrightarrow e, \qquad \text{d'où } gf = p : X \longrightarrow e,
\]
il est immédiat, par le calcul standard
\[
p_!(F) = \pi_0(X/F) \qquad (F \in \mathrm{Ob}\, X^\wedge)
\]
(où on a identifié $e^\wedge$ à (Ens)), qu'on a
\[
\pi_0(X) = p_!(e_X) = g_!(f_!(e_X)).
\]
Or le préfaisceau
\[
i \longmapsto \pi_0(X_i) \qquad (\simeq \pi_0(i\backslash X), \text{ car } X\ W\text{-lisse sur } I)
\]
\msp{page 117}
n'est autre que $f_!(e_X)$. D'autre part, pour un préfaisceau quelconque $G$ sur $I$, on a
\[
g_!(G) \simeq \pi_0(I/G) \simeq \varinjlim_{I} G(i),
\]
et appliquant ceci à $f_!(e_X)$, on trouve la formule $(*)$ de la proposition 12.

Revenons à la situation avec $f$ \textit{$W$-propre,} et considérons un foncteur
\[
p : X \longrightarrow Y,
\]
nous sommes intéressés à ‘calculer’
\[
p_!(e_X) = (y \longmapsto \pi_0(y\backslash X)) \in \mathrm{Ob}\, Y^\wedge.
\]
Considérons pour ceci le diagramme suivant
% foot: Version 7 décembre 2002
\printedpage{80}
% head: A. Grothendieck : Dérivateurs, Ch. XII, Édition des Universités de Montpellier II et Paris VII
\[
\begin{tikzcd}
y\backslash X_i \arrow[r] \arrow[d] & y\backslash X \arrow[r] \arrow[d] & y\backslash Y \arrow[d] \\
X_i \arrow[r] \arrow[d] & X \arrow[r, "p"] \arrow[d, "f"] & Y \\
\underset{\simeq\, e}{\{i\}} \arrow[r] & I\rlap{\quad ,} &
\end{tikzcd}
\]
où les trois carrés sont cartésiens. Comme $y\backslash Y \longrightarrow Y$ est $W$-propre, il en est de même de $y\backslash X \longrightarrow X$ par changement de base, donc du composé $y\backslash X \longrightarrow X \longrightarrow I$ par transitivité. Le diagramme montre d'autre part que la fibre de $y\backslash X$ en $i \in \mathrm{Ob}\, I$ s'identifie à $y\backslash X_i$.

\msp{page 118}
La proposition 12 nous donne donc la formule
\[
\pi_0(y\backslash X) = \varinjlim_{I} \pi_0(y\backslash X_i),
\]
laquelle formule ne peut s'empêcher d'être fonctorielle en $y$, tout comme celle de prop. 12 l'est dans $\underline{\mathrm{Prop}}_W I$. Or pour $y$ variable dans $Y$, le préfaisceau sur $Y$
\[
y \longmapsto \pi_0(y\backslash X_i)
\]
n'est autre que $p_{i\,!}(e_{X_i})$, où
\[
p_i : X_i \longrightarrow Y
\]
est le composé $X_i \longrightarrow X \xrightarrow{p} Y$ (induit par $p$). On a ainsi prouvé

\textbf{Corollaire 1} (de prop. 12). \textit{Soit $f : X \longrightarrow I$ $W$-propre} (\footnote[82]{dire $W_0$-propre \dots}), \textit{et $p : X \longrightarrow Y$ un foncteur quelconque. Soit, pour tout $i \in \mathrm{Ob}\, I$, $p_i : X_i \longrightarrow Y$ le foncteur induit sur la fibre $X_i$ de $X$ en $i$. On a alors une loi} fonctorielle \textit{naturelle}
\[
i \longmapsto p_{i\,!}(e_{X_i}) : I \longrightarrow Y^\wedge,
\]
\textit{et on a la formule}
\[
p_!(e_X) \xleftarrow{\ \sim\ } \varinjlim_{I} p_{i\,!}(e_{X_i}), \tag*{$(*)$}
\]
\textit{où la flèche est induite par les morphismes canoniques}
\[
p_{i\,!}(e_{X_i}) \longrightarrow p_!(e_X)
\]
% foot: Version 7 décembre 2002
\printedpage{81}
% head: A. Grothendieck : Dérivateurs, Ch. XII, Édition des Universités de Montpellier II et Paris VII
\msp{page 119}
\textit{induites par les inclusions $X_i \longrightarrow X$.}

La loi fonctorielle de $p_{i\,!}(e_{X_i})$ en $i$ a été explicitée en cours de démonstration.

Notons que lorsque $X$ sur $I$ est cofibrée et scindée, i.e. provient d'un foncteur
\[
i \longmapsto X_i, \quad I \longrightarrow \mathrm{Cat},
\]
alors la loi fonctorielle de $p_{i\,!}(e_{X_i})$ en $i$ s'explicite de façon particulièrement agréable : à une flèche
\[
\alpha : i \longrightarrow j, \qquad \text{d'où } \alpha_* : X_i \longrightarrow X_j,
\]
est associé
\[
p_!(\alpha) : p_{i\,!}(e_{X_i}) \longrightarrow p_{j\,!}(e_{X_j}),
\]
en notant qu'on a un diagramme
\[
\begin{tikzcd}
X_i \arrow[rr, "\alpha_*"] \arrow[dr, "p_i"', ""{name=A}] & & X_j \arrow[dl, "p_j", ""'{name=B}] \\
& Y &
\arrow[from=A, to=B, bend left, "\lambda_\alpha"] % [resolved: see DER_XIIb-a_corrections.csv]
\end{tikzcd}
\]
avec une flèche de commutativité
\[
\lambda_\alpha : p_i \longrightarrow p_j\alpha_*
\]
(pas \texttt{[un]} isomorphisme en général), d'où
\[
p_{i\,!}(e_{X_i}) \xrightarrow{(\lambda_\alpha)_!} p_{j\,!}\underbrace{\alpha_{*\,!}(e_{X_i})}_{\longrightarrow\, e_{X_j}} \longrightarrow p_{j\,!}(e_{X_j}),
\]
plus de doute c'est lui! (\footnote[83]{Cette flèche \texttt{[la flèche }$(\lambda_\alpha)_!$\texttt{]} devrait être explicitée avec soin!}).

\dgrblock{dgrb548}{The comparison map requested in footnote 83}
\begin{dcomment}{exposition}{p119:plus de doute c'est lui!}{The comparison map requested in footnote 83}
The point of the split-opfibration calculation is to see the transition map before passing to components.  This also answers the request in footnote~83 without relying on a convention for mates of adjunctions.

Let $\alpha:i\to j$, let $\alpha_*:X_i\to X_j$ be the chosen transport, and let $\lambda_x:p_i(x)\to p_j(\alpha_*x)$ be the image under $p$ of the chosen cocartesian arrow over $\alpha$.  For every $y\in Y$ there is a functor
\[
 y\backslash X_i\longrightarrow y\backslash X_j,
 \qquad (x,a)\longmapsto(\alpha_*x,\lambda_x a).
 \tag{A3}
\]
A morphism $v:(x,a)\to(x',a')$ is sent to $\alpha_*v$.  This is a morphism of the target comma category because
\[
 p_j(\alpha_*v)\lambda_x a
   =\lambda_{x'}p_i(v)a=\lambda_{x'}a'.
\]
Passing to components gives the displayed map $p_!(\alpha)$ in the manuscript.  It is natural in $y$ because both routes for a map $b:y'\to y$ replace $a$ by $ab$.  For two composable base arrows, the split identities and the cocartesian composition law identify the two composites of (A3).  Hence these maps form the asserted covariant diagram of presheaves.

To recover the factorization written by Grothendieck, first keep the intermediate object $\alpha_{*!}1_{X_i}$ and apply the natural transformation of left Kan extensions induced by $\lambda:p_i\Rightarrow p_j\alpha_*$.  Then use the counit-induced map $\alpha_{*!}1_{X_i}\to1_{X_j}$.  The first arrow is the mate of $(p_j\alpha_*)^*\Rightarrow p_i^*$, whose component on a presheaf $F$ is $F(\lambda_x)$.  Formula (A3) proves that this mate has the stated direction.  An invertible $\lambda$ is unnecessary.

If $I$ has terminal object $i_0$, the colimit in (A2) is its value at $i_0$.  We therefore obtain, naturally for every $p:X\to Y$,
\[
 p_!1_X\cong p_{i_0!}1_{X_{i_0}}.
\]
If $X_{i_0}$ is totally $W$-aspherical, the category of elements of the right-hand presheaf is $W$-aspherical, which proves Proposition~11.  The printed $p_{i_0!}(X_{i_0})$ on the next page is a source slip: the argument must be $1_{X_{i_0}}$.  Taking $X=Z\times I$ proves Corollary~3, since the product projection is a split opfibration with final fibre $Z$.  The argument in fact requires only $W_0$-properness of $f$ together with the stated total $W$-asphericity of its final fibre.
\end{dcomment}


\msp{page 120}
\textbf{Corollaire 2 :} \textit{Sous les conditions du corollaire 1, supposons que $I$ ait un objet final $i_0$. Alors}
\[
(p_{i_0})_!(e_{X_{i_0}}) \longrightarrow p_!(e_X)
\]
\textit{est un isomorphisme dans $Y^\wedge$.}

\dgrblock{dgrb549}{The transition map on representatives}
\begin{dcomment}{exposition}{p119}{The transition map on representatives}
The footnote asks that the arrow induced by $\lambda_\alpha$ be written out. This matters because the triangle is only laxly commutative: replacing its natural transformation by an equality would erase the data that moves a fibre over $i$ to a fibre over $j$.

Suppose $X\to I$ is the split opfibration associated to a strict functor $D:I\to\Cat$, and write $X_i=D(i)$. An arrow $\alpha:i\to j$ gives $\alpha_*:X_i\to X_j$ and chosen opcartesian arrows $x\to\alpha_*x$ in $X$. Applying $p:X\to Y$ gives a natural transformation
\[
 \lambda_\alpha:p_i\longrightarrow p_j\alpha_*.
\]
For $y\in Y$, an object of $y\backslash X_i$ is a pair $(x,\beta:y\to p_i(x))$. Define
\[
 L_{\alpha,y}(x,\beta)
   =\bigl(\alpha_*x,\lambda_{\alpha,x}\beta\bigr)
     \quad\text{in }y\backslash X_j.                 \tag{C3}
\]
If $h:x\to x'$ is a morphism in the first comma category, so that $p_i(h)\beta=\beta'$, its image is $\alpha_*h$. Naturality of $\lambda_\alpha$ gives
\[
 p_j(\alpha_*h)\lambda_{\alpha,x}\beta
 =\lambda_{\alpha,x'}p_i(h)\beta
 =\lambda_{\alpha,x'}\beta',
\]
which verifies the comma-category condition. Hence (C3) is a functor. Taking connected components yields the value at $y$ of the required presheaf map
\[
 (p_i)_!e_{X_i}\longrightarrow(p_j)_!e_{X_j}.
\]
Precomposition with a map $y'\to y$ commutes strictly with (C3), proving naturality in $y$.

This also identifies the two arrows in the manuscript's displayed composite. First keep the object $x$ in $X_i$ and replace $\beta$ by $\lambda_{\alpha,x}\beta$. That is the map induced by $\lambda_\alpha$ from $(p_i)_!e_{X_i}$ to $(p_j\alpha_*)_!e_{X_i}$. Next apply $\alpha_*$ to $x$; under transitivity of Kan extension this is the counit map
\[
 (p_j)_!(\alpha_*)_!e_{X_i}\longrightarrow(p_j)_!e_{X_j}.
\]
Their composite is exactly (C3), so its direction is fixed without an appeal to an unspecified mate convention.

Functoriality in $i$ follows from the split identities. For $\alpha:i\to j$ and $\gamma:j\to k$ one has
\[
 (\gamma\alpha)_*=\gamma_*\alpha_*,
 \qquad
 \lambda_{\gamma\alpha,x}
 =\lambda_{\gamma,\alpha_*x}\lambda_{\alpha,x}.
\]
Substitution in (C3) gives $L_{\gamma,y}L_{\alpha,y}=L_{\gamma\alpha,y}$; identity arrows act identically. Thus the transition maps constitute an actual functor $I\to\widehat Y$, not merely objectwise maps in a localization. With a non-split opfibration the comparison isomorphisms of a cleavage make the same formulas agree after taking components.

For a general $W$-proper functor, still assuming W(5), the manuscript obtains these transitions through the comma-category construction. The projection $y\backslash Y\to Y$ is a discrete opfibration, hence proper; base change and composition make $y\backslash X\to I$ proper. Its fibre over $i$ is $y\backslash X_i$. Proposition~12 applied there gives
\[
 \pi_0(y\backslash X)
   \cong\varinjlim_{i\in I}\pi_0(y\backslash X_i).
\]
The construction is natural in $y$ and so proves the presheaf isomorphism of Corollary~1. If $I$ has a terminal object, the colimit is its value at that object, recovering (C1). In the product case $X=Z\times I$, the terminal fibre is $Z$, which proves Corollary~3 as well.
\end{dcomment}


D'où aussitôt la proposition 11 : si $X_{i_0}$ est totalement $W$-asphérique, alors $X$ l'est, car pour tout $p : X \longrightarrow Y$, $p_!(e_X) \simeq p_{i_0\,!}(X_{i_0})$, et ce dernier élément de $Y^\wedge$ est $W$-asphérique % [sic] p_{i_0 !}(X_{i_0}) -> p_{i_0 !}(e_{X_{i_0}}), elsewhere the argument of p_{i!} is e_{X_i}
par hypothèse sur $X_{i_0}$.

Notons le cas particulier :

\textbf{Corollaire 3.} \textit{Soit $Z$ un objet totalement $W$-asphérique dans} Cat. \textit{Alors son produit $Z\times I$ par toute $I$ dans} Cat \textit{ayant un objet final est aussi totalement $W$-asphérique.}
% foot: Version 7 décembre 2002
\printedpage{82}
% head: A. Grothendieck : Dérivateurs, Ch. XII, Édition des Universités de Montpellier II et Paris VII

\dgrblock{dgrb550}{Total asphericity depends on more than homotopy type}
\begin{dcomment}{development}{p120:ayant un objet final est aussi totalement}{Total asphericity depends on more than homotopy type}
The transfer result just proved has a stronger content than invariance of the classifying space.  The finite construction developed below, in ``A strict cocone can repair the wrong pushout'', makes this distinction explicit.

\textbf{Corollary (failure of homotopy invariance).}\ours
Total $W_\infty$-asphericity in the sense of this chapter is not invariant under weak homotopy equivalences of categories.

\emph{Proof.}
The cited construction supplies finite face posets $E,S$ and a functor $f:E\to S$ such that $BE$ is contractible, $BS\simeq S^2$, and every $E/s$ is nonempty and connected.  Their construction and all these properties are proved there.  Set $p=f^{\mathrm{op}}:E^{\mathrm{op}}\to S^{\mathrm{op}}$.  For each $s\in S$,
\[
 (p_!1_{E^{\mathrm{op}}})(s)
   =\pi_0(s\backslash E^{\mathrm{op}})
   =\pi_0((E/s)^{\mathrm{op}})=\{*\}.
\]
Thus $p_!1_{E^{\mathrm{op}}}$ is the terminal presheaf on $S^{\mathrm{op}}$, whose category of elements is $S^{\mathrm{op}}$.  Its classifying space is a two-sphere, so $E^{\mathrm{op}}$ is not totally $W_\infty$-aspherical.  Nevertheless $E^{\mathrm{op}}\to e$ is a weak homotopy equivalence.  The category $e$ is totally $W_\infty$-aspherical: along any functor selecting $y\in Y$, its terminal presheaf is sent by left Kan extension to the representable $Y(-,y)$, and the category of elements $Y/y$ has a terminal object.\hfill$\square$

This corollary is a consequence of the explicit finite calculation below, not an assertion that the underlying topological phenomenon is new.  It explains why Proposition~11 uses properness to identify ordinary Kan extensions rather than inferring total asphericity merely from $X_{i_0}\to X\in W$.
\end{dcomment}


\msp{page 121}

\dgrblock{dgrb551}{A span is remembered by its transition arrows}
\begin{dcomment}{exposition}{p120}{A span is remembered by its transition arrows}
The next construction turns an arbitrary span into a category in which the gluing instructions remain visible as morphisms.  The ordinary pushout identifies their images.  Integration keeps those images separate and inserts comparison arrows between them.  This small change is what retains the homotopy information of the gluing.

Write $\Phi=(1\to0\leftarrow2)$ and let the span $X_1\xleftarrow{\alpha_1}X_0\xrightarrow{\alpha_2}X_2$ be the functor $\Phi^{\mathrm{op}}\to\Cat$.  Its fibred Grothendieck construction $\mathcal X$ has the disjoint union of the three sets of objects.  Within a fibre the morphisms are unchanged.  The only morphisms between distinct fibres go from a side fibre to the middle fibre, and are
\[
 \operatorname{Hom}_{\mathcal X}(a,x)
     =\operatorname{Hom}_{X_k}(a,\alpha_kx),
 \qquad a\in X_k,\ x\in X_0,\ k=1,2.
 \tag{A4}
\]
Composition on the source side is ordinary composition in $X_k$; composition with $v:x\to x'$ on the target side uses $\alpha_k(v)$.  These rules satisfy associativity because $\alpha_k$ is a functor.  They also show directly that $\mathcal X\to\Phi$ is a split fibration.

The identity of $\alpha_kx$ determines a morphism
\[
 i_k\alpha_k(x)\longrightarrow i_0(x).
\]
Naturality follows from the composition rule in (A4).  These two transformations are the missing sides of the apparently noncommutative square on manuscript p.~123.  Their sources and targets are usually different objects, so the square is not strictly commutative.  Natural transformations give homotopies on nerves; for every fundamental localizer the resulting functors become equal after localization.  The assertion that the square commutes in $\mathrm{Hot}_W$ is therefore correct, but does not replace the two transformations that constructed it.

This is the variance of the \emph{smooth} integral used in this section.  An opfibred construction of the covariant span places the transition arrows in the opposite direction.  Grothendieck compares those two constructions later in \S3.3.  Reversing only the base while leaving the fibres untouched is not an innocuous notational change.

For $W=W_\infty$, realization of the nerve of $\mathcal X$ is a homotopy pushout of $BX_1\leftarrow BX_0\to BX_2$, where $BC=|NC|$.  The closed-piece calculation below gives a direct proof in this particular shape.  The general Grothendieck-construction theorem is due to \citet[Theorem~1.2]{Thomason1979}; a primary modern proof and an explicit comparison map are given in \citet[\S3, equation~(9) and Theorem~3]{Cegarra2020}.  The manuscript expresses the corresponding fact intrinsically as a left Kan extension in the category of coherent diagrams.  Its footnote~84 correctly retains the stronger localization hypothesis needed by the preceding construction of that adjoint.
\end{dcomment}


\section*{3\quad Sommes amalgamées et intégration\quad \begin{tikzcd}[ampersand replacement=\&, row sep=small, column sep=small, baseline=(\tikzcdmatrixname-2-1.base)]
\& X_1 \\
\int X_0 \arrow[ur] \arrow[dr] \& \\
\& X_2
\end{tikzcd}}

\subsection*{3.1\quad Réduction canonique au cas d'un diagramme d'immersions fermées}

Considérons un diagramme
\[
\begin{tikzcd}[row sep=small, column sep=small]
& X_1 \\
X_0 \arrow[ur, "\alpha_1"] \arrow[dr, "\alpha_2"'] & \\
& X_2
\end{tikzcd}
\tag*{$(*)$}
\]
dans Cat. Nous pouvons le regarder comme un contrafoncteur
\[
\Phi^\circ \longrightarrow \mathrm{Cat},
\]
où
\[
\Phi = \left( \vcenter{\hbox{\begin{tikzcd}[row sep=small, column sep=small]
& 1 \arrow[dl] \\
0 & \\
& 2 \arrow[ul]
\end{tikzcd}}} \right),
\]
donc comme définissant une catégorie fibrée (scindée) $\mathscr{X}$ sur $\Phi$. Il est parfois commode de représenter cette catégorie $\mathscr{X}$ par le ‘signe d'integration’ % [sic] integration -> intégration, accent missing in the print
\[
\begin{tikzcd}[row sep=small, column sep=small]
& X_1 \\
\int X_0 \arrow[ur] \arrow[dr] & \\
& X_2,
\end{tikzcd}
\]
et la regarder comme une sorte de ‘somme amalgamée homotopique’ de $X_1$ et $X_2$ sous $X_0$, ou comme une ‘limite inductive homotopique’ du diagramme de départ $(*)$.

\msp{page 122}
Il faut entendre par là que (pour $W \subseteq \mathrm{Fl}(\mathrm{Cat})$ donné) l'image $\mathrm{hot}_W(\mathscr{X})$ de $\mathscr{X}$ dans $\mathrm{Hot}_W$ a un ‘sens homotopique’ remarquable en termes de $(*)$, de façon précise, le diagramme $(*)$ définit un objet $\xi$ de $\mathrm{HOT}_W(\Phi)$ (qui est une catégorie de fractions de la catégorie de tous les diagrammes du type $(*)$ \dots), et si
\[
\pi : \Phi \longrightarrow e
\]
désigne la projection absolue de $\Phi$, on a
\[
\pi_!(\xi) = \mathrm{hot}_W(\mathscr{X}),
\]
% foot: Version 7 décembre 2002


\dgrblock{dgrb552}{Integration as a calculus of gluing}
\begin{dcomment}{context}{p121}{Integration as a calculus of gluing}
Grothendieck's integral is designed to make a family and its transition maps into one geometric object. Grothendieck's discussion of integration in Chapter~VI, manuscript pages~81 and~146, gives the general background. Here the additional issue is quite specific: the total category must retain enough information to decide whether a given cocone is a homotopy pushout.

In \emph{Pursuing Stacks}, \S69, \citet{Grothendieck1983} connects diagram categories to the trace problems of SGA~5, then describes integration by a total cofibred category. He also singles out presheaf categories as a setting where actual amalgamated sums are available. These two concerns meet in the present pages. The integral supplies the homotopical comparison; the last theorem shows that a pushout of presheaves along a monomorphism gives an actual categorical pushout with the required homotopy type. The intervening analysis of open immersions is the bridge between the two constructions.

The letter to Thomason of 2 April 1991 gives a closely related explanation of why one diagram category cannot suffice: a derived category ``ne vient jamais seule'' \citep[\S2]{Grothendieck1991Letter}. In \S1 of the same letter, Grothendieck places weak equivalences and localization before the auxiliary classes used in model structures. This explains the order of the present argument. He constructs a comparison map first and asks whether it belongs to $W$; cofibration criteria are then ways of proving that assertion.

The source is also a record of the distinctions being worked out. A natural transformation provides a homotopy without becoming an equality of functors; a weak equivalence of total categories need not preserve every ordinary Kan extension; a gluing theorem need not reflect equivalences backwards. The proofs below keep the transition arrows, the presheaf adjunctions, and the excision hypotheses visible at precisely the steps where each is used.
\end{dcomment}


\dgrblock{dgrb553}{Integration retains the homotopy of the span}
\begin{dcomment}{exposition}{p122}{Integration retains the homotopy of the span}
The shift to integration changes what is being glued. An ordinary pushout identifies the common category immediately. The integrated category instead retains a copy of that category and arrows connecting it to both sides. These arrows provide the homotopies needed for a derived colimit.

For the span $X_1\xleftarrow{\alpha_1}X_0\xrightarrow{\alpha_2}X_2$, the manuscript uses the indexing category $\Phi=(1\to0\leftarrow2)$ and a contravariant functor on it. The associated fibred category $\mathscr X$ has objects $(k,x)$ with $x\in X_k$. It contains each $X_k$ fully, and its additional morphisms are
\[
 \Hom_{\mathscr X}((k,x),(0,z))
   =\Hom_{X_k}(x,\alpha_kz),\qquad k=1,2.
\]
There are no arrows between the $1$- and $2$-fibres, nor from the $0$-fibre to either of them. Composition with an arrow in $X_k$ is precomposition; composition with an arrow $z\to z'$ in $X_0$ is postcomposition by its image under $\alpha_k$. These rules give the split Grothendieck construction with precisely the variance used by Grothendieck. This continues the discussion of variance in Chapter~VI, manuscript p.~81, rather than replacing it by a covariant convention without explanation.

For this span one can prove the homotopy-pushout assertion directly. Let $M_k$ be the full subcategory of $\mathscr X$ on its $0$- and $k$-fibres. It is a categorical mapping cylinder for $\alpha_k$. The functor
\[
 r_k:M_k\longrightarrow X_k
\]
is the identity on $X_k$ and sends $z\in X_0$ to $\alpha_kz$. If $s_k:X_k\hookrightarrow M_k$ is the inclusion, then $r_ks_k=\id$, and the canonical arrow $\alpha_kz\to z$ in $M_k$ gives a natural transformation $s_kr_k\to\id_{M_k}$, equal to the identity on $X_k$. Thus $Br_k$ is a homotopy equivalence, and its restriction to $BX_0$ is $B\alpha_k$.

Every composable string in $\mathscr X$ lies in $M_1$ or $M_2$, and their intersection is exactly $X_0$. Therefore
\[
 N\mathscr X=NM_1\amalg_{NX_0}NM_2.
\]
Both maps out of $NX_0$ are monomorphisms, so this is a homotopy pushout of simplicial sets. The two homotopy equivalences $Br_k$, together with the identity on $BX_0$, now identify its homotopy type with
\[
 B\mathscr X\simeq
 \operatorname{hocolim}\bigl(BX_1\leftarrow BX_0\to BX_2\bigr).
\]
This is the span instance of \citet[Theorem~1.2]{Thomason1979}, whose statement is also recorded by \citet[Corollary~26.9]{ChacholskiScherer2001}; the direct proof records the orientation and the connecting morphisms needed in these pages.

For example, take $X_0$ to be two discrete objects and $X_1=X_2=e$. The ordinary pushout is $e$. The integrated category has four objects and four nonidentity arrows: each of the two side objects maps to each of the two middle objects. Its nerve is a square with no two-dimensional nondegenerate simplices, hence has realization $S^1$. The integration remembers the circle formed by the two independent identifications; the ordinary pushout has forgotten it.

The final formula $\pi_!(\xi)=\mathrm{hot}_W(\mathscr X)$ expresses this construction in the derivator already under discussion. It is a canonical isomorphism in the localized category, rather than a literal equality of chosen representatives. For $W_\infty$ the preceding argument gives its concrete meaning. For another $W$ the formula presupposes the standing hypotheses under which the diagram localization and its homotopy Kan extension are available; an arbitrary class of arrows would not by itself justify that notation. The advertised reduction to a diagram of closed immersions begins here; it is not completed on manuscript page~122.
\end{dcomment}
\begin{otherlanguage}{english}
\dgroverview{XIIb-a}{Chapter XII \textperiodcentered\ \textup{XIIb-a} \textperiodcentered\ Overview}
\begingroup
\selectlanguage{english}
\noindent\textbf{Attachments, ordinary direct images, and homotopy.}

These pages, Chapter~XII, section~V, manuscript pp.~59--122, begin with a concrete problem: how to attach one category to another along a full subcategory, and how to recognize when the attachment respects weak equivalences. The first part answers the categorical question completely. The second discovers how much more is needed for the homotopical question. Along the way Grothendieck introduces total asphericity, isolates the Dwyer condition, and reaches the Grothendieck construction as a replacement for an ordinary pushout. The text is exploratory in a mathematically consequential sense. Some tentative arguments are explicitly retracted; one proposed theorem survives in the text even though its asserted reverse implication is false. The commentary keeps these different situations separate.

An ``open'' subcategory $U\subseteq X$ is a sieve: an arrow whose target lies in $U$ has its source there as well. Its complementary full subcategory $T$ is a cosieve. There are no arrows from $T$ to $U$, so the information not already contained in these two categories is exactly the bifunctor
\[
 H:U^{\mathrm{op}}\times T\longrightarrow\mathrm{Set},
 \qquad H(u,t)=\Hom_X(u,t).
\]
Its two functorial actions encode composition on either side of a crossing arrow. Conversely any such bifunctor constructs a category on the disjoint union of the object sets. This is the collage of a profunctor. Maps and natural transformations out of the collage have equally explicit descriptions. The formulas on manuscript p.~61 must be read relatively: changing $U$ and $T$ by pullback of $H$ gives a cartesian lifting over the category of pairs of categories, not the unqualified representability asserted in the printed displays. The corrected universal property is proved in the commentary.

This description also explains the relation with open--closed gluing of topoi. A presheaf on the collage is a pair of presheaves on $U$ and $T$, together with the restriction data across $H$. The gluing functor is
\[
 \varphi(F)(t)=\Hom_{\widehat U}(H(-,t),F).
\]
It preserves limits and is right adjoint to the cocontinuous extension of $t\mapsto H(-,t)$. The converse recovers the collage from an accessible gluing functor preserving small limits. This is already the construction of SGA~4, Expos\'e~IV, Exercise~9.5.10; the comparison is exact, including the characterization of the resulting topos as a presheaf topos \citep{SGA4I1972}. Natural transformations between right adjoints reverse the direction of transformations between their left adjoints, a convention that must be retained when stating the category-level equivalence. The historical connection helps explain why the categorical attachment problem first appears here in the language of an open and its closed complement.

The advantage of the collage is that a pushout can be computed by an ordinary Kan extension. Given $p:U\to U'$ and $r:T\to T'$, the new crossing bifunctor is
\[
 H'=\operatorname{Lan}_{p^{\mathrm{op}}\times r}H.
\]
Its pointwise coend description records a crossing arrow together with paths in the new categories to its two ends, modulo the relations imposed by composition. When only $U$ changes, this becomes
\[
 \Hom_{X'}(u',t)\cong\pi_0(u'\backslash U/t),
 \qquad X'=X\amalg_U U'.
\]
The category on the right remembers factorizations through $U$; the ordinary pushout retains only their connected components. This formula proves that $U'$ is again a sieve, that the complement remains $T$, and that the support of the gluing weight is unchanged. It is also the source of the later difficulty: taking components does not preserve the higher homotopy of a factorization category.

\medskip
\Needspace{5\baselineskip}
\noindent\textbf{The gluing theorem that survives, and the reverse implication that fails.}

If a category is the union of two sieves, every simplex of its nerve lies in one of them. Its nerve is therefore the simplicial union along the intersection, and the square is a homotopy pushout. Consequently weak equivalences on the intersection and on both pieces induce a weak equivalence on the union. This is the valid gluing direction. It does not imply the symmetric rule W(7), which asks that any two among equivalence on the intersection, equivalence on both pieces, and equivalence on the union imply the third.

Theorem~2 on manuscript p.~78 asserts that this stronger rule holds for $W_\infty$. A complete counterexample uses a connected acyclic noncontractible finite complex. Take a finite poset $A$ with that homotopy type and adjoin two incomparable terminal cone points. The two one-cone subcategories are sieves, each contractible, and their union has the homotopy type of the suspension of $BA$. The suspension is simply connected and acyclic, hence contractible, whereas $BA$ remains noncontractible. Thus the map from a cone to the union is a weak equivalence, but the opposite map from the intersection to the other cone is not. The counterexample refutes both the reverse part of W(7 bis) and the full W(7); it does not merely exhibit a missing proof.

The commentary gives the classical two-cell input explicitly, verifies its homology and a nontrivial fundamental-group quotient, and passes to a finite simplicial model and its face poset \citep[Example~2.38 and Theorem~2C.5]{Hatcher2002}. A further development supplies finite tests against every bounded homotopy-group version of the axiom. It also explains why the proposed Mayer--Vietoris argument succeeds for fixed-coefficient homology equivalences: the three mapping-cone complexes form a short exact sequence. The failure occurs in passing from that homological statement to the asserted unstable homotopical converse.

This finding changes how the subsequent pages must be read. Assertions deduced under W(7) or W(7 bis) remain conditional assertions. Their application to $W_\infty$ cannot use Theorem~2. In particular, the proposed one-point converse to asphericity on pp.~84--85 is not an unconditional theorem about weak homotopy equivalences. The audit records the actual hypotheses and, where possible, supplies a direct proof of a useful consequence that avoids the false reverse principle.

\medskip
\Needspace{5\baselineskip}
\noindent\textbf{Total asphericity measures ordinary images.}

For a presheaf $F$ on $U$, total $W$-asphericity means that the category of elements of $p_!F$ is $W$-aspherical for every functor $p:U\to U'$. These are ordinary set-valued Kan extensions. For a category $V$, the definition applies to its terminal presheaf. It is substantially stronger than $V\to e\in W$, and it is not interchangeable with other uses of ``totally aspherical'' in the literature on test categories.

The relevant ordinary factorization sends $f:V\to V'$ to
\[
 V\longrightarrow V'/f_!e_V\longrightarrow V'.
\]
The second map is a discrete fibration; the first is final in the ordinary categorical sense. Its undercategories are nonempty and connected, rather than necessarily contractible. This is the comprehensive factorization. The commentary proves its universal property and relates it to the treatment of Street--Walters in Kelly's account \citep[\S4.7]{Kelly1982}. It yields a useful formulation developed here: a category is totally $W$-aspherical exactly when every target of an ordinary final functor out of it is $W$-aspherical. In particular the class of totally aspherical categories is closed under ordinary final images. This reformulation identifies the quantified condition without confusing ordinary finality with homotopy finality.

Grothendieck's attempted stability under unions fails because a left Kan extension need not preserve the monomorphisms describing a cover. His attempted use of retractions does not repair that defect: a categorical retraction does not supply a splitting of the associated inclusion of presheaves. A finite example makes both issues explicit. A final functor from the contractible outgoing three-object fork lands in a four-object category whose classifying space is $S^2$. Its relevant comma category is a square, connected but not contractible. Two distinct representatives of a Kan-extension value become equal after adjoining one branch. The construction was located in Grothendieck's Appendix~XII and is credited there. It also produces the pushout failures on pp.~104--107; adjoining a further cone gives a classifying space $S^3$, so their validity can be checked independently of the disputed one-point converse.

The converse from totality of a presheaf to totality of its category of elements is established here for subterminal presheaves, as the source proposes. A general answer was not located in the primary literature checked. The same limitation applies to certain composition and classification questions posed later in these pages. An unsuccessful search is not turned into a claim of present-day openness.

\medskip
\Needspace{5\baselineskip}
\noindent\textbf{Good attachments and the passage to integration.}

A $W$-Dwyer sieve asks that every nonempty crossing weight be totally $W$-aspherical. Ordinary Kan-extension transitivity proves that this condition is preserved by pushout. The stronger version, expressed through totality of the categories of elements themselves, uses the final first map of the comprehensive factorization. Classical Dwyer immersions are the representable case: the inclusion of $U$ into the cosieve it generates admits a right adjoint. Filtered examples show why total asphericity is broader than the existence of a terminal object. The interval-product calculation on p.~114, however, shows directly why the classical adjoint condition is too rigid for every intended homotopy attachment.

The manuscript's $W$-cofibrations have since been studied as $h$-cofibrations, defined by preservation of weak equivalences under cobase change on undercategories. In left proper model categories they characterize attachments whose pushouts compute homotopy pushouts; the theory also recovers the equivalence between trivial $h$-cofibrations and couniversal weak equivalences \citep[Proposition~1.6, Corollary~1.9 and Lemma~1.10]{BataninBerger2017}. This is a precise later form of the useful gluing principle. For classical Dwyer maps, the nerve comparison is now known to be a weak categorical equivalence, with a generalization to discretely flat cofibrations \citep[Theorems~1.5--1.6]{HackneyOzornovaRiehlRovelli2024}. These later notions are compared with the definitions here without identifying their different ambient weak equivalences.

The final properness criterion is independent of the failed cube axiom. Under the necessary connectivity hypothesis, the inclusion of a terminal fibre is ordinary final. Its terminal presheaf therefore has direct image equal to the terminal presheaf of the total category. Composing Kan extensions proves that total asphericity of that fibre implies total asphericity of the total category. The commentary gives the general component-colimit formula and writes out the natural transformation requested in the footnote on p.~119, including its action on representatives and its compatibility with composition.

The text then begins a new construction. Integrating a span retains its middle category and the arrows connecting it to both sides. For two points mapping to a point on each side, the ordinary pushout is a point, whereas the integrated category has classifying space a circle. Two categorical mapping cylinders give a direct proof that the integrated span models the homotopy pushout. The final displayed identity $\pi_!\xi\cong\mathrm{hot}_W(\mathscr X)$ states the intended derivator interpretation. The announced reduction to closed immersions continues after that identity, beyond manuscript page~122.

\endgroup
\end{otherlanguage}
\dgrchunkend
